\chapter{Elliptic curves over fields}\label{ch:ellcurves}

Every object in the paper is built from an elliptic curve $E/\Q$: its Weierstrass invariants and $v_q(j_E)$, its torsion points and isogenies, its Tate module $T_\ell E$ with the Galois
representation $\rho_{E,\ell}$, the numbers $a_p$ and the roots $\alpha,\beta$ of $X^2-a_pX+p$, its quadratic twists, its period lattice, its $3$-division polynomial. This chapter develops
these objects over an arbitrary field, then over $\C$, over finite fields and over $\Q$, proving what the paper uses and stating the deep theorems (Mordell-Weil, Mazur, Kenku, Faltings,
Deuring, modularity) precisely with references. It uses \Cref{ch:groups} (finite abelian groups, $\GL_2$, inverse limits), \Cref{ch:rings} (finite fields, quadratic characters),
\Cref{ch:reps} (twists, contragredients, adjoints), \Cref{ch:galois} (Galois theory of polynomials, profinite groups), \Cref{ch:ant} (Frobenius elements, imaginary quadratic orders), and
above all \Cref{ch:curves} (curves, divisors, Riemann-Roch, Weierstrass models, morphisms and their degree, the Abel-Jacobi theorem). Local questions (minimal models, reduction types,
the Tate curve) are \Cref{ch:reduction}; the images of the Galois representations are \Cref{ch:galrep}.

\section*{Core results at a glance}

Throughout, $K$ is a field, $\bar K$ an algebraic closure, $G_K=\Gal(K^{\mathrm{sep}}/K)$, $E$ an elliptic curve over $K$ given by a Weierstrass equation, and $\ell\ne\operatorname{char}K$ a prime.

\begin{enumerate}[label=\textbf{\thechapter.\arabic*}]
\item \textbf{Conventions.} Elliptic curves are Weierstrass cubics with $O=[0:1:0]$; over $\Q$, global minimal models as in Cremona's tables; $\chi_d$ is the character of $K(\sqrt d)$,
$\chi_{q^*}$ the quadratic character of conductor $q$ (never $\chi_q$), $\chi_\ell$ the cyclotomic character; $e_m$ is Silverman's Weil pairing. \Cref{conv:ellcurves:conventions}.
\item \textbf{Weierstrass invariants} (definition): $b_2,b_4,b_6,b_8$, $c_4=b_2^2-24b_4$, $c_6=-b_2^3+36b_2b_4-216b_6$, $\Delta$, $j=c_4^3/\Delta$, $\omega=dx/(2y+a_1x+a_3)$.
\Cref{def:ellcurves:weierstrass}.
\item \textbf{Identities.} $4b_8=b_2b_6-b_4^2$, $1728\Delta=c_4^3-c_6^2$; $(2y+a_1x+a_3)^2=4x^3+b_2x^2+2b_4x+b_6$, a cubic of discriminant $16\Delta$; for $\operatorname{char}K\ne2,3$ the short form
$Y^2=X^3-27c_4X-54c_6$. \emph{(proved here)} \Cref{lem:ellcurves:identities}.
\item \textbf{Smoothness.} A Weierstrass cubic is smooth iff $\Delta\ne0$; if $\Delta=0$ it has one singular point, a node if $c_4\ne0$, a cusp if $c_4=0$; every elliptic curve is a smooth
Weierstrass cubic. \emph{(proved here for $\operatorname{char}K\ne2$; reference in characteristic $2$)} \Cref{prop:ellcurves:smooth}.
\item \textbf{Changes of coordinates.} $x=u^2x'+r$, $y=u^3y'+su^2x'+t$ give $u^4c_4'=c_4$, $u^6c_6'=c_6$, $u^{12}\Delta'=\Delta$, $j'=j$, $\omega'=u\omega$; every isomorphism fixing $O$ is of this form; over $\Q$,
$v_q(j)=3v_q(c_4)-v_q(\Delta)$ is independent of the model and $v_q(\Delta)$ changes by multiples of $12$. \emph{(proved here)} \Cref{prop:ellcurves:coordchange}.
\item \textbf{The $j$-invariant.} Curves are isomorphic over $\bar K$ iff they have the same $j$; every $j_0\in K$ occurs. \emph{(proved here for $\operatorname{char}K\ne2,3$; reference in general)}
\Cref{thm:ellcurves:jinv}.
\item \textbf{Automorphisms.} For $\operatorname{char}K\ne2,3$, $\#\Aut_{\bar K}(E)=2,4,6$ for $j\ne0,1728$, $j=1728$, $j=0$; $[-1](x,y)=(x,-y-a_1x-a_3)$. \emph{(proved here)} \Cref{prop:ellcurves:aut}.
\item \textbf{The chord-tangent law} (construction): $P+Q+R=O$ for the three intersections of a line; explicit formulas for $P_1+P_2$ and $x(2P)=(x^4-b_4x^2-2b_6x-b_8)/(4x^3+b_2x^2+2b_4x+b_6)$; for
$\operatorname{char}K\ne2$, $\#E[2]=4$. \emph{(proved here)} \Cref{con:ellcurves:chordtangent}, \Cref{prop:ellcurves:addition}.
\item \textbf{Chord-tangent $=\Pic^0$.} $P\mapsto[(P)-(O)]$ is additive for the chord-tangent law. \emph{(proved here, given \Cref{ch:curves})} \Cref{prop:ellcurves:pic0}.
\item \textbf{The group law.} $E(\bar K)$ is an abelian group, $E(L)$ a subgroup for $K\subseteq L\subseteq\bar K$, $G_K$ acts by automorphisms, addition and translations are morphisms. \emph{(proved here;
the morphism property at the excluded points by reference)} \Cref{thm:ellcurves:grouplaw}.
\item \textbf{Principal divisors.} $\sum n_P(P)$ is principal iff $\sum n_P=0$ and $\sum[n_P]P=O$. \emph{(proved here)} \Cref{cor:ellcurves:principal}.
\item \textbf{The invariant differential.} $\operatorname{div}\omega=0$, $\tau_Q^*\omega=\omega$, $(\varphi+\psi)^*\omega=\varphi^*\omega+\psi^*\omega$, $[m]^*\omega=m\omega$. \emph{(proved here, in every characteristic; one polynomial identity checked by
machine; a second proof in characteristic $0$ in \Cref{rem:ellcurves:invdiff-C})} \Cref{prop:ellcurves:invdiff}.
\item \textbf{Isogenies} (definition): morphisms with $\varphi(O)=O$; degree $[\bar K(E_1):\varphi^*\bar K(E_2)]$; isogenous over $K$. \Cref{def:ellcurves:isogeny}.
\item \textbf{First properties.} Isogenies are $0$ or surjective; $\Hom$ is a group, $\deg$ is multiplicative; $[m]\ne0$ for $m\ne0$, so $\Hom$ is torsion-free. \emph{(proved here; characteristic $2$ by reference)}
\Cref{prop:ellcurves:isogeny-basic}.
\item \textbf{Kernels.} Isogenies are homomorphisms; $\#\ker\varphi=\deg_s\varphi$; a separable isogeny is unramified, with $\#\ker\varphi=\deg\varphi$ and Galois function field extension with group
$\ker\varphi$. \emph{(proof sketched here)} \Cref{thm:ellcurves:isog-hom}.
\item \textbf{Quotients.} For a finite subgroup $H$ there is a separable $E\to E/H$ with kernel $H$, over $K$ if $H$ is $G_K$-stable (V\'elu); an isogeny factors through any separable isogeny
whose kernel is contained in its kernel. \emph{(stated; reference)} \Cref{thm:ellcurves:quotient}.
\item \textbf{The dual isogeny.} $\hat\varphi\circ\varphi=[\deg\varphi]$, $\varphi\circ\hat\varphi=[\deg\varphi]$, $\widehat{\varphi+\psi}=\hat\varphi+\hat\psi$, $\widehat{[m]}=[m]$, $\deg\hat\varphi=\deg\varphi$; $\deg$ is a positive definite
quadratic form on $\Hom$. \emph{(proof sketched here)} \Cref{thm:ellcurves:dual}.
\item \textbf{Torsion.} $\deg[m]=m^2$; if $\operatorname{char}K\nmid m$, $[m]$ is separable and $E[m]\simeq(\Z/m)^2$; in characteristic $p$, $E[p^e]$ is $0$ or $\Z/p^e$. \emph{(proved here, given
\Cref{thm:ellcurves:dual}; the last statement stated, reference)} \Cref{cor:ellcurves:degm}.
\item \textbf{Cyclic factorisation} ($\operatorname{char}K=0$). $\ker\varphi\simeq\Z/a\times\Z/ab$ and $\varphi=\psi\circ[a]$ with $\psi$ cyclic over $K$; a cyclic $n$-isogeny over $K$ has $G_K$-stable kernels of every
order $d\mid n$, and factors into cyclic isogenies of prime degree over $K$. \emph{(proved here)} \Cref{prop:ellcurves:cyclicfactor}.
\item \textbf{Division polynomials} (definition): $\psi_2=2y+a_1x+a_3$, $\psi_3=3x^4+b_2x^3+3b_4x^2+3b_6x+b_8$, $\psi_4$, and the recursions. \Cref{def:ellcurves:divpol}.
\item \textbf{Degrees.} For odd $m$, $\psi_m\in\Z[a_i][x]$ has degree $(m^2-1)/2$ and leading coefficient $m$; for even $m$, $\psi_m/\psi_2$ has degree $(m^2-4)/2$ and leading coefficient $m/2$.
\emph{(proved here)} \Cref{prop:ellcurves:divpol-deg}.
\item \textbf{Multiplication by $m$.} For $\operatorname{char}K\nmid m$: $mP=O$ iff $\psi_m(P)=0$, and $x(mP)=x-\psi_{m-1}\psi_{m+1}/\psi_m^2$. \emph{(proof sketched here)} \Cref{thm:ellcurves:multm}.
\item \textbf{Roots of $\psi_m$.} For odd $m$, $\operatorname{char}K\nmid m$: the roots of $\psi_m$ are the $(m^2-1)/2$ distinct values $x(P)$, $P\in E[m]\smallsetminus\{O\}$, $G_K$-equivariantly. \emph{(proved here,
given \Cref{thm:ellcurves:multm})} \Cref{cor:ellcurves:divpol-roots}.
\item \textbf{The $3$-division polynomial.} $x(2P)-x(P)=-\psi_3/\psi_2^2$; for $\operatorname{char}K\ne2,3$ the four roots of $\psi_3$ are the $x$-coordinates of the four lines of $E[3]$, $G_K$-equivariantly, so the
factorisation pattern of $\psi_3$ is the orbit pattern of $\bar\rho_{E,3}(G_K)$ on $\PP^1(\F_3)$, and a root in $K$ means a $K$-rational $3$-isogeny; over $\Q$, $\disc\psi_3\in-3(\Q^\times)^2$.
\emph{(proved here)} \Cref{prop:ellcurves:psi3}.
\item \textbf{The Weil pairing.} $e_m\colon E[m]\times E[m]\to\mu_m$ is bilinear, alternating, nondegenerate, Galois-equivariant, compatible in $m$, and $e_m(S,\hat\varphi T)=e_m(\varphi S,T)$. \emph{(proof
sketched here)} \Cref{thm:ellcurves:weil}.
\item \textbf{The determinant.} $G_K$ acts on $E[m]$ with determinant $\chi_m^{\mathrm{cyc}}$; $E[m]\subseteq E(K)$ forces $\mu_m\subseteq K$; for a stable line $L\subset E[\ell]$, $E[\ell]/L\simeq\Hom(L,\mu_\ell)$, so a
rational point of order $\ell$ gives $0\to\Z/\ell\to E[\ell]\to\mu_\ell\to0$. \emph{(proved here)} \Cref{cor:ellcurves:det}.
\item \textbf{Tate modules} (definition and structure): $T_\ell E=\varprojlim E[\ell^n]$ is free of rank $2$ over $\Z_\ell$, $T_\ell E/\ell^n=E[\ell^n]$, $V_\ell E/T_\ell E=E[\ell^\infty]$; the Weil pairings give a perfect
alternating $G_K$-equivariant $e\colon T_\ell E\times T_\ell E\to\Z_\ell(1)$; in characteristic $p$, $T_pE$ is $0$ or $\Z_p$. \emph{(proved here; the last statement stated, reference)} \Cref{def:ellcurves:tate}, \Cref{prop:ellcurves:tate}.
\item \textbf{The Galois representations} (definition and properties): $\rho_{E,\ell}\colon G_K\to\GL_2(\Z_\ell)$, $\bar\rho_{E,\ell}$ on $E[\ell]$, $\rho\bmod\ell^k$ on $E[\ell^k]$; continuous, $\det\rho_{E,\ell}=\chi_\ell$; an
isogeny over $K$ gives $T_\ell E\hookrightarrow T_\ell E'$ with cokernel $(\ker\varphi)[\ell^\infty]$, so $V_\ell E\simeq V_\ell E'$, and $T_\ell E\simeq T_\ell E'$ if $\ell\nmid\deg\varphi$. \emph{(proved here)}
\Cref{def:ellcurves:rho}, \Cref{prop:ellcurves:rho}.
\item \textbf{Self-duality.} $T_\ell E\simeq(T_\ell E)^\vee(1)$, so $H^1(E,\Z_\ell)=(T_\ell E)^\vee\simeq T_\ell E(-1)$ and $\ad$ of $T_\ell E$ and of its dual agree. \emph{(proved here)} \Cref{cor:ellcurves:dual-T}.
\item \textbf{Complex conjugation.} For $E/\Q$ and $\ell$ odd, $\rho_{E,\ell}(c)\sim\operatorname{diag}(1,-1)$: eigenvalues $\pm1$, trace $0$, determinant $-1$. \emph{(proved here)} \Cref{cor:ellcurves:complex-conj}.
\item \textbf{Frobenius conventions.} $\Frob_p$ arithmetic; $\rho_{E,\ell}(\Frob_p)$ has eigenvalues $\alpha_p,\beta_p$; at good $p$, $L_p(E,s)=\det(1-\Frob_pp^{-s}\mid V_\ell E)^{-1}=L_p(V_\ell E^\vee,s)$ and
$L_p(V_\ell E,s)=L_p(E,s+1)$, with $L$-factors defined by the geometric Frobenius. \Cref{conv:ellcurves:frobenius}.
\item \textbf{Endomorphism rings.} $\End_{\bar K}(E)$ is $\Z$, an order in an imaginary quadratic field, or (characteristic $p$ only) a maximal order in the quaternion algebra ramified at $p$, $\infty$.
\emph{(stated; reference)} \Cref{thm:ellcurves:end}.
\item \textbf{Normalised endomorphisms.} In characteristic $0$, $\alpha\mapsto c_\alpha$ ($\alpha^*\omega=c_\alpha\omega$) embeds $\End(E)$ in $\bar K$; $\End_\Q(E)=\Z$; for $E/\Q$ with CM by $\mathcal O\subset F$,
$\sigma[c]\sigma^{-1}=[\sigma(c)]$, endomorphisms are defined over $F$, and $\rho(\sigma)[c]_\ell\rho(\sigma)^{-1}=[\sigma c]_\ell$. \emph{(proved here, given \Cref{prop:ellcurves:invdiff} and \Cref{thm:ellcurves:end})}
\Cref{prop:ellcurves:normalised}.
\item \textbf{Complex multiplication} (definition): $\End_{\bar K}(E)\ne\Z$; it is an isogeny invariant. \emph{(the invariance proved here)} \Cref{def:ellcurves:cm}, \Cref{lem:ellcurves:cm-isog}.
\item \textbf{The thirteen CM $j$-invariants.} $E/\Q$ with CM by $\mathcal O$ has $h(\mathcal O)=1$, $D_{\mathcal O}$ is one of $-3$, $-4$, $-7$, $-8$, $-11$, $-12$, $-16$, $-19$, $-27$, $-28$, $-43$, $-67$, $-163$, and $j(E)$ is the corresponding value of the table.
\emph{(stated; reference)} \Cref{thm:ellcurves:13j}.
\item \textbf{CM and $a_p$.} For $E/\Q$ with CM by $F$: $\rho\simeq\rho\otimes\chi_F$; $a_p=0$ at good $p$ inert in $F$; $a_p=\pi_p+\bar\pi_p$ with $\pi_p\bar\pi_p=p$ at good split $p$; every prime ramified in $F$
divides $N$ (so CM by $\Q(\sqrt{\ell^*})$ forces $\ell\mid N$); conversely $a_p=0$ at almost all inert $p$ forces CM by $F$. \emph{(proof sketched here; (i), (ii), (iv) proved in \Cref{ch:ant})}
\Cref{thm:ellcurves:cm-ap}.
\item \textbf{Lattice sums.} $\sum'\abs\omega^{-s}<\infty$ for $s>2$; $G_{2k}$, $g_2=60G_4$, $g_3=140G_6$, $\Delta(\Lambda)$, $j(\Lambda)$ are homogeneous of weights $2k,4,6,12,0$. \emph{(proved here)}
\Cref{def:ellcurves:lattice}, \Cref{lem:ellcurves:latticesum}.
\item \textbf{The Weierstrass function.} $\wp$ is even, $\Lambda$-periodic, with double poles at $\Lambda$; $\wp=z^{-2}+\sum(2k+1)G_{2k+2}z^{2k}$; $\wp'^2=4\wp^3-g_2\wp-g_3$; $\Delta(\Lambda)\ne0$; periodic functions with poles of
order $\le2$ (resp.\ $\le3$) at $\Lambda$ only are $a\wp+b$ (resp.\ $a\wp'+b\wp+c$). \emph{(proved here)} \Cref{thm:ellcurves:wp}.
\item \textbf{Uniformisation.} $z\mapsto(\wp,\wp')$ is a biholomorphic group isomorphism $\C/\Lambda\to E_\Lambda(\C)$ with $dx/y\leftrightarrow dz$; every $E/\C$ is $E_{\Lambda_E}$ for its period lattice; for $4A^3+27B^2\ne0$ there is a
unique $\Lambda$ with $g_2=-4A$, $g_3=-4B$. \emph{(proved here, given the Abel-Jacobi theorem of \Cref{ch:curves})} \Cref{thm:ellcurves:cplxuniform}.
\item \textbf{Maps of tori.} Holomorphic maps $\C/\Lambda_1\to\C/\Lambda_2$ fixing $0$ are $z\mapsto\alpha z$, $\alpha\Lambda_1\subseteq\Lambda_2$, of degree $[\Lambda_2:\alpha\Lambda_1]$; $E[m]=\frac1m\Lambda/\Lambda$; $\End(\C/\Lambda)$ is $\Z$ or an
imaginary quadratic order, the latter iff $\tau$ is imaginary quadratic; $e_m=\zeta^{\det}$ on $\frac1m\Lambda/\Lambda$. \emph{(proved here; the sign of $\zeta$ by reference)} \Cref{prop:ellcurves:cplxisog}.
\item \textbf{Moduli over $\C$.} Lattices up to homothety $=$ curves over $\C$ up to isomorphism $=\SL_2(\Z)\backslash\HH$; $j(\tau)$ is holomorphic and invariant, $j(i)=1728$, $j(\zeta_3)=0$; $j$ is a bijection
$\SL_2(\Z)\backslash\HH\to\C$. \emph{(proved here; the bijectivity and the $q$-expansion in \Cref{ch:modforms})} \Cref{thm:ellcurves:moduli}.
\item \textbf{The period lattice} (definition and properties): $\Lambda_E$ of the N\'eron differential of the global minimal model; conjugation-stable, rectangular iff $\Delta>0$; $\vol(\Lambda_E)=\frac i2\int\omega\wedge\bar\omega$;
computed by the AGM; commensurable for isogenous curves. \emph{(proof sketched here)} \Cref{def:ellcurves:periodlattice}, \Cref{prop:ellcurves:periods}.
\item \textbf{Frobenius endomorphism} (definition and properties): $\pi(x,y)=(x^{\mathfrak q},y^{\mathfrak q})$, purely inseparable of degree $\mathfrak q$, acting as $\Frob_{\mathfrak q}$; $E(\F_{\mathfrak q^n})=\ker(1-\pi^n)$; $1-\pi$
separable, $\#E(\F_{\mathfrak q})=\deg(1-\pi)$; $a_{\mathfrak q}=\mathfrak q+1-\#E(\F_{\mathfrak q})$. \emph{(proved here, given \Cref{ch:curves})} \Cref{def:ellcurves:frobendo}, \Cref{prop:ellcurves:frobendo}.
\item \textbf{Hasse.} $\abs{a_{\mathfrak q}}\le2\sqrt{\mathfrak q}$. \emph{(proved here)} \Cref{thm:ellcurves:hasse}.
\item \textbf{Characteristic polynomial.} $\det\varphi_\ell=\deg\varphi$ on $T_\ell E$; $\pi_\ell$ has characteristic polynomial $X^2-a_{\mathfrak q}X+\mathfrak q$, roots $\alpha,\beta$ with $\alpha\beta=\mathfrak q$,
$\abs\alpha=\abs\beta=\sqrt{\mathfrak q}$; $\#E(\F_{\mathfrak q^n})=\mathfrak q^n+1-\alpha^n-\beta^n$; the zeta function. \emph{(proved here)} \Cref{thm:ellcurves:charpoly}.
\item \textbf{Supersingular curves} (definition and criteria): $E[p]=0$; iff $p\mid a_{\mathfrak q}$; for $\mathfrak q=p\ge5$ iff $a_p=0$. \emph{(the first criterion stated, reference; the second proved here)}
\Cref{def:ellcurves:supersingular}, \Cref{prop:ellcurves:supersingular}.
\item \textbf{Twists over $\F_{\mathfrak q}$.} $a_{\mathfrak q}(E^{\mathrm{tw}})=-a_{\mathfrak q}(E)$ and $(\mathfrak q+1)^2-a_{\mathfrak q}^2=\#E(\F_{\mathfrak q})\#E^{\mathrm{tw}}(\F_{\mathfrak q})$. \emph{(proved here)} \Cref{prop:ellcurves:twistFq}.
\item \textbf{Good reduction.} For $E/\Q$ good at $p\ne\ell$: $\rho_{E,\ell}$ is unramified at $p$ and $\rho_{E,\ell}(\Frob_p)$ has characteristic polynomial $X^2-a_pX+p$; $\tr\bar\rho_{E,\ell}(\Frob_p)\equiv a_p$.
\emph{(proof sketched here, using \Cref{ch:reduction})} \Cref{thm:ellcurves:frobQ}.
\item \textbf{Mordell-Weil.} $E(\Q)\simeq\Z^r\oplus E(\Q)_{\mathrm{tors}}$. \emph{(stated; reference)} \Cref{thm:ellcurves:mordellweil}.
\item \textbf{Torsion and reduction.} Reduction is injective on prime-to-$p$ torsion at good $p$ (on all torsion if $p\ge3$); Nagell-Lutz; $E(\Q)_{\mathrm{tors}}$ of \curve{11a1}, \curve{20a1}. \emph{(stated, reference;
the examples proved here)} \Cref{prop:ellcurves:torsred}.
\item \textbf{Mazur's torsion theorem.} $E(\Q)_{\mathrm{tors}}$ is $\Z/n$, $n\le10$ or $12$, or $\Z/2\times\Z/2n$, $n\le4$. \emph{(stated; reference)} \Cref{thm:ellcurves:mazur}.
\item \textbf{Isogeny classes} (definition and properties): an equivalence relation; isogenous curves have the same $V_\ell$, $a_p$, $L$-function, rank and conductor; Faltings's converse; at most $8$ curves
(Kenku). \emph{(the first two proved here, except the conductor; the rest stated, reference)} \Cref{def:ellcurves:isogclass}, \Cref{thm:ellcurves:isogclass}.
\item \textbf{Rational $\ell$-isogenies.} $K$-rational $\ell$-isogenies $\leftrightarrow$ $G_K$-stable lines of $E[\ell]$; $E[\ell]$ reducible iff a rational $\ell$-isogeny exists; then the two characters multiply to
$\bar\chi_\ell$. \emph{(proved here, given \Cref{thm:ellcurves:quotient})} \Cref{prop:ellcurves:isogline}.
\item \textbf{Mazur and Kenku.} Rational $\ell$-isogenies only for $\ell=2$, $3$, $5$, $7$, $11$, $13$, $17$, $19$, $37$, $43$, $67$, $163$; cyclic $n$-isogenies only for $n\le19$ or $n=21$, $25$, $27$, $37$, $43$, $67$, $163$,
with $19$, $27$, $43$, $67$, $163$ only for CM curves. \emph{(stated; reference)} \Cref{thm:ellcurves:mazurkenku}.
\item \textbf{$\Z/\ell\oplus\mu_\ell$.} With a rational point of order $\ell$ (odd), $\bar\rho=\sm1*0{\bar\chi_\ell}$; splitting $\Leftrightarrow$ image $\{\operatorname{diag}(1,d)\}$ $\Leftrightarrow$ two rational $\ell$-isogenies;
$a_p\equiv1+p\pmod\ell$. \emph{(proved here)} \Cref{prop:ellcurves:zlmul}.
\item \textbf{No point of order $9$} (the paper's Proposition prop:no9). A non-CM $E/\Q$ with $E[3]\simeq\Z/3\oplus\mu_3$ has no rational point of order $9$; so
$v_3(\#E(\Q)_{\mathrm{tors}})+[E[3]\simeq\Z/3\oplus\mu_3]\le2$. \emph{(proved here, given \Cref{thm:ellcurves:mazurkenku})} \Cref{prop:ellcurves:no27}.
\item \textbf{Isogenies prime to $\ell$.} If $E[\ell]$ is irreducible, every rational isogeny is $\psi\circ[a]$ with $\psi$ cyclic of degree prime to $\ell$; all curves of the class have isomorphic $T_\ell$, the same
$\ad T_\ell$ and $h_\ell$, and the same $t_\ell(E,q)$. \emph{(proved here; the last statement sketched, using \Cref{ch:reduction})} \Cref{lem:ellcurves:primetol}.
\item \textbf{Trivial isogeny classes.} Then $E[\ell]$ is irreducible for every $\ell$ and $E$ is optimal. \emph{(proved here)} \Cref{cor:ellcurves:trivialclass}.
\item \textbf{Quadratic twists} (definition and properties): $E^{(d)}\colon y^2=x^3+da_2x^2+d^2a_4x+d^3a_6$, for a general model $Y^2=X^3+db_2X^2+8d^2b_4X+16d^3b_6$; depends on $d$ modulo squares;
$\psi\colon E^{(d)}\to E$ over $K(\sqrt d)$ with $\psi^\sigma=[\chi_d(\sigma)]\psi$; the twisted Galois action. \emph{(proved here; the classification of all twists stated)} \Cref{def:ellcurves:twist},
\Cref{prop:ellcurves:twistbasic}.
\item \textbf{Invariants of twists.} $d^2c_4$, $d^3c_6$, $d^6\Delta$, same $j$; conversely equal $j\ne0,1728$ means a quadratic twist by $d=c_6'c_4/(c_6c_4')$; the twists by $q$ and $q^*$ differ by the twist by $-1$.
\emph{(proved here)} \Cref{prop:ellcurves:twistinv}.
\item \textbf{Galois representation of a twist.} $\rho_{E^{(d)},\ell}\simeq\rho_{E,\ell}\otimes\chi_d$; same $\ad$, same $\det$, same $E[2]$, irreducibility preserved. \emph{(proved here)} \Cref{thm:ellcurves:twistgal}.
\item \textbf{$a_p$ of a twist.} $a_p(E^{(d)})=\chi_d(p)a_p(E)$ for good $p\nmid D$. \emph{(proved here)} \Cref{cor:ellcurves:twistap}.
\item \textbf{The twist of good reduction.} At a type $\kod{I}_0^*$ prime $q$, $E'=E^{(q^*)}$ is good at $q$, $V_\ell E\simeq V_\ell E'\otimes\chi_{q^*}$, and $a_q(E')^2$ does not depend on the choice of twist. \emph{(proved here,
given \Cref{ch:reduction})} \Cref{prop:ellcurves:twistgood}.
\item \textbf{Composite twists} (definition): $\chi_S=\prod_{q\in S}\chi_{q^*}$, $E_S=E^{(d_S)}$, $d_S=\prod q^*$; $a_p(E_S)=\chi_S(p)a_p(E)$. \Cref{def:ellcurves:compositetwist}.
\item \textbf{What a twist changes.} $j$, $t_\ell$, $\ad\rho$ are twist-invariant; $e_\ell(E,q)$ is invariant only under twists unramified at $q$. \Cref{rem:ellcurves:twistinvariance}.
\item \textbf{The $L$-function} (definition and properties): $L(E,s)=\prod_{p\nmid N}(1-a_pp^{-s}+p^{1-2s})^{-1}\prod_{p\mid N}(1-a_pp^{-s})^{-1}$, absolutely convergent for $\operatorname{Re}s>3/2$, with multiplicative coefficients
and the recursion $a_{p^{k+1}}=a_pa_{p^k}-pa_{p^{k-1}}$; continuation and functional equation by modularity. \emph{(proved here, except (iv): stated, reference)} \Cref{def:ellcurves:Lfunction},
\Cref{prop:ellcurves:Lfunction}.
\item \textbf{Labels.} Cremona's $N\,c\,k$, curve $1$ optimal except in \curve{990h}; LMFDB labels differ. \Cref{conv:ellcurves:labels}, \Cref{ex:ellcurves:running}.
\end{enumerate}

\section*{Where the paper uses this chapter}

\begin{itemize}
\item \textbf{Section 1.1} ($V=V_\ell E$, $T=T_\ell E$, $\rho$, $\bar\rho$, the Frobenius conventions): \Cref{def:ellcurves:tate}, \Cref{prop:ellcurves:tate}, \Cref{def:ellcurves:rho}, \Cref{prop:ellcurves:rho},
\Cref{conv:ellcurves:frobenius}, \Cref{pl:ellcurves:section11}.
\item \textbf{Definition def:terms} ($j_E$, $v_q(j_E)$, the twist-invariance of $t_\ell$ and the failure for $e_\ell$): \Cref{prop:ellcurves:coordchange}(iii), \Cref{pl:ellcurves:terms},
\Cref{rem:ellcurves:twistinvariance}.
\item \textbf{The table of Section 1 and Lemma lem:euler} (the twist $E'$ of good reduction at a type $\kod{I}_0^*$ prime, $\alpha,\beta$ with $\alpha\beta=q$, $\alpha+\beta=a_q(E')$, $V=V'\otimes\chi_{q^*}$, $\ad V=\ad V'$):
\Cref{thm:ellcurves:charpoly}, \Cref{prop:ellcurves:twistFq}, \Cref{thm:ellcurves:twistgal}, \Cref{prop:ellcurves:twistgood}, \Cref{pl:ellcurves:table}, \Cref{pl:ellcurves:twist}.
\item \textbf{Lemma lem:twist (ii)} (the model and the invariants $(q^*)^2c_4$, $(q^*)^3c_6$, $(q^*)^6\Delta$ of the twist; $v_q(\Delta)$ modulo $12$): \Cref{def:ellcurves:twist}, \Cref{prop:ellcurves:twistinv},
\Cref{prop:ellcurves:coordchange}(iii), \Cref{warn:ellcurves:twists}.
\item \textbf{Section 1.4} (the twist-symmetrised weight $w_\ell(E,q)$, which compares $E$ with its twist by $\chi_{q^*}$; the paper writes this twist $E^{(q)}$, the book
$E^{(q^*)}$): \Cref{def:ellcurves:twist}(c), \Cref{prop:ellcurves:twistinv}(iii), \Cref{warn:ellcurves:twists}(a).
\item \textbf{Theorems thm:twist and thm:quant} ($E_S$, $\chi_S$, $a_p(f)=\chi_S(p)a_p(g)$, irreducibility of $\bar\rho\otimes\chi_S$): \Cref{def:ellcurves:compositetwist}, \Cref{thm:ellcurves:twistgal},
\Cref{cor:ellcurves:twistap}.
\item \textbf{Theorem thm:qual and the Example of Section 3} (ex:qual: twists of \curve{121a1}, \curve{539c1}, \curve{980i1}; labels; $a_p$): \Cref{cor:ellcurves:twistap}, \Cref{warn:ellcurves:twists},
\Cref{ex:ellcurves:aptable}, \Cref{ex:ellcurves:running}.
\item \textbf{Lemma lem:h0} (complex conjugation with eigenvalues $\pm1$, $\det\rho=\chi_\ell$, $\bar\chi_3=\chi_{-3}$; in the original version, CM): \Cref{cor:ellcurves:complex-conj}, \Cref{prop:ellcurves:rho},
\Cref{thm:ellcurves:cm-ap}, \Cref{warn:ellcurves:chi-l}.
\item \textbf{Proposition prop:ns} (the $3$-division polynomial, $\det\bar\rho_{E,3}=\chi_{-3}$, complex conjugation): \Cref{prop:ellcurves:psi3}, \Cref{ex:ellcurves:psi3-examples}, \Cref{cor:ellcurves:det},
\Cref{cor:ellcurves:complex-conj}, \Cref{pl:ellcurves:ns}.
\item \textbf{Conjecture conj:eis, Section sec:eis and Proposition prop:no9} (torsion, rational $\ell$-isogenies, $\Z/3\oplus\mu_3$, no point of order $9$, ``not CM''): \Cref{thm:ellcurves:mazur},
\Cref{prop:ellcurves:isogline}, \Cref{thm:ellcurves:mazurkenku}, \Cref{prop:ellcurves:zlmul}, \Cref{prop:ellcurves:no27}, \Cref{lem:ellcurves:cm-isog}, \Cref{pl:ellcurves:eis}, \Cref{pl:ellcurves:cm}.
\item \textbf{Section 1 (optimality and $\ell$; the area of the period lattice) and Section sec:oot (trivial isogeny classes)}: \Cref{lem:ellcurves:primetol}, \Cref{cor:ellcurves:trivialclass},
\Cref{thm:ellcurves:isogclass}, \Cref{def:ellcurves:periodlattice}, \Cref{ex:ellcurves:11a1-periods}, \Cref{pl:ellcurves:optimal}.
\item \textbf{Section 4.2 and Proposition prop:pullback} ($M_{f,\lambda}=H^1(E,\Z_\ell)=(T_\ell E)^\vee\simeq T_\ell E(-1)$, with the same adjoint): \Cref{cor:ellcurves:dual-T},
\Cref{pl:ellcurves:section11}.
\item \textbf{Lemma lem:old} ($X^2-a_p(h)X+p$, multiplicativity and recursion of the coefficients): \Cref{thm:ellcurves:frobQ}, \Cref{prop:ellcurves:Lfunction}.
\item \textbf{Section sec:two} (Calegari-Emerton; CM curves with $N\in\{27,32,49,243\}$): \Cref{thm:ellcurves:13j}, \Cref{pl:ellcurves:cm}.
\item \textbf{The appendix and Section sec:anomaly} (\texttt{ellinit}, \texttt{elldivpol}, \texttt{ns3}, \texttt{valuation(E.j, q)}, \curve{121a1}, \curve{726a1}): \Cref{def:ellcurves:weierstrass},
\Cref{prop:ellcurves:psi3}, \Cref{ex:ellcurves:psi3-examples}, \Cref{conv:ellcurves:labels}, \Cref{ex:ellcurves:running}.
\end{itemize}

\section{Weierstrass equations and their invariants}\label{sec:ellcurves:weierstrass}

Every elliptic curve in the paper is given by a Weierstrass equation, in the form that PARI's \texttt{ellinit} expects, and
every local quantity the paper reads off at a bad prime $q$ starts from the invariants $c_4$, $c_6$, $\Delta$ and $j$ of such an
equation. This section defines them, proves that a Weierstrass cubic is smooth exactly when $\Delta\ne0$, records how the
invariants change under a change of coordinates, and proves that $j$ classifies curves over an algebraically closed field.
The paper uses these facts in Definition def:terms (the valuation $v_q(j_E)$, which must not depend on the model), in the proof of
Lemma lem:twist (ii) (the invariants of a twisted model, and the rule that $v_q(\Delta)$ changes by multiples of $12$), and in its
appendix (\texttt{valuation(E.j, q)}).

\begin{convention}[conventions of this chapter]\label{conv:ellcurves:conventions}
\leavevmode
\begin{enumerate}[label=(\alph*)]
\item \emph{Fields and curves.} $K$ is a field, $\bar K$ a fixed algebraic closure, $K^{\mathrm{sep}}$ the separable closure in it and
$G_K=\Gal(K^{\mathrm{sep}}/K)$. An \emph{elliptic curve over $K$} is a pair $(E,O)$ of a smooth projective curve $E$ of genus one over $K$
and a point $O\in E(K)$. By \Cref{ch:curves} (the theorem that a genus-one curve with a rational point is a Weierstrass
cubic) every elliptic curve is isomorphic to a smooth Weierstrass cubic with $O\mapsto[0:1:0]$, and we always work with such a model.
Whenever a prime $\ell$ appears, $\ell\ne\operatorname{char}K$.
\item \emph{Over $\Q$.} ``$E/\Q$'' means a curve given by a Weierstrass equation with rational coefficients; from
\Cref{sec:ellcurves:overQ} on, by its global minimal model (\Cref{ch:reduction}), which is the model printed in Cremona's tables
\cite{Cremona97,ecdata}. Such models have $a_1,a_3\ne0$ in general (\curve{121a1}, \curve{726a1}, \curve{99a1}), so ``a model'' always
means a general Weierstrass equation.
\item \emph{Letters.} $p$, $q$, $\ell$ are primes; as in the paper, $q$ is usually a bad prime of $E/\Q$, $p$ any prime and $\ell$ the
coefficient prime. The number of elements of a finite field that need not be prime is written $\mathfrak q=p^r$
(\Cref{sec:ellcurves:finite}); for $\F_p$ we write $a_p$. In two places only, \Cref{thm:ellcurves:moduli}(iii) and the last sentence of
\Cref{ex:ellcurves:aptable}, $q=e^{2\pi i\tau}$ is the variable of $q$-expansions, as in \Cref{ch:modforms}; no prime $q$ occurs there.
\item \emph{Characters.} For $d\in K^\times$ and $\operatorname{char}K\ne2$, $\chi_d\colon G_K\to\{\pm1\}$ is the character of $K(\sqrt d)/K$:
$\sigma(\sqrt d)=\chi_d(\sigma)\sqrt d$. For $K=\Q$ and $d$ not a square, $\chi_d(\Frob_p)=\bigl(\frac Dp\bigr)$ (Kronecker symbol) for
$p\nmid D$, where $D$ is the discriminant of $\Q(\sqrt d)$ (\Cref{ch:rings,ch:ant}). For an odd prime $q$,
$q^*=(-1)^{(q-1)/2}q$ and $\chi_{q^*}$ is the quadratic character of conductor $q$. The paper's $\chi_q$ is written $\chi_{q^*}$; the
symbol $\chi_\ell$ with a letter subscript is the $\ell$-adic cyclotomic character, $\sigma(\zeta)=\zeta^{\chi_\ell(\sigma)}$ on
$\mu_{\ell^\infty}$, and $\bar\chi_\ell$ is its reduction modulo $\ell$; at a numerical prime we write $\chi_3^{\mathrm{cyc}}$, so that
$\bar\chi_3=\chi_{-3}$. For an integer $m$ prime to $\operatorname{char}K$, $\chi_m^{\mathrm{cyc}}\colon G_K\to(\Z/m)^\times$ is defined by
$\sigma(\zeta)=\zeta^{\chi_m^{\mathrm{cyc}}(\sigma)}$ for $\zeta\in\mu_m$.
\item \emph{Labels.} Cremona's labels are written as the paper writes them: \curve{11a1} is a curve, \curve{121c} an isogeny class. The
capital letters of Cremona's book and of \cite{ARS12} (\curve{99A1}) denote the same classes as the lower-case letters of the data files.
LMFDB labels are written \curve{11.a2} (\Cref{sec:ellcurves:labels}).
\item \emph{The Weil pairing} is Silverman's $e_m$ \cite[Section III.8]{SilvermanAEC}. Some books use its inverse; the only
statement of this chapter that depends on the choice is the comparison with $e^{2\pi i/m}$ over $\C$ in
\Cref{prop:ellcurves:cplxisog}(iv).
\item \emph{Duals and maps.} $M^\vee=\Hom(M,R)$ is the linear dual; $\varphi^*$ and $\varphi_*$ are pullback and pushforward along a
map $\varphi$ of curves; $\hat\varphi$ is the dual isogeny. Maps of curves are written $\varphi$, $\psi$; the letter $\psi_m$ is also used
for division polynomials (\Cref{sec:ellcurves:divpol}), which is never ambiguous in context.
\end{enumerate}
The Frobenius conventions are fixed in \Cref{conv:ellcurves:frobenius}.
\end{convention}

\begin{definition}[Weierstrass equations and their invariants]\label{def:ellcurves:weierstrass}
A \emph{Weierstrass equation} over $K$ is
\begin{equation}\label{eq:ellcurves:weierstrass}
y^2+a_1xy+a_3y=x^3+a_2x^2+a_4x+a_6,\qquad a_1,a_2,a_3,a_4,a_6\in K .
\end{equation}
The curve it defines is its projective closure $Y^2Z+a_1XYZ+a_3YZ^2=X^3+a_2X^2Z+a_4XZ^2+a_6Z^3$ in $\PP^2$, which has exactly one point
on the line $Z=0$, namely $O=[0:1:0]$. Put
\begin{align*}
b_2&=a_1^2+4a_2, & b_4&=2a_4+a_1a_3, & b_6&=a_3^2+4a_6,\\
b_8&=a_1^2a_6+4a_2a_6-a_1a_3a_4+a_2a_3^2-a_4^2, & c_4&=b_2^2-24b_4, & c_6&=-b_2^3+36b_2b_4-216b_6,
\end{align*}
\[
\Delta=-b_2^2b_8-8b_4^3-27b_6^2+9b_2b_4b_6,\qquad j=\frac{c_4^3}{\Delta}\ \text{ when }\Delta\ne0 .
\]
$\Delta$ is the \emph{discriminant} and $j$ the \emph{$j$-invariant}. The \emph{invariant differential} is
\[
\omega=\frac{dx}{2y+a_1x+a_3}=\frac{dy}{3x^2+2a_2x+a_4-a_1y}.
\]
The two expressions agree because differentiating \eqref{eq:ellcurves:weierstrass} gives
$(2y+a_1x+a_3)\,dy=(3x^2+2a_2x+a_4-a_1y)\,dx$. These are the conventions of \cite[Section III.1]{SilvermanAEC}, and PARI's
\texttt{ellinit} returns $a_1,\dots,a_6,b_2,b_4,b_6,b_8,c_4,c_6,\Delta,j$ in this order (\Cref{ch:computation}).
\end{definition}

\begin{lemma}[identities]\label{lem:ellcurves:identities}
For every Weierstrass equation:
\begin{enumerate}[label=(\roman*)]
\item $4b_8=b_2b_6-b_4^2$ and $1728\Delta=c_4^3-c_6^2$.
\item (Completing the square.) With $\eta=2y+a_1x+a_3$ the equation \eqref{eq:ellcurves:weierstrass} is equivalent, when $\operatorname{char}K\ne2$, to
\[
\eta^2=4x^3+b_2x^2+2b_4x+b_6 .
\]
The cubic on the right is written $F(x)$ below; its discriminant (as a polynomial in $x$) is $16\Delta$.
\item (Short form.) If $\operatorname{char}K\ne2,3$, the substitution $X=36x+3b_2$, $Y=108(2y+a_1x+a_3)$ transforms \eqref{eq:ellcurves:weierstrass} into
\[
Y^2=X^3-27c_4X-54c_6 .
\]
\end{enumerate}
\end{lemma}

\begin{proof}
All three are identities between polynomials in $a_1,\dots,a_6$ (and $x$, $y$). (i) Expanding $b_2b_6-b_4^2=(a_1^2+4a_2)(a_3^2+4a_6)-(2a_4+a_1a_3)^2$
gives $a_1^2a_3^2+4a_1^2a_6+4a_2a_3^2+16a_2a_6-4a_4^2-4a_1a_3a_4-a_1^2a_3^2$, which is $4b_8$. For the second identity expand both cubes as polynomials in $b_2,b_4,b_6$: the terms $b_2^6$ and $b_2^4b_4$ cancel, and
\begin{align*}
c_4^3-c_6^2&=432b_2^2b_4^2-432b_2^3b_6-13824b_4^3+15552b_2b_4b_6-46656b_6^2\\
&=1728\Bigl(-\tfrac14b_2^2(b_2b_6-b_4^2)-8b_4^3-27b_6^2+9b_2b_4b_6\Bigr),
\end{align*}
which is $1728\Delta$ by the first identity. (ii) $\eta^2=4y^2+4a_1xy+4a_3y+(a_1x+a_3)^2$, and replacing $4(y^2+a_1xy+a_3y)$ by $4(x^3+a_2x^2+a_4x+a_6)$ gives
$4x^3+(a_1^2+4a_2)x^2+(2a_1a_3+4a_4)x+(a_3^2+4a_6)$. Conversely, when $2$ is invertible, $y=(\eta-a_1x-a_3)/2$ recovers $y$. The
discriminant of $4x^3+b_2x^2+2b_4x+b_6$ is, by the formula for the discriminant of a cubic (\Cref{ch:galois}),
$b_2^2(2b_4)^2-4\cdot4(2b_4)^3-4b_2^3b_6-27\cdot16\,b_6^2+18\cdot4\,b_2(2b_4)b_6$, which equals $16\Delta$ after replacing $b_2^2b_4^2-b_2^3b_6$
by $-4b_2^2b_8$. (iii) Substituting $x=(X-3b_2)/36$ into $108^2F(x)$ and expanding gives $X^3-27c_4X-54c_6$, and $108^2\eta^2=Y^2$.
Each expansion was also checked by machine (PARI/GP, with $a_1,\dots,a_6$ as indeterminates).
\end{proof}

\begin{proposition}[smoothness]\label{prop:ellcurves:smooth}
Let $E$ be the curve defined by a Weierstrass equation over $K$.
\begin{enumerate}[label=(\roman*)]
\item The point $O$ is smooth, and $E$ is smooth if and only if $\Delta\ne0$.
\item If $\Delta=0$, $E$ has exactly one singular point. It is a \emph{node} (two distinct tangent lines) if $c_4\ne0$ and a \emph{cusp} (one
tangent line) if $c_4=0$.
\item A smooth Weierstrass cubic, with the point $O$, is an elliptic curve; every elliptic curve over $K$ is isomorphic over $K$ to one, with its
origin going to $O$.
\end{enumerate}
\end{proposition}

\begin{proof}
(i) That $O$ is smooth is checked in \Cref{ch:curves} (the partial derivative of the homogeneous equation with respect to $Z$ is $1$ at $O$).
Assume $\operatorname{char}K\ne2$. By \Cref{lem:ellcurves:identities}(ii) we may use the coordinates $(x,\eta)$, in which the equation is
$G=\eta^2-F(x)=0$. A point $(x_0,\eta_0)$ is singular iff $G=\partial G/\partial\eta=\partial G/\partial x=0$ there, that is, iff $\eta_0=0$ and
$F(x_0)=F'(x_0)=0$, iff $x_0$ is a double root of $F$. The cubic $F$ has a double root in $\bar K$ iff its discriminant $16\Delta$ vanishes. In
characteristic $2$ the statement is \cite[Section III.1]{SilvermanAEC}; we do not prove it.

(ii) Let $S$ be a singular point; it is the only one, for a line through two singular points would meet the cubic with multiplicity at least
$4$ and so lie in it (\Cref{ch:curves}, Bezout), while a Weierstrass cubic is irreducible (\Cref{ch:curves}). Move $S$ to $(0,0)$ by
$x\mapsto x+x_0$, $y\mapsto y+y_0$ (a change of coordinates with $u=1$, $s=0$, which by \Cref{prop:ellcurves:coordchange} below does not
change $c_4$ or $\Delta$). Since $(0,0)$ lies on the curve and both partial derivatives vanish there, $a_6=a_4=a_3=0$, and the equation is
$y^2+a_1xy-a_2x^2=x^3$. Now $b_4=b_6=b_8=0$, so $\Delta=0$ and $c_4=b_2^2=(a_1^2+4a_2)^2$. The tangent lines at $(0,0)$ are the factors of the
quadratic part $y^2+a_1xy-a_2x^2$, which are distinct iff its discriminant $a_1^2+4a_2=b_2$ is nonzero (in characteristic $2$: iff $a_1\ne0$,
iff $b_2=a_1^2\ne0$). So the point is a node iff $c_4\ne0$.

(iii) A smooth Weierstrass cubic has genus one and $O$ is a $K$-point (\Cref{ch:curves}); the converse is the theorem of \Cref{ch:curves}
quoted in \Cref{conv:ellcurves:conventions}(a).
\end{proof}

\begin{proposition}[changes of coordinates]\label{prop:ellcurves:coordchange}
\leavevmode
\begin{enumerate}[label=(\roman*)]
\item For $u\in K^\times$ and $r,s,t\in K$, the substitution
\[
x=u^2x'+r,\qquad y=u^3y'+su^2x'+t
\]
transforms a Weierstrass equation into a Weierstrass equation, with
\begin{align*}
ua_1'&=a_1+2s, & u^2a_2'&=a_2-sa_1+3r-s^2,\\
u^3a_3'&=a_3+ra_1+2t, & u^4a_4'&=a_4-sa_3+2ra_2-(t+rs)a_1+3r^2-2st,
\end{align*}
\begin{align*}
u^6a_6'&=a_6+ra_4+r^2a_2+r^3-ta_3-t^2-rta_1,\\
u^2b_2'&=b_2+12r,\qquad u^4b_4'=b_4+rb_2+6r^2,\qquad u^6b_6'=b_6+2rb_4+r^2b_2+4r^3,\\
u^8b_8'&=b_8+3rb_6+3r^2b_4+r^3b_2+3r^4,\qquad u^4c_4'=c_4,\qquad u^6c_6'=c_6,
\end{align*}
\[
u^{12}\Delta'=\Delta,\qquad j'=j,\qquad \omega'=u\,\omega .
\]
These substitutions form a group under composition.
\item Every isomorphism between two smooth Weierstrass cubics over $K$ that sends $O$ to $O$ is given by such a substitution.
\item Let $K=\Q$ and $q$ a prime. For any two models of $E$, $v_q(\Delta')-v_q(\Delta)\in12\Z$, and $v_q(j)=3v_q(c_4)-v_q(\Delta)$ is the same for
both. In particular $v_q(j_E)$ depends only on the isomorphism class of $E$ over $\Q$.
\end{enumerate}
\end{proposition}

\begin{proof}
(i) Substitute and compare coefficients: with $W(x,y)$ the left side minus the right side of \eqref{eq:ellcurves:weierstrass} and $W'$ the same
for the primed coefficients, one has $W(u^2x'+r,u^3y'+su^2x'+t)=u^6W'(x',y')$, which defines the $a_i'$ and is equivalent to the five formulas
displayed. The formulas for the $b_i'$, $c_4'$, $c_6'$, $\Delta'$ follow by substitution into \Cref{def:ellcurves:weierstrass}; all were checked
by machine with $a_i,u,r,s,t$ indeterminate. Then $j'=c_4'^3/\Delta'=u^{-12}c_4^3/(u^{-12}\Delta)=j$. Finally $dx=u^2dx'$ and
$2y+a_1x+a_3=u^3(2y'+a_1'x'+a_3')$ (from the formulas for $a_1'$ and $a_3'$), so $\omega=u^{-1}\omega'$. The composite of two substitutions is
the substitution with $u=u_1u_2$, $r=r_1+u_1^2r_2$, $s=s_1+u_1s_2$, $t=t_1+u_1^2s_1r_2+u_1^3t_2$, and $(u,r,s,t)$ has inverse
$(u^{-1},-ru^{-2},-su^{-1},(rs-t)u^{-3})$, as one checks by composing.

(ii) Let $\varphi\colon E\to E'$ be such an isomorphism, with coordinates $x,y$ on $E$ and $x',y'$ on $E'$. By \Cref{ch:curves}, $\mathcal L(2(O))$ has
basis $1,x$ and $\mathcal L(3(O))$ has basis $1,x,y$, and likewise on $E'$. Since $\varphi$ is an isomorphism with $\varphi(O)=O$, $\varphi^*x'$ has a
double pole at $O$ and no other pole, and $\varphi^*y'$ a triple pole; so $\varphi^*x'=\alpha x+\rho$ and $\varphi^*y'=\beta y+\gamma x+\tau$ with
$\alpha,\beta\in K^\times$, $\rho,\gamma,\tau\in K$ (all in $K$ when $\varphi$ is defined over $K$). Substituting these into the equation of $E'$ gives a relation
among $1,x,y,x^2,xy,y^2,x^3$ in $K(E)$. These seven functions lie in $\mathcal L(6(O))$, of dimension $6$, and have pole orders $0,2,3,4,5,6,6$ at $O$; the first six have
distinct pole orders and are therefore linearly independent. So the relations among the seven form a one-dimensional space, spanned by the
equation of $E$ (\Cref{ch:curves}, proof of the Weierstrass form). Comparing the coefficients of $y^2$ and $x^3$ gives $\beta^2=\alpha^3$. Put $u=\alpha/\beta$; then $u^2=\alpha^2/\beta^2=\alpha^{-1}$ and
$u^3=\beta^{-1}$, so $x=u^2\varphi^*x'-u^2\rho$ and $y=u^3\varphi^*y'-u^3\gamma x-u^3\tau$, which after substituting $x$ is of the stated form.

(iii) If both models have rational coefficients, $u\in\Q^\times$ and $v_q(\Delta')=v_q(\Delta)-12v_q(u)$, $v_q(c_4')=v_q(c_4)-4v_q(u)$; so
$3v_q(c_4')-v_q(\Delta')=3v_q(c_4)-v_q(\Delta)$, which is $v_q(j)$ (with $v_q(0)=+\infty$ when $c_4=0$).
\end{proof}

\begin{theorem}[the $j$-invariant]\label{thm:ellcurves:jinv}
\leavevmode
\begin{enumerate}[label=(\roman*)]
\item Two elliptic curves over $K$ are isomorphic over $\bar K$ if and only if they have the same $j$-invariant.
\item Every $j_0\in K$ is the $j$-invariant of an elliptic curve over $K$: one can take $y^2+y=x^3$ for $j_0=0$ (if $\operatorname{char}K\ne3$),
$y^2=x^3+x$ for $j_0=1728$ (if $\operatorname{char}K\ne2$), and
\[
y^2+xy=x^3-\frac{36}{j_0-1728}x-\frac1{j_0-1728}\qquad(j_0\ne0,1728),
\]
which has $c_4=j_0/(j_0-1728)$ and $\Delta=j_0^2/(j_0-1728)^3$.
\end{enumerate}
\end{theorem}

\begin{proof}
(i) ``Only if'' is \Cref{prop:ellcurves:coordchange}(i). For ``if'' we give the proof when $\operatorname{char}K\ne2,3$; the general case is
\cite[Appendix A]{SilvermanAEC}. By \Cref{lem:ellcurves:identities}(iii) (a change of coordinates with $u=1/6$) both curves have short models
$y^2=x^3-27c_4x-54c_6$, $y^2=x^3-27c_4'x-54c_6'$, and the change $(u,0,0,0)$ replaces $(c_4,c_6)$ by $(u^{-4}c_4,u^{-6}c_6)$. So it suffices to find
$u\in\bar K^\times$ with $c_4=u^4c_4'$ and $c_6=u^6c_6'$. If $j=j'=0$ then $c_4=c_4'=0$ and $c_6,c_6'\ne0$ (as $1728\Delta=-c_6^2$); take $u$ with
$u^6=c_6/c_6'$. If $j=j'=1728$ then $c_6=c_6'=0$ (as $1728\Delta=c_4^3-c_6^2$ and $j=1728$ means $c_4^3=1728\Delta$); take $u^4=c_4/c_4'$. Otherwise all of
$c_4,c_6,c_4',c_6'$ are nonzero. Choose $u$ with $u^4=c_4/c_4'$. Then $j=j'$ gives $\Delta/\Delta'=c_4^3/c_4'^3=u^{12}$, and
$c_6^2=c_4^3-1728\Delta=u^{12}(c_4'^3-1728\Delta')=u^{12}c_6'^2$, so $c_6=\pm u^6c_6'$; if the sign is $-$, replace $u$ by $iu$ with $i^2=-1$, which keeps $u^4$
and changes the sign of $u^6$.

(ii) The first two curves have $(c_4,\Delta)=(0,-27)$ and $(-48,-64)$. For the third, $a_1=1$, $a_2=a_3=0$, $a_4=-36/(j_0-1728)$, $a_6=-1/(j_0-1728)$, and
\Cref{def:ellcurves:weierstrass} gives the stated $c_4$ and $\Delta$, so $j=c_4^3/\Delta=j_0$; $\Delta\ne0$ because $j_0\ne0$.
\end{proof}

\begin{proposition}[automorphisms]\label{prop:ellcurves:aut}
Let $\operatorname{char}K\ne2,3$ and let $E$ be an elliptic curve over $K$. The group $\Aut_{\bar K}(E)$ of automorphisms of $E$ fixing $O$ is
cyclic of order $2$ if $j\ne0,1728$, of order $4$ if $j=1728$, and of order $6$ if $j=0$. On a short model $y^2=x^3+Ax+B$ the automorphisms are
$(x,y)\mapsto(u^2x,u^3y)$ with $u^4A=A$, $u^6B=B$; and $[-1](x,y)=(x,-y-a_1x-a_3)$ on a general model.
\end{proposition}

\begin{proof}
By \Cref{prop:ellcurves:coordchange}(ii) an automorphism of the short model is a substitution $(u,r,s,t)$ preserving the shape: $a_1'=0$
forces $s=0$, then $a_2'=0$ forces $r=0$, and then $a_3'=0$ forces $t=0$. So it is $(u,0,0,0)$ with $A=u^4A$, $B=u^6B$ (here $x=u^2x'$, so the map
$(x',y')\mapsto(x,y)$ is $(u^2x',u^3y')$). If $AB\ne0$ ($j\ne0,1728$) this means $u^2=1$; if $B=0$ ($j=1728$), $u^4=1$; if $A=0$ ($j=0$),
$u^6=1$. Distinct $u$ give distinct maps. The formula for $[-1]$ is the substitution $(u,r,s,t)=(-1,0,-a_1,-a_3)$, which sends $(x,y)$ to the other point
with the same $x$.
\end{proof}

\begin{example}[the invariants of the running examples]\label{ex:ellcurves:running-invariants}
The table lists the global minimal models of Cremona's tables \cite{ecdata} and their invariants, all recomputed from the $a$-invariants
with PARI/GP.
\begin{center}\small
\begin{tabular}{llllll}
\toprule
curve & $[a_1,a_2,a_3,a_4,a_6]$ & $c_4$ & $c_6$ & $\Delta$ & $j$\\
\midrule
\curve{11a1} & $[0,-1,1,-10,-20]$ & $2^4\cdot31$ & $20008$ & $-11^5$ & $-2^{12}31^3/11^5$\\
\curve{11a3} & $[0,-1,1,0,0]$ & $2^4$ & $-152$ & $-11$ & $-2^{12}/11$\\
\curve{20a1} & $[0,1,0,4,4]$ & $-2^4\cdot11$ & $-2368$ & $-2^8\cdot5^2$ & $2^4\cdot11^3/5^2$\\
\curve{99a1} & $[1,-1,1,-2,0]$ & $3^4$ & $135$ & $3^3\cdot11$ & $3^9/11$\\
\curve{121a1} & $[1,1,1,-30,-76]$ & $11\cdot131$ & $54703$ & $-11^2$ & $-11\cdot131^3$\\
\curve{121c1} & $[1,1,0,-2,-7]$ & $11^2$ & $5203$ & $-11^4$ & $-11^2$\\
\curve{121c2} & $[1,1,0,-3632,82757]$ & $11^3\cdot131$ & $-72809693$ & $-11^8$ & $-11\cdot131^3$\\
\curve{121d1} & $[0,-1,1,-40,-221]$ & $2^4\cdot11^2$ & $202312$ & $-11^7$ & $-2^{12}/11$\\
\curve{539c1} & $[1,0,1,170,-3237]$ & $-7^2\cdot167$ & $2808827$ & $-7^9\cdot11^2$ & $167^3/(7^3\cdot11^2)$\\
\curve{726a1} & $[1,1,0,-35,-51]$ & $5\cdot11\cdot31$ & $31339$ & $2^6\cdot3^3\cdot11^3$ & $5^3\cdot31^3/(2^6\cdot3^3)$\\
\curve{980i1} & $[0,0,0,1568,-72716]$ & $-2^9\cdot3\cdot7^2$ & $62826624$ & $-2^8\cdot5\cdot7^{11}$ & $2^{19}3^3/(5\cdot7^5)$\\
\curve{77c1} & $[1,1,0,4,11]$ & $-167$ & $-8189$ & $-7^3\cdot11^2$ & $167^3/(7^3\cdot11^2)$\\
\curve{140b1} & $[0,0,0,32,212]$ & $-2^9\cdot3$ & $-183168$ & $-2^8\cdot5\cdot7^5$ & $2^{19}3^3/(5\cdot7^5)$\\
\curve{539d1} & $[0,1,1,-16,-66]$ & $2^4\cdot7^2$ & $52136$ & $-7^6\cdot11$ & $-2^{12}/11$\\
\bottomrule
\end{tabular}
\end{center}
For instance, for \curve{11a1}: $b_2=-4$, $b_4=-20$, $b_6=-79$, $b_8=-21$, so $c_4=16+480=496$, $c_6=64+2880+17064=20008$ and
$\Delta=-16\cdot(-21)+64000-27\cdot6241+9\cdot(-4)(-20)(-79)=-161051=-11^5$. The valuations $v_q(j)$ feed the paper's term $t_\ell(E,q)$:
$v_7(j(\curve{539c1}))=-3$, so $t_3(\curve{539c1},7)=v_3(3)=1$; $v_7(j(\curve{980i1}))=-5$, so $t_5(\curve{980i1},7)=1$; and $v_{11}(j(\curve{121a1}))=1\ge0$, so
$t_\ell(\curve{121a1},11)=0$. Curves with the same $j$ in the table (\curve{539c1} and \curve{77c1}; \curve{980i1} and \curve{140b1}; \curve{11a3},
\curve{121d1} and \curve{539d1}; \curve{121a1} and \curve{121c2}) are quadratic twists of one another (\Cref{sec:ellcurves:twists}).
\end{example}

\begin{paperlink}\label{pl:ellcurves:terms}
Definition def:terms of the paper defines $t_\ell(E,q)=v_\ell(-v_q(j_E))$ when $v_q(j_E)<0$. That this is a property of $E$ and not of the
chosen equation is \Cref{prop:ellcurves:coordchange}(iii); the appendix computes it as \texttt{valuation(E.j, q)}. In the proof of Lemma
lem:twist (ii) the twisted model ``is minimal exactly when the result is less than $12$, otherwise $12$ is subtracted'': that $v_q(\Delta)$ can
only change by multiples of $12$ is \Cref{prop:ellcurves:coordchange}(iii), and the minimality criterion itself is \Cref{ch:reduction}.
\end{paperlink}

\section{The group law}\label{sec:ellcurves:grouplaw}

The points of an elliptic curve form an abelian group, and every group the paper writes down is a group of points or a subgroup of one:
$E(\Q)_{\mathrm{tors}}$ in Conjecture conj:eis, $E[\ell]$ in Section 1.1, $E(\F_q)$ behind the numbers $a_q$. This section gives the
chord-tangent construction with its explicit formulas, identifies it with the group $\Pic^0(E)$ of \Cref{ch:curves} (which gives
associativity for free), and proves the translation invariance of the invariant differential, which is what makes multiplication by $m$
separable when $\operatorname{char}K\nmid m$.

\begin{construction}[the chord-tangent law]\label{con:ellcurves:chordtangent}
Let $E$ be a smooth Weierstrass cubic over $K$ and $P,Q\in E(\bar K)$. The line through $P$ and $Q$ (the tangent line at $P$ if $P=Q$) meets
$E$ in exactly three points counted with multiplicity (\Cref{ch:curves}, Bezout); let $R$ be the third. The line through $R$ and $O$ meets $E$ in a
third point, which is defined to be $P+Q$. For an affine point $R=(x,y)$, the line through $R$ and $O=[0:1:0]$ is the vertical line through $x$,
and its third point is
\[
-R=(x,-y-a_1x-a_3).
\]
Hence $P+Q=-R$, and $P+Q+R=O$ whenever $P$, $Q$, $R$ are the three intersections of a line with $E$ (the line at infinity $Z=0$ meets $E$ only at
$O$, with multiplicity $3$).
\end{construction}

\begin{proposition}[the addition formulas]\label{prop:ellcurves:addition}
Let $P_1=(x_1,y_1)$ and $P_2=(x_2,y_2)$ be affine points of $E$.
\begin{enumerate}[label=(\roman*)]
\item If $x_1=x_2$ and $y_2=-y_1-a_1x_1-a_3$, then $P_1+P_2=O$.
\item Otherwise put, if $x_1\ne x_2$,
\[
\lambda=\frac{y_2-y_1}{x_2-x_1},\qquad\nu=\frac{y_1x_2-y_2x_1}{x_2-x_1},
\]
and, if $P_1=P_2$ (so $2y_1+a_1x_1+a_3\ne0$),
\[
\lambda=\frac{3x_1^2+2a_2x_1+a_4-a_1y_1}{2y_1+a_1x_1+a_3},\qquad\nu=\frac{-x_1^3+a_4x_1+2a_6-a_3y_1}{2y_1+a_1x_1+a_3}.
\]
Then $P_1+P_2=(x_3,y_3)$ with $x_3=\lambda^2+a_1\lambda-a_2-x_1-x_2$ and $y_3=-(\lambda+a_1)x_3-\nu-a_3$.
\item (Duplication.) If $2P\ne O$, then with $F$ as in \Cref{lem:ellcurves:identities},
\[
x(2P)=\frac{x^4-b_4x^2-2b_6x-b_8}{4x^3+b_2x^2+2b_4x+b_6}\quad(x=x(P)).
\]
\item The points $P\ne O$ with $2P=O$ are those with $2y+a_1x+a_3=0$; when $\operatorname{char}K\ne2$ their $x$-coordinates are the three distinct roots of
$F(x)=4x^3+b_2x^2+2b_4x+b_6$, so $E[2]$ has four elements.
\end{enumerate}
\end{proposition}

\begin{proof}
(i) $P_2=-P_1$ by the formula for $-R$. (ii) The line $y=\lambda x+\nu$ passes through $P_1$ and $P_2$ (for $P_1=P_2$, $\lambda$ is the slope of the
tangent, by implicit differentiation of \eqref{eq:ellcurves:weierstrass}, and $\nu=y_1-\lambda x_1$, which after using the equation of $E$ at $P_1$ is
the displayed expression). Substituting $y=\lambda x+\nu$ into \eqref{eq:ellcurves:weierstrass} gives a monic cubic in $x$,
\[
x^3+(a_2-\lambda^2-a_1\lambda)x^2+\cdots=0,
\]
whose roots are $x_1,x_2,x_3$ (with multiplicity). Comparing the $x^2$-coefficients gives $x_1+x_2+x_3=\lambda^2+a_1\lambda-a_2$. The third point is
$R=(x_3,\lambda x_3+\nu)$, and $P_1+P_2=-R$. (iii) Substitute the tangent slope into $x_3=\lambda^2+a_1\lambda-a_2-2x_1$; clearing the denominator
$(2y_1+a_1x_1+a_3)^2=F(x_1)$ (\Cref{lem:ellcurves:identities}(ii)) and eliminating $y_1$ with the equation of $E$ leaves the stated quotient (an
identity checked by machine; on \curve{11a1} at $P=(5,5)$ it gives $x(2P)=1936/121=16$). (iv) $2P=O$ means $P=-P$, i.e.\ $y=-y-a_1x-a_3$. In
characteristic not $2$ this is $\eta=0$, so $F(x)=0$ by \Cref{lem:ellcurves:identities}(ii); conversely each root of $F$ gives exactly one point with
$\eta=0$. The roots are distinct because $\operatorname{disc}F=16\Delta\ne0$.
\end{proof}

\begin{proposition}[the chord-tangent law is the law of $\Pic^0$]\label{prop:ellcurves:pic0}
Let $\kappa\colon E(\bar K)\to\Pic^0(E)$, $P\mapsto$ the class of $(P)-(O)$, be the bijection of \Cref{ch:curves}. Then $\kappa(P+Q)=\kappa(P)+\kappa(Q)$ for
the chord-tangent law of \Cref{con:ellcurves:chordtangent}.
\end{proposition}

\begin{proof}
Let $P,Q$ be given, $R$ the third point of the line $L$ through them, and $P+Q=-R$ the third point of the line $L'$ through $R$ and $O$. Write $L$ and
$L'$ as linear forms in $X,Y,Z$. The rational functions $L/Z$ and $L'/Z$ have divisors
\[
\operatorname{div}(L/Z)=(P)+(Q)+(R)-3(O),\qquad\operatorname{div}(L'/Z)=(R)+(-R)+(O)-3(O)=(R)+(-R)-2(O),
\]
because $L$ meets $E$ in $P,Q,R$, and $Z$ meets $E$ only at $O$ with multiplicity $3$ (if $L$ or $L'$ is the line $Z=0$ itself the corresponding
statement is trivial). Hence $\operatorname{div}(L/L')=(P)+(Q)-(P+Q)-(O)$, so $(P)-(O)+(Q)-(O)\sim(P+Q)-(O)$, which is the claim.
\end{proof}

\begin{theorem}[the group law]\label{thm:ellcurves:grouplaw}
\leavevmode
\begin{enumerate}[label=(\roman*)]
\item $(E(\bar K),+)$ is an abelian group with identity $O$ and inverse $P\mapsto-P$.
\item For every field $L$ with $K\subseteq L\subseteq\bar K$, $E(L)$ is a subgroup, and $\sigma(P+Q)=\sigma P+\sigma Q$ for every automorphism $\sigma$ of $\bar K$
over $K$.
\item The maps $E\times E\to E$, $(P,Q)\mapsto P+Q$, and $E\to E$, $P\mapsto-P$, are morphisms (given by rational functions with coefficients in
$K$); in particular each translation $\tau_Q\colon P\mapsto P+Q$ is an isomorphism of curves, defined over $K$ if $Q\in E(K)$.
\end{enumerate}
\end{theorem}

\begin{proof}
(i) Commutativity is built into the construction (the line through $P$ and $Q$ is the line through $Q$ and $P$); $P+O=P$ because the line
through $P$ and $O$ meets $E$ again at $-P$, and the line through $-P$ and $O$ meets it at $P$; and $P+(-P)=O$ because the vertical line through
$P$ meets $E$ in $P$, $-P$, $O$, and the tangent at $O$ is $Z=0$, which meets $E$ only at $O$. Associativity holds because $\kappa$ is a bijection onto
the group $\Pic^0(E)$ that carries $+$ to the group law there (\Cref{prop:ellcurves:pic0}). A direct verification from the formulas is possible but
long; see \cite[Section III.2]{SilvermanAEC} and the references there. (ii) The formulas of \Cref{prop:ellcurves:addition} have coefficients in $K$, so
they commute with $\sigma$ and preserve coordinates in $L$. (iii) The formulas of \Cref{prop:ellcurves:addition} are rational functions; that they
define a morphism also at the excluded points (where $x_1=x_2$, or a point is $O$) is \cite[Section III.3]{SilvermanAEC}; for $\tau_Q$ with $Q$ fixed it
follows from \Cref{ch:curves} (a rational map from a smooth curve to a projective variety is a morphism), and $\tau_{-Q}$ is its inverse.
\end{proof}

\begin{corollary}[principal divisors]\label{cor:ellcurves:principal}
A divisor $D=\sum n_P(P)$ on $E$ is principal if and only if $\deg D=0$ and $\sum[n_P]P=O$ in $E(\bar K)$.
\end{corollary}

\begin{proof}
For $\deg D=0$, $D=\sum n_P\bigl((P)-(O)\bigr)$ has class $\sum n_P\kappa(P)=\kappa\bigl(\sum[n_P]P\bigr)$ by \Cref{prop:ellcurves:pic0}. Since $\kappa$ is
injective, the class is $0$ iff $\sum[n_P]P=O$. Principal divisors have degree $0$ (\Cref{ch:curves}).
\end{proof}

\begin{proposition}[the invariant differential]\label{prop:ellcurves:invdiff}
Let $\omega$ be the invariant differential of $E$.
\begin{enumerate}[label=(\roman*)]
\item $\omega$ is holomorphic and has no zeros: $\operatorname{div}(\omega)=0$.
\item $\tau_Q^*\omega=\omega$ for every $Q\in E(\bar K)$.
\item Let $E'$ be an elliptic curve and $\varphi,\psi\colon E'\to E$ morphisms with $\varphi(O)=\psi(O)=O$ (isogenies, \Cref{def:ellcurves:isogeny}). Then
$(\varphi+\psi)^*\omega=\varphi^*\omega+\psi^*\omega$, where $\varphi+\psi$ is the pointwise sum and the pullback by the zero map is $0$. Hence $[m]^*\omega=m\,\omega$ for every $m\in\Z$.
\end{enumerate}
\end{proposition}

\begin{proof}
(i) is \Cref{ch:curves} (the divisor of $dx/(2y+a_1x+a_3)$ on a smooth Weierstrass cubic). For (ii) and (iii) put $\eta=2y+a_1x+a_3$ and
$\xi=3x^2+2a_2x+a_4-a_1y$. Differentiating \eqref{eq:ellcurves:weierstrass} gives $\eta\,dy=\xi\,dx$, and $\omega=dx/\eta$. For a morphism $\varphi$ from a
smooth curve to $E$ we put $\varphi^*\omega=0$ if $\varphi$ is constant; if $\varphi$ is nonconstant, $\varphi^*\eta\ne0$ (pullback is injective on function fields) and
$\varphi^*\omega=d(\varphi^*x)/\varphi^*\eta$.

\emph{Step 1 (an identity).} Let $x_1,y_1,x_2,y_2$ be indeterminates, $W(x,y)=y^2+a_1xy+a_3y-x^3-a_2x^2-a_4x-a_6$, and let $X_3$, $Y_3$ be the
expressions for $x_3$, $y_3$ in the chord case of \Cref{prop:ellcurves:addition}(ii), so that their denominators are powers of $x_2-x_1$. Put
$H_3=2Y_3+a_1X_3+a_3$, and $\eta_i$, $\xi_i$ for $\eta$, $\xi$ evaluated at $(x_i,y_i)$. Then for $i=1,2$
\begin{equation}\label{eq:ellcurves:invdiff-identity}
\eta_i\,\frac{\partial X_3}{\partial x_i}+\xi_i\,\frac{\partial X_3}{\partial y_i}=H_3
\end{equation}
in the ring $\Z[a_1,a_2,a_3,a_4,a_6][x_1,y_1,x_2,y_2,(x_1-x_2)^{-1}]$ modulo the ideal generated by $W(x_1,y_1)$ and $W(x_2,y_2)$. Since $W$ is monic of degree $2$
in $y$, the difference of the two sides can be reduced, by subtracting multiples of $W(x_1,y_1)$ and $W(x_2,y_2)$, to an expression of degree at most $1$ in $y_1$ and
in $y_2$, without new denominators. We carried out this reduction with PARI/GP, with $a_1,\dots,a_6$ indeterminate, and it gives $0$. So
\eqref{eq:ellcurves:invdiff-identity} holds after any specialisation of the variables into a field in which $W(x_1,y_1)=W(x_2,y_2)=0$ and $x_1\ne x_2$, in every
characteristic.

\emph{Step 2 (sums of morphisms).} Let $C$ be a smooth curve over $\bar K$ and $\varphi,\psi\colon C\to E$ morphisms, neither of them the constant map to $O$, such that
$x\circ\varphi\ne x\circ\psi$ and the pointwise sum $\varphi+\psi$ is a nonconstant morphism. We show $(\varphi+\psi)^*\omega=\varphi^*\omega+\psi^*\omega$. Put
$x_1=\varphi^*x$, $y_1=\varphi^*y$, $x_2=\psi^*x$, $y_2=\psi^*y$ in $\bar K(C)$; then $W(x_i,y_i)=0$ and $x_1-x_2\ne0$. The sets $\varphi^{-1}(O)$ and $\psi^{-1}(O)$
are finite, and $x_1-x_2$ has finitely many zeros, so at all but finitely many $P\in C$ the points $\varphi(P)$, $\psi(P)$ are affine with different $x$-coordinates,
and there \Cref{prop:ellcurves:addition}(ii) gives $(\varphi+\psi)(P)=\bigl(X_3,Y_3\bigr)\bigl(x_1(P),y_1(P),x_2(P),y_2(P)\bigr)$. Two rational functions on $C$ that agree at
infinitely many points are equal, so $(\varphi+\psi)^*x=X_3(x_1,y_1,x_2,y_2)$ and $(\varphi+\psi)^*\eta=H_3(x_1,y_1,x_2,y_2)$, the latter nonzero because $\varphi+\psi$
is nonconstant. By the chain rule for the derivation $d$ of $\bar K(C)$,
\[
d\bigl((\varphi+\psi)^*x\bigr)=\sum_{i=1,2}\Bigl(\frac{\partial X_3}{\partial x_i}\,dx_i+\frac{\partial X_3}{\partial y_i}\,dy_i\Bigr).
\]
If $\varphi$ is nonconstant, $\eta_1=\varphi^*\eta\ne0$ and pulling back $\eta\,dy=\xi\,dx$ gives $dy_1=(\xi_1/\eta_1)\,dx_1$; so the term $i=1$ is
$\bigl(\eta_1\partial X_3/\partial x_1+\xi_1\partial X_3/\partial y_1\bigr)dx_1/\eta_1=H_3\,\varphi^*\omega$ by \eqref{eq:ellcurves:invdiff-identity}. If $\varphi$ is constant,
$dx_1=dy_1=0$ and the term is $0=H_3\,\varphi^*\omega$. The same holds for $i=2$. Dividing by $H_3=(\varphi+\psi)^*\eta$ gives the claim.

(ii) For $Q=O$ there is nothing to prove. For $Q\ne O$ apply Step 2 with $C=E$, $\varphi$ the identity and $\psi$ the constant map to $Q$: then $\varphi+\psi=\tau_Q$ is an
isomorphism (\Cref{thm:ellcurves:grouplaw}(iii)) and $x\ne x(Q)$, so $\tau_Q^*\omega=\omega+0=\omega$.

(iii) First, $x\circ[-1]=x$ and $y\circ[-1]=-y-a_1x-a_3$ (\Cref{con:ellcurves:chordtangent}), so $[-1]^*\eta=-\eta$ and $[-1]^*\omega=-\omega$. The maps $\varphi\pm\psi$ are
morphisms sending $O$ to $O$ (\Cref{thm:ellcurves:grouplaw}(iii)), and such a morphism is either the zero map or nonconstant, in which case its fibres, in particular its
kernel, are finite (\Cref{ch:curves}). If $\varphi=0$ or $\psi=0$ there is nothing to prove. If $\psi=-\varphi$, then $(\varphi+\psi)^*\omega=0$ and
$\psi^*\omega=\varphi^*[-1]^*\omega=-\varphi^*\omega$. If $\psi=\varphi$: the set $E[2]$ is finite, since its points other than $O$ are the zeros of the nonzero function $\eta$
(\Cref{prop:ellcurves:addition}(iv)); choose $Q$ with $2Q\ne O$. Then $\tau_Q+\tau_{-Q}=[2]$ is nonconstant, and $x\circ\tau_Q\ne x\circ\tau_{-Q}$: otherwise for every
$P$ the points $P+Q$ and $P-Q$ would be equal or opposite, so $2Q=O$ or $2P=O$, and $E(\bar K)=E[2]$ would be finite. Step 2 and (ii) give
$[2]^*\omega=\tau_Q^*\omega+\tau_{-Q}^*\omega=2\omega$, hence $(2\varphi)^*\omega=([2]\circ\varphi)^*\omega=\varphi^*[2]^*\omega=2\varphi^*\omega$. In the remaining case
$\varphi,\psi\ne0$ and $\psi\ne\pm\varphi$, so $\varphi+\psi$ is nonconstant; and $x\circ\varphi\ne x\circ\psi$, for otherwise $\varphi(P)=\pm\psi(P)$ for every $P$ (two points
with the same $x$-coordinate in $\PP^1$ are equal or opposite), and $E'(\bar K)$, which is infinite, would be the union of the finite kernels of $\varphi-\psi$ and
$\varphi+\psi$. Step 2 applies.

Finally $[0]^*\omega=0$, $[m+1]^*\omega=[m]^*\omega+\omega$ for $m\ge0$ by what was just proved, so $[m]^*\omega=m\,\omega$ for $m\ge0$, and
$[-m]^*\omega=([-1]\circ[m])^*\omega=[m]^*(-\omega)=-m\,\omega$. Silverman \cite[Section III.5]{SilvermanAEC} argues differently, proving (iii) through differentials on
$E\times E$; the explicit identity of Step 1 replaces that argument here.
\end{proof}

\begin{example}[torsion points on the running examples]\label{ex:ellcurves:torsion-points}
All computed with \Cref{prop:ellcurves:addition} (and checked with PARI).
\begin{enumerate}[label=(\alph*)]
\item On \curve{11a1}, $P=(5,5)$: $2P=(16,-61)$, $3P=(16,60)$, $4P=(5,-6)$, $5P=O$. Indeed $3P=-2P$ because $-(16,-61)=(16,61-1)=(16,60)$.
\item On \curve{20a1} $=[0,1,0,4,4]$: $(-1,0)$ has order $2$ (it satisfies $2y=0$), $(0,2)$ has order $3$, and $(4,10)$ has order $6$, with $2(4,10)=(0,2)$ and
$3(4,10)=(-1,0)$.
\item On \curve{14a1} $=[1,0,1,4,-6]$: $(2,-5)$ has order $3$ (with $2(2,-5)=(2,2)=-(2,-5)$), and $(9,23)$ has order $6$, with $2(9,23)=(2,2)$, $3(9,23)=(1,-1)$.
\item On \curve{27a1} $=[0,0,1,0,-7]$: $(3,4)$ has order $3$, $2(3,4)=(3,-5)$.
\item On \curve{99a1}, $(0,0)$ has infinite order: its multiples $(3,2)$, $(-8/9,13/27)$, $(105/64,-897/512)$ have growing denominators, which
torsion points cannot have (\Cref{prop:ellcurves:torsred}); its canonical height is $0.3025\ldots>0$.
\end{enumerate}
\end{example}
\section{Isogenies}\label{sec:ellcurves:isogenies}

The paper works with isogeny classes: the $X_0(N)$-optimal curve is one curve of a class, a rational $\ell$-isogeny is the Eisenstein
case of Conjecture conj:eis, and Proposition prop:no9 is an argument about a cyclic isogeny of degree $27$. This section defines
isogenies and their degree, states the existence of quotients by finite subgroups and of the dual isogeny, and proves the consequences
the book needs: $\deg[m]=m^2$, $E[m]\simeq(\Z/m)^2$ when $\operatorname{char}K\nmid m$, and the factorisation of an isogeny into a
multiplication map and a cyclic isogeny.

\begin{definition}[isogenies]\label{def:ellcurves:isogeny}
Let $E_1$, $E_2$ be elliptic curves over $K$. An \emph{isogeny} $\varphi\colon E_1\to E_2$ is a morphism of curves with $\varphi(O)=O$. It is
\emph{defined over $K$} if it is given by rational functions with coefficients in $K$; for perfect $K$ this holds iff $\varphi(\sigma P)=\sigma\varphi(P)$
for all $\sigma\in G_K$ and $P\in E_1(\bar K)$ (\Cref{ch:curves}). Its \emph{degree} is $\deg\varphi=[\bar K(E_1):\varphi^*\bar K(E_2)]$ if $\varphi$ is
nonconstant and $\deg0=0$; $\deg_s\varphi$ and $\deg_i\varphi$ are the separable and inseparable degrees, and $\varphi$ is \emph{separable} if
$\deg_i\varphi=1$ (\Cref{ch:curves}). $\Hom(E_1,E_2)$ is the set of isogenies (over $\bar K$ unless stated otherwise), $\End(E)=\Hom(E,E)$, and
$[m]\colon E\to E$ is multiplication by $m\in\Z$. $E_1$ and $E_2$ are \emph{isogenous over $K$} if there is a nonconstant isogeny $E_1\to E_2$ defined
over $K$. For $m\ge1$, $E[m]=\{P\in E(\bar K):mP=O\}$.
\end{definition}

\begin{proposition}[first properties]\label{prop:ellcurves:isogeny-basic}
\leavevmode
\begin{enumerate}[label=(\roman*)]
\item An isogeny is either constant (the zero map) or surjective on $\bar K$-points.
\item $\Hom(E_1,E_2)$ is an abelian group under $(\varphi+\psi)(P)=\varphi(P)+\psi(P)$, $\End(E)$ is a ring under composition, and
$\deg(\psi\circ\varphi)=\deg\psi\cdot\deg\varphi$.
\item For $m\ne0$, $[m]\ne0$. Hence $\Hom(E_1,E_2)$ is torsion-free, and $\End(E)$ has characteristic $0$.
\end{enumerate}
\end{proposition}

\begin{proof}
(i) is \Cref{ch:curves} (a morphism of curves is constant or surjective). (ii) $\varphi+\psi$ is the composite of $(\varphi,\psi)\colon E_1\to E_2\times E_2$
with the addition morphism (\Cref{thm:ellcurves:grouplaw}(iii)), hence a morphism, and it sends $O$ to $O$. The group and ring axioms hold pointwise.
For nonconstant $\varphi$, $\psi$ the degree is multiplicative in the tower $\bar K(E_1)\supseteq\varphi^*\bar K(E_2)\supseteq\varphi^*\psi^*\bar K(E_3)$; if one of
them is constant both sides are $0$.
(iii) Assume first $\operatorname{char}K\ne2$ and write $m=\pm2^am'$ with $m'$ odd; as degrees multiply, it suffices to show $[2]\ne0$ and $[m']\ne0$. By
\Cref{prop:ellcurves:addition}(iv) the points with $2P=O$ are $O$ and three points with $2y+a_1x+a_3=0$; since $E(\bar K)$ is infinite, $[2]\ne0$. Let $T$ be one of
the three points of order $2$ (a root of the cubic $F$). Then $m'T=T\ne O$, so $[m']\ne0$. In characteristic $2$ the same argument works with $[3]$ in place of $[2]$
and a point of order $3$ in place of $T$ (the $x$-coordinates of the points of order $3$ are roots of $\psi_3$, of degree $4$ in characteristic $2$; \Cref{prop:ellcurves:psi3}(i));
see \cite[Section III.4]{SilvermanAEC} for the details. Finally, if $n\varphi=0$ with $n\ne0$, then $[n]\circ\varphi=0$; a nonzero $\varphi$ is surjective by (i), so $[n]$ would vanish on $E_2(\bar K)$, which we have
excluded. Hence $\varphi=0$.
\end{proof}

\begin{theorem}[isogenies are homomorphisms; kernels]\label{thm:ellcurves:isog-hom}
Let $\varphi\colon E_1\to E_2$ be a nonconstant isogeny.
\begin{enumerate}[label=(\roman*)]
\item $\varphi(P+Q)=\varphi(P)+\varphi(Q)$ for all $P,Q$.
\item $\ker\varphi$ is finite of order $\deg_s\varphi$. Every fibre $\varphi^{-1}(Q)$ has $\deg_s\varphi$ points, each with ramification index $\deg_i\varphi$.
\item If $\varphi$ is separable, then $\#\ker\varphi=\deg\varphi$, $\varphi$ is unramified, and $\bar K(E_1)/\varphi^*\bar K(E_2)$ is a Galois extension whose group is
$\{\tau_T^*:T\in\ker\varphi\}\simeq\ker\varphi$.
\end{enumerate}
\end{theorem}

\begin{proof}[Sketch of proof]
(i) The pushforward $\varphi_*$ on divisors maps principal divisors to principal divisors (\Cref{ch:curves}), so it induces a homomorphism
$\Pic^0(E_1)\to\Pic^0(E_2)$, and $\varphi_*\bigl((P)-(O)\bigr)=(\varphi P)-(O)$. Through the bijections $\kappa$ of \Cref{prop:ellcurves:pic0} this says that
$\varphi$ is a homomorphism. (ii) For $T\in\ker\varphi$, $\varphi\circ\tau_T=\varphi$ by (i), so $\tau_T$ permutes each fibre and preserves ramification indices:
$e_\varphi(P+T)=e_\varphi(P)$. Hence all points of a fibre $\varphi^{-1}(Q)=P+\ker\varphi$ have the same index $e$, and the degree formula of \Cref{ch:curves} gives
$\#\ker\varphi\cdot e=\deg\varphi$. Since all but finitely many fibres have exactly $\deg_s\varphi$ points (\Cref{ch:curves}), $\#\ker\varphi=\deg_s\varphi$ and
$e=\deg_i\varphi$. (iii) If $\varphi$ is separable, $e=1$ everywhere and $\#\ker\varphi=\deg\varphi$. The $\deg\varphi$ automorphisms $\tau_T^*$, $T\in\ker\varphi$, of
$\bar K(E_1)$ fix $\varphi^*\bar K(E_2)$ (as $\varphi\circ\tau_T=\varphi$) and are distinct; an extension of degree $n$ with $n$ automorphisms over the base is
Galois with that group (\Cref{ch:galois}). This is \cite[Section III.4]{SilvermanAEC}. In characteristic $0$ every morphism of genus-one curves is
separable and unramified (\Cref{ch:curves}, Riemann-Hurwitz), which gives (iii) directly.
\end{proof}

\begin{theorem}[quotients by finite subgroups]\label{thm:ellcurves:quotient}
\leavevmode
\begin{enumerate}[label=(\roman*)]
\item Let $H\subset E(\bar K)$ be a finite subgroup. There are an elliptic curve $E'$ and a separable isogeny $\varphi\colon E\to E'$ with $\ker\varphi=H$;
the pair is unique up to a unique isomorphism of $E'$. If $H$ is stable as a set under $G_K$, then $E'$ and $\varphi$ can be taken to be defined over $K$.
We write $E'=E/H$.
\item (Factorisation.) Let $\varphi\colon E\to E_1$ and $\psi\colon E\to E_2$ be nonconstant isogenies with $\psi$ separable and $\ker\psi\subseteq\ker\varphi$.
Then there is a unique isogeny $\lambda\colon E_2\to E_1$ with $\varphi=\lambda\circ\psi$; if $\varphi$ and $\psi$ are defined over $K$, so is $\lambda$, and
$\ker\lambda=\psi(\ker\varphi)$.
\end{enumerate}
\end{theorem}

\noindent We do not prove this; see \cite[Section III.4 and Exercise 3.13]{SilvermanAEC}. The idea of (ii): by \Cref{thm:ellcurves:isog-hom}(iii),
$\varphi^*\bar K(E_1)$ is fixed by the translations by $\ker\psi$, hence contained in $\psi^*\bar K(E_2)$, and an inclusion of function fields is a morphism
(\Cref{ch:curves}); $\lambda$ is unique because $\psi$ is surjective, and $\lambda^\sigma\circ\psi=\varphi^\sigma=\varphi$ gives $\lambda^\sigma=\lambda$. The quotient in (i)
is the curve whose function field is the fixed field of the translations by $H$. V\'elu's formulas \cite{Velu71} give an explicit Weierstrass equation
for $E/H$ and for $\varphi$ in terms of the coordinates of the points of $H$ (PARI's \texttt{ellisogeny}).

\begin{theorem}[the dual isogeny]\label{thm:ellcurves:dual}
Let $\varphi\colon E_1\to E_2$ be a nonconstant isogeny of degree $m$. There is a unique isogeny $\hat\varphi\colon E_2\to E_1$ with
$\hat\varphi\circ\varphi=[m]_{E_1}$. Moreover:
\begin{enumerate}[label=(\roman*)]
\item $\varphi\circ\hat\varphi=[m]_{E_2}$, $\deg\hat\varphi=\deg\varphi$ and $\hat{\hat\varphi}=\varphi$;
\item $\widehat{\psi\circ\varphi}=\hat\varphi\circ\hat\psi$, $\widehat{\varphi+\psi}=\hat\varphi+\hat\psi$ (with $\hat0=0$), and $\widehat{[m]}=[m]$;
\item if $\varphi$ is defined over $K$, so is $\hat\varphi$;
\item $\deg\colon\Hom(E_1,E_2)\to\Z$ is a positive definite quadratic form: $\deg(n\varphi)=n^2\deg\varphi$, $\deg\varphi\ge0$ with equality only for $\varphi=0$, and
$\langle\varphi,\psi\rangle=\deg(\varphi+\psi)-\deg\varphi-\deg\psi$ is bilinear.
\end{enumerate}
\end{theorem}

\begin{proof}[Sketch of proof]
$\hat\varphi$ is the composite $E_2\xrightarrow{\kappa}\Pic^0(E_2)\xrightarrow{\varphi^*}\Pic^0(E_1)\xrightarrow{\kappa^{-1}}E_1$: since
$\varphi^*\bigl((Q)\bigr)=\sum_{\varphi(P)=Q}e_\varphi(P)(P)$ and all points of a fibre $P+\ker\varphi$ have index $\deg_i\varphi$ (\Cref{thm:ellcurves:isog-hom}(ii)),
$\varphi^*\bigl((\varphi P)-(O)\bigr)=\deg_i\varphi\sum_{T\in\ker\varphi}\bigl((P+T)-(T)\bigr)$, whose image under $\kappa^{-1}$ is $[\deg_i\varphi\cdot\#\ker\varphi]P=[\deg\varphi]P$ by
\Cref{cor:ellcurves:principal}. Uniqueness: if $\psi\circ\varphi=\psi'\circ\varphi$ then $(\psi-\psi')\circ\varphi=0$ and $\psi=\psi'$ as $\varphi$ is surjective. The
additivity $\widehat{\varphi+\psi}=\hat\varphi+\hat\psi$ is the deepest property; see \cite[Section III.6]{SilvermanAEC}. The rest follows formally: from
$\varphi\circ\hat\varphi\circ\varphi=\varphi\circ[m]=[m]\circ\varphi$ and surjectivity of $\varphi$ one gets $\varphi\circ\hat\varphi=[m]$; then $\deg\varphi\deg\hat\varphi=\deg[m]=m^2$
(\Cref{cor:ellcurves:degm}(i), whose proof uses only $\widehat{[m]}=[m]$, itself a consequence of additivity and $\hat{[1]}=[1]$), so $\deg\hat\varphi=m$. For (iv),
$[\deg(\varphi+\psi)]=\widehat{(\varphi+\psi)}\circ(\varphi+\psi)=[\deg\varphi]+[\deg\psi]+\hat\psi\circ\varphi+\hat\varphi\circ\psi$, so $\langle\varphi,\psi\rangle$ corresponds to
$\hat\psi\circ\varphi+\hat\varphi\circ\psi$, which is visibly biadditive; and $\End(E_1)$ is torsion-free (\Cref{prop:ellcurves:isogeny-basic}(iii)), so the
identification of integers with multiplication maps is injective.
\end{proof}

\begin{corollary}[multiplication by $m$ and the $m$-torsion]\label{cor:ellcurves:degm}
Let $m\ge1$.
\begin{enumerate}[label=(\roman*)]
\item $\deg[m]=m^2$.
\item If $\operatorname{char}K\nmid m$, then $[m]$ is separable and $E[m]\simeq\Z/m\times\Z/m$.
\item If $\operatorname{char}K=p>0$, then either $E[p^e]=0$ for all $e\ge1$, or $E[p^e]\simeq\Z/p^e$ for all $e\ge1$.
\end{enumerate}
\end{corollary}

\begin{proof}
(i) By \Cref{thm:ellcurves:dual}, $\widehat{[m]}=[m]$ and $\widehat{[m]}\circ[m]=[\deg[m]]$, so $[m^2]=[\deg[m]]$ and $[m^2-\deg[m]]=0$; as $\End(E)$ is torsion-free
(\Cref{prop:ellcurves:isogeny-basic}(iii)), $\deg[m]=m^2$. (ii) By \Cref{prop:ellcurves:invdiff}(iii), $[m]^*\omega=m\omega\ne0$, so $[m]$ is separable (\Cref{ch:curves}:
a morphism is separable iff it pulls back a nonzero differential to a nonzero one). By \Cref{thm:ellcurves:isog-hom}(iii), $\#E[d]=d^2$ for every $d\mid m$. Let
$A=E[m]$; it is a finite abelian group killed by $m$. For each prime $p\mid m$, $A[p]=E[p]$ has $p^2$ elements, so the $p$-primary part of $A$ is
$\Z/p^a\times\Z/p^b$ with exactly two factors (\Cref{ch:groups}, structure theorem), and $a,b\le v_p(m)$. Since $\#A=m^2$, $a+b=2v_p(m)$, so $a=b=v_p(m)$, and
$A\simeq(\Z/m)^2$. (iii) We do not prove this; see \cite[Section III.6]{SilvermanAEC}. Over $\C$, \Cref{prop:ellcurves:cplxisog}(ii) gives a second proof of (i)
and (ii).
\end{proof}

\begin{proposition}[cyclic factorisation]\label{prop:ellcurves:cyclicfactor}
Let $\operatorname{char}K=0$ and let $\varphi\colon E\to E'$ be a nonconstant isogeny.
\begin{enumerate}[label=(\roman*)]
\item $\ker\varphi\simeq\Z/a\times\Z/ab$ for unique integers $a,b\ge1$, and $a$ is the largest integer with $E[a]\subseteq\ker\varphi$.
\item $\varphi=\psi\circ[a]$ for a unique isogeny $\psi\colon E\to E'$, and $\ker\psi=a\cdot\ker\varphi\simeq\Z/b$ is cyclic. If $\varphi$ is defined over $K$, so is $\psi$.
\item An isogeny over $K$ with cyclic kernel of order $n$ is called a \emph{cyclic $n$-isogeny}. Its kernel $C$ contains exactly one subgroup $C_d$ of each order
$d\mid n$, and $C_d$ is $G_K$-stable. Consequently a cyclic $n$-isogeny over $K$ is a composite of cyclic isogenies over $K$ of prime degrees, and if $\ell\mid n$
then $E$ has a $K$-rational cyclic $\ell$-isogeny, with kernel $C_\ell$.
\end{enumerate}
\end{proposition}

\begin{proof}
(i) $H=\ker\varphi$ is a finite subgroup of $E[N]\simeq(\Z/N)^2$ for $N=\#H$ (\Cref{cor:ellcurves:degm}(ii)). For each prime $p$, $H[p]\subseteq E[p]\simeq(\Z/p)^2$,
so the $p$-primary part of $H$ has at most two cyclic factors, and $H\simeq\Z/d_1\times\Z/d_2$ with $d_1\mid d_2$ (invariant factors, \Cref{ch:groups}); put $a=d_1$,
$b=d_2/d_1$. Then $H[a]\simeq(\Z/a)^2$ has $a^2$ elements and lies in $E[a]$, which has $a^2$ elements, so $E[a]\subseteq H$; and if $E[a']\subseteq H$ then
$(\Z/a')^2\hookrightarrow\Z/d_1\times\Z/d_2$, which forces $a'\mid d_1$.
(ii) $[a]$ is separable with kernel $E[a]\subseteq\ker\varphi$, so \Cref{thm:ellcurves:quotient}(ii) gives $\psi$ with $\varphi=\psi\circ[a]$, unique, defined over $K$ if
$\varphi$ is, and with kernel $[a](\ker\varphi)=a\cdot(\Z/a\times\Z/ab)\simeq\Z/b$.
(iii) A cyclic group of order $n$ has exactly one subgroup of each order $d\mid n$. For $\sigma\in G_K$, $\sigma(C_d)$ is a subgroup of order $d$ of $\sigma(C)=C$, so
$\sigma(C_d)=C_d$. For a prime $p\mid n$, the isogeny $E\to E/C_p$ is defined over $K$ (\Cref{thm:ellcurves:quotient}(i)), and $\varphi$ factors through it
(\Cref{thm:ellcurves:quotient}(ii)) as $E\to E/C_p\to E'$, the second map having cyclic kernel $C/C_p$ of order $n/p$, $G_K$-stable. Induction on $n$ finishes the proof.
\end{proof}

\begin{warning}[isogenous over which field]\label{warn:ellcurves:isogenous}
In the paper and in this book ``isogenous'' and ``isogeny class'' for curves over $\Q$ always mean over $\Q$. Two curves over $\Q$ can be isogenous over
$\Qbar$ without being isogenous over $\Q$: a curve and its quadratic twist by $\Q(\sqrt d)$ are isomorphic over $\Q(\sqrt d)$
(\Cref{def:ellcurves:twist}), but in general not isogenous over $\Q$: \curve{11a1} has $a_2=-2$, while its twist \curve{121d2} lies in the class of \curve{121d1}, which has
$a_2=2$ (\Cref{ex:ellcurves:aptable}), so the two are not isogenous over $\Q$ (\Cref{thm:ellcurves:isogclass}(ii)).
\end{warning}

\section{Division polynomials}\label{sec:ellcurves:divpol}

The paper decides whether the mod $3$ image of $E$ is the normaliser $N_s(\F_3)$ of a split Cartan subgroup by factoring the $3$-division polynomial
(Proposition prop:ns and the function \texttt{ns3} of its appendix, which calls \texttt{elldivpol(E,3)}). This section defines the division polynomials,
proves their degrees, and proves that the four roots of $\psi_3$ are the $x$-coordinates of the four lines of $E[3]$, Galois-equivariantly. The group-theoretic
side (which orbit patterns occur for which subgroups of $\GL_2(\F_3)$) is in \Cref{ch:groups,ch:galois}; the deduction of Proposition prop:ns is in
\Cref{ch:galrep}.

\begin{definition}[division polynomials]\label{def:ellcurves:divpol}
Let $R=\Z[a_1,a_2,a_3,a_4,a_6]$, $F(x)=4x^3+b_2x^2+2b_4x+b_6\in R[x]$, and let $A=R[x][\psi_2]/(\psi_2^2-F(x))$, a free $R[x]$-module with basis $1,\psi_2$. Define
$\psi_m\in A$ by
\begin{gather*}
\psi_0=0,\quad\psi_1=1,\quad\psi_2=\psi_2,\quad\psi_3=3x^4+b_2x^3+3b_4x^2+3b_6x+b_8,\\
\psi_4=\psi_2\bigl(2x^6+b_2x^5+5b_4x^4+10b_6x^3+10b_8x^2+(b_2b_8-b_4b_6)x+b_4b_8-b_6^2\bigr),
\end{gather*}
and recursively, for $m\ge2$ and $m\ge3$ respectively,
\[
\psi_{2m+1}=\psi_{m+2}\psi_m^3-\psi_{m-1}\psi_{m+1}^3,\qquad\psi_2\psi_{2m}=\psi_m\bigl(\psi_{m+2}\psi_{m-1}^2-\psi_{m-2}\psi_{m+1}^2\bigr),
\]
where the second formula defines $\psi_{2m}$ because its right side is $\psi_2$ times an element of $A$ (\Cref{prop:ellcurves:divpol-deg}), and
multiplication by $\psi_2$ is injective on $A$. On a curve $E$ we evaluate $\psi_m$ at the coefficients of $E$ and put $\psi_2=2y+a_1x+a_3$, which satisfies
$\psi_2^2=F(x)$ on $E$ (\Cref{lem:ellcurves:identities}(ii)); so $\psi_m$ becomes a function on $E$. These are the conventions of \cite[Exercise 3.7]{SilvermanAEC}
and of PARI; \cite[Section 3.2]{WashingtonEC} uses short models.
\end{definition}

\begin{proposition}[degrees]\label{prop:ellcurves:divpol-deg}
For odd $m\ge1$, $\psi_m\in R[x]$ is a polynomial of degree $(m^2-1)/2$ with leading coefficient $m$. For even $m\ge2$, $\psi_m=\psi_2f_m$ with $f_m\in R[x]$ of
degree $(m^2-4)/2$ and leading coefficient $m/2$.
\end{proposition}

\begin{proof}
For $m\le4$ this is read off the definition ($f_2=1$; $f_4$ has degree $6$ and leading coefficient $2$). Assume it for all indices below the one considered;
we use $\psi_2^2=F$, of degree $3$ with leading coefficient $4$. Write $\psi_k$ for odd $k$ and $\psi_k=\psi_2f_k$ for even $k$.

\emph{$\psi_{2n+1}$, $n\ge2$.} If $n$ is even, $\psi_{n+2}\psi_n^3=\psi_2^4f_{n+2}f_n^3=F^2f_{n+2}f_n^3$ has degree $6+\frac{(n+2)^2-4}2+3\frac{n^2-4}2=2n^2+2n$ and
leading coefficient $16\cdot\frac{n+2}2\cdot\frac{n^3}8=(n+2)n^3$, while $\psi_{n-1}\psi_{n+1}^3$ (odd indices) has degree $\frac{(n-1)^2-1}2+3\frac{(n+1)^2-1}2=2n^2+2n$ and
leading coefficient $(n-1)(n+1)^3$. If $n$ is odd the roles are exchanged: $\psi_{n+2}\psi_n^3$ has leading coefficient $(n+2)n^3$ and
$\psi_{n-1}\psi_{n+1}^3=F^2f_{n-1}f_{n+1}^3$ has $16\cdot\frac{n-1}2\cdot\frac{(n+1)^3}8=(n-1)(n+1)^3$, both of degree $2n^2+2n$. In both cases $\psi_{2n+1}\in R[x]$ has degree at most
$2n^2+2n=\frac{(2n+1)^2-1}2$ and coefficient of that power
\[
(n+2)n^3-(n-1)(n+1)^3=2n+1 .
\]
\emph{$\psi_{2n}$, $n\ge3$.} If $n$ is odd, $\psi_{n\pm1}^2=Ff_{n\pm1}^2$, so the bracket is $F\bigl(\psi_{n+2}f_{n-1}^2-\psi_{n-2}f_{n+1}^2\bigr)$ and
$\psi_{2n}=\psi_2\cdot\psi_n\bigl(\psi_{n+2}f_{n-1}^2-\psi_{n-2}f_{n+1}^2\bigr)$. The two products in the bracket have degree $\frac{(n\pm2)^2-1}2+(n\mp1)^2-4=\frac{3n^2-3}2$ and
leading coefficients $(n+2)\frac{(n-1)^2}4$ and $(n-2)\frac{(n+1)^2}4$, whose difference is $\frac14\bigl((n+2)(n-1)^2-(n-2)(n+1)^2\bigr)=\frac44=1$. So $f_{2n}$ has degree
$\frac{n^2-1}2+\frac{3n^2-3}2=2n^2-2=\frac{(2n)^2-4}2$ and leading coefficient $n$. If $n$ is even, $\psi_n=\psi_2f_n$ and $\psi_{n\pm2}=\psi_2f_{n\pm2}$, so
$\psi_{2n}=\psi_2\cdot f_n\bigl(f_{n+2}\psi_{n-1}^2-f_{n-2}\psi_{n+1}^2\bigr)$; the bracket has degree $\frac{(n\pm2)^2-4}2+(n\mp1)^2-1=\frac{3n^2}2$ and leading coefficient
$\frac12\bigl((n+2)(n-1)^2-(n-2)(n+1)^2\bigr)=2$, so $f_{2n}$ has degree $\frac{n^2-4}2+\frac{3n^2}2=2n^2-2$ and leading coefficient $\frac n2\cdot2=n$. (For $n=4$ the term
$f_{n-2}=f_2=1$; the case $n-2=0$ does not arise since $n\ge3$.) In particular the right side of the even recursion is always $\psi_2$ times an element of
$R[x]$, as \Cref{def:ellcurves:divpol} requires.
\end{proof}

For \curve{11a1} the degrees $0,3,4,9,12,19,24,33,40$ and leading coefficients $1,4,3,8,5,12,7,16,9$ of PARI's \texttt{elldivpol(E,m)}, $m=1,\dots,9$, agree with
\Cref{prop:ellcurves:divpol-deg}: for even $m$ PARI returns $\psi_2\psi_m=Ff_m$, of degree $3+(m^2-4)/2$ and leading coefficient $4\cdot m/2$.

\begin{theorem}[multiplication by $m$]\label{thm:ellcurves:multm}
Let $m\ge2$ with $\operatorname{char}K\nmid m$, and $P=(x,y)\in E(\bar K)$. Then $mP=O$ if and only if $\psi_m(P)=0$, and if $mP\ne O$,
\[
x(mP)=x-\frac{\psi_{m-1}\psi_{m+1}}{\psi_m^2}(P).
\]
There is a similar formula for $y(mP)$ with denominator $\psi_m^3$. Equivalently,
\[
\operatorname{div}(\psi_m)=\sum_{T\in E[m]}(T)-m^2(O).
\]
\end{theorem}

\begin{proof}[Sketch of proof]
The formulas are proved by induction on $m$ from the addition formulas of \Cref{prop:ellcurves:addition}, which is what the recursions of
\Cref{def:ellcurves:divpol} encode; see \cite[Section 3.2]{WashingtonEC} for short models and \cite[Exercise 3.7]{SilvermanAEC} in general. For $m=2$ the
formula is the duplication formula, as $x-\psi_1\psi_3/\psi_2^2=x-\psi_3/F=x(2P)$ by \Cref{prop:ellcurves:psi3}(i); for $m=3$ the statement ``$3P=O$ iff
$\psi_3(P)=0$'' is \Cref{prop:ellcurves:psi3}(i). The form with divisors: $\psi_m$ has poles only at $O$, of order $m^2-1$ (by \Cref{prop:ellcurves:divpol-deg},
as $\ord_Ox=-2$, $\ord_Oy=-3$), and its zeros are the $m^2-1$ points of $E[m]\smallsetminus\{O\}$ (\Cref{cor:ellcurves:degm}(ii)), each simple.
\end{proof}

\begin{corollary}[roots of $\psi_m$ for odd $m$]\label{cor:ellcurves:divpol-roots}
Let $m$ be odd with $\operatorname{char}K\nmid m$. The roots of the polynomial $\psi_m(x)$ in $\bar K$ are exactly the $x$-coordinates of the points of
$E[m]\smallsetminus\{O\}$; they are simple, there are $(m^2-1)/2$ of them, and $P\mapsto x(P)$ is a $G_K$-equivariant bijection
$(E[m]\smallsetminus\{O\})/\{\pm1\}\to\{\text{roots of }\psi_m\}$.
\end{corollary}

\begin{proof}
By \Cref{thm:ellcurves:multm} every $x(P)$, $P\in E[m]\smallsetminus\{O\}$, is a root. Two points have the same $x$-coordinate iff they are equal or opposite (the
vertical line meets $E$ in $P$, $-P$, $O$). No point of $E[m]\smallsetminus\{O\}$ has order $2$ since $m$ is odd, so $P\ne-P$, and the number of distinct values $x(P)$ is
$(m^2-1)/2$ (\Cref{cor:ellcurves:degm}(ii)). This is the degree of $\psi_m$, whose leading coefficient $m$ is nonzero in $K$ (\Cref{prop:ellcurves:divpol-deg});
so these values are all the roots, each simple. Equivariance: $\psi_m$ has coefficients in $K$ and $x(\sigma P)=\sigma(x(P))$.
\end{proof}

\begin{proposition}[the $3$-division polynomial]\label{prop:ellcurves:psi3}
Let $\psi_3=3x^4+b_2x^3+3b_4x^2+3b_6x+b_8$.
\begin{enumerate}[label=(\roman*)]
\item As rational functions on $E$, $x(2P)-x(P)=-\psi_3(x)/\psi_2(P)^2$. Hence for $P\ne O$ with $2P\ne O$: $3P=O$ if and only if $\psi_3(x(P))=0$.
\item Let $\operatorname{char}K\ne2,3$. Then $E[3]$ is a two-dimensional $\F_3$-vector space with four lines $\{O,P,-P\}$, and $\{O,P,-P\}\mapsto x(P)$ is a $G_K$-equivariant
bijection from the set of lines of $E[3]$ onto the set of the four (distinct) roots of $\psi_3$. Consequently the degrees of the irreducible factors of $\psi_3$ over
$K$ are the sizes of the orbits of $\bar\rho_{E,3}(G_K)$ on $\PP^1(\F_3)$ (\Cref{def:ellcurves:rho}); $\psi_3$ has a root in $K$ iff $E[3]$ has a $G_K$-stable line, iff $E$
has a $3$-isogeny defined over $K$ (\Cref{prop:ellcurves:isogline}).
\item Let $K=\Q$. Then $\disc(\psi_3)\in-3\cdot(\Q^\times)^2$, and the splitting field of $\psi_3$ contains $\Q(\sqrt{-3})$.
\end{enumerate}
\end{proposition}

\begin{proof}
(i) By the duplication formula (\Cref{prop:ellcurves:addition}(iii)),
\begin{align*}
x(2P)-x&=\frac{x^4-b_4x^2-2b_6x-b_8-x(4x^3+b_2x^2+2b_4x+b_6)}{F(x)}\\
&=-\frac{3x^4+b_2x^3+3b_4x^2+3b_6x+b_8}{F(x)},
\end{align*}
and $F(x)=\psi_2^2$ on $E$. If $P\ne O$ and $2P\ne O$, then $3P=O$ iff $2P=-P$ iff $x(2P)=x(P)$ (as $2P\ne P$), iff $\psi_3(x(P))=0$.
(ii) $E[3]\simeq(\Z/3)^2$ by \Cref{cor:ellcurves:degm}(ii); its lines are the four subgroups of order $3$, each of the form $\{O,P,-P\}$. By \Cref{cor:ellcurves:divpol-roots} with
$m=3$, $x$ gives a $G_K$-equivariant bijection from the lines onto the four distinct roots of $\psi_3$. For a separable polynomial over $K$, the irreducible factors over $K$
correspond to the $G_K$-orbits on its roots (\Cref{ch:galois}); after a choice of basis, $G_K$ acts on the lines of $E[3]$ through $\bar\rho_{E,3}$ and the action of
$\GL_2(\F_3)$ on $\PP^1(\F_3)$. A root in $K$ is a fixed root, that is, a $G_K$-stable line.
(iii) The roots $r_1,\dots,r_4$ are distinct, so $\delta=\prod_{i<j}(r_i-r_j)\ne0$, $\disc\psi_3=3^6\delta^2$, and $\sigma(\delta)=\sgn(\pi_\sigma)\delta$ where $\pi_\sigma$ is the permutation of
the roots induced by $\sigma\in G_\Q$. By (ii), $\pi_\sigma$ is the permutation of the four lines of $E[3]$ induced by $\bar\rho_{E,3}(\sigma)$, whose sign is $\det\bar\rho_{E,3}(\sigma)\in\{\pm1\}=\F_3^\times$
(\Cref{ch:groups}); and $\det\bar\rho_{E,3}=\bar\chi_3=\chi_{-3}$ (\Cref{cor:ellcurves:det}). So $\sigma$ fixes $\delta$ iff $\sigma$ fixes $\sqrt{-3}$, which means
$\Q(\delta)=\Q(\sqrt{-3})$ and $\delta^2\in-3(\Q^\times)^2$.
\end{proof}

\begin{example}[factorisations of $\psi_3$]\label{ex:ellcurves:psi3-examples}
The following factorisations over $\Q$ and discriminants were computed with PARI (\texttt{elldivpol(E,3)}, \texttt{factor}, \texttt{poldisc}); the image of $\bar\rho_{E,3}$
is taken from Cremona's \texttt{galrep} file (labels of \cite{Sutherland16}; no label means surjective; their meaning is \Cref{ch:galrep}).
\begin{center}\small
\begin{tabular}{llll}
\toprule
curve & $\psi_3$ over $\Q$ & $\disc\psi_3$ & image at $3$\\
\midrule
\curve{11a1} & $3x^4-4x^3-60x^2-237x-21$ (irreducible) & $-3^3\cdot11^{10}$ & $\GL_2(\F_3)$\\
\curve{121a1} & $3x^4+5x^3-177x^2-909x-1249$ (irreducible) & $-3^3\cdot11^4$ & $\GL_2(\F_3)$\\
\curve{726a1} & $(x^2-x-74)(3x^2+8x+20)$ & $-2^{12}3^9 11^6$ & \texttt{3Ns}\\
\curve{338d1} & $(x^2+29x+1858)(3x^2-82x-172)$ & $-2^63^313^{18}$ & \texttt{3Ns}\\
\curve{6912d1} & $3(x^2-4x-2)(x^2+4x+6)$ & $-2^{18}3^9$ & \texttt{3Ns}\\
\curve{20a1} & $x(3x^3+4x^2+24x+48)$ & $-2^{16}3^35^4$ & \texttt{3B.1.1}\\
\curve{14a1} & $(x-2)(3x+1)(x^2+2x+13)$ & $-2^{12}3^37^6$ & \texttt{3Cs.1.1}\\
\curve{27a1} (CM) & $3x(x-3)(x^2+3x+9)$ & $-3^{21}$ & \texttt{3Cs.1.1}\\
\curve{121b1} (CM) & $(x^2+3x-6)(3x^2-13x+15)$ & $-3^3\cdot11^6$ & $N_s(\F_3)$, not listed\\
\bottomrule
\end{tabular}
\end{center}
Every discriminant is $-3$ times a square, as \Cref{prop:ellcurves:psi3}(iii) says. For \curve{726a1} the quadratic factors have discriminants $297=3^3\cdot11$ and
$-176=-2^4\cdot11$, so their roots generate $\Q(\sqrt{33})$ and $\Q(\sqrt{-11})$, and the third quadratic subfield of the compositum is $\Q(\sqrt{-3})$. For
\curve{20a1} the root $x=0$ is the $x$-coordinate of the rational point $(0,2)$ of order $3$ (\Cref{ex:ellcurves:torsion-points}); the cubic factor is irreducible, so the
other three lines form one orbit. For \curve{14a1} the root $2$ is the $x$-coordinate of $(2,-5)$, of order $3$; the root $-1/3$ gives a second rational line, whose
points $(-1/3,y)$ satisfy $y^2+\frac23y+\frac{199}{27}=0$, of discriminant $-\frac{784}{27}=-3\cdot\bigl(\frac{28}9\bigr)^2$, so they are defined over $\Q(\sqrt{-3})$ and
not over $\Q$: this line is $\mu_3$ (\Cref{prop:ellcurves:zlmul}). The curve \curve{121b1} has complex multiplication by $\Q(\sqrt{-11})$, in which $3$ splits. Its
$\psi_3$ has no root in $\Q$, so $E[3]$ is irreducible (\Cref{prop:ellcurves:psi3}(ii)); the image has orbits of sizes $2,2$ on $\PP^1(\F_3)$ and contains the image of complex
conjugation (determinant $-1$, trace $0$), so by the classification recalled in \Cref{pl:ellcurves:ns} it is exactly $N_s(\F_3)$, the normaliser of the split Cartan subgroup
attached to the two primes above $3$ in $\Q(\sqrt{-11})$. The \texttt{galrep} file gives no label at $3$ for this curve: for CM curves it often omits labels (it
does give \texttt{3Cs.1.1} for the CM curves \curve{27a1} and \curve{27a3}).
\end{example}

\begin{paperlink}\label{pl:ellcurves:ns}
Proposition prop:ns says that, for $E[3]$ irreducible, the $3$-division polynomial factors as a product of two irreducible quadratics when the image is $N_s(\F_3)$
and is irreducible otherwise. The bridge between the image and the polynomial is \Cref{prop:ellcurves:psi3}(ii): the orbits of $N_s(\F_3)$ on the four lines have sizes
$2,2$, and the other irreducible subgroups of $\GL_2(\F_3)$ containing an element of determinant $-1$ and trace $0$ (complex conjugation,
\Cref{cor:ellcurves:complex-conj}), namely the normaliser of a nonsplit Cartan subgroup and $\GL_2(\F_3)$ itself, act transitively (\Cref{ch:groups}; the
deduction is in \Cref{ch:galrep}). The paper's function \texttt{ns3} applies this to \texttt{elldivpol(E,3)}; a reducible $E[3]$ gives a linear factor. The CM curve \curve{121b1} of the
table has the pattern $2+2$ and image exactly $N_s(\F_3)$, consistently with Proposition prop:ns; the paper excludes CM curves from its data, and \texttt{ns3}
would classify \curve{121b1} correctly. Section sec:anomaly uses \curve{726a1}. The proof of Proposition prop:ns itself is \Cref{ch:galrep}.
\end{paperlink}
\section[The Weil pairing, Tate modules, Galois representations]{The Weil pairing, the Tate module and the Galois representations}\label{sec:ellcurves:weil}

Section 1.1 of the paper begins with $V=V_\ell E$, $T=T_\ell E$, $\rho=\rho_{E,\ell}$ and its reduction $\bar\rho$; every later statement is about these
objects or about the adjoint $\ad T$ built from them. This section constructs them, proves that $\det\rho_{E,\ell}$ is the cyclotomic character
(through the Weil pairing), identifies $T_\ell E$ with its own dual twisted by $\Z_\ell(1)$, and fixes the Frobenius conventions. The paper uses these facts in
Section 1.1, in Section 4.2 (the module $T^D$ and the identification $M_{f,\lambda}=H^1(E,\Z_\ell)$ after Proposition prop:pullback), in Lemma lem:h0 (complex
conjugation, $\det\rho=\chi_\ell$), and in Proposition prop:ns ($\det\bar\rho_{E,3}=\chi_{-3}$).

\begin{theorem}[the Weil pairing]\label{thm:ellcurves:weil}
Let $m\ge1$ with $\operatorname{char}K\nmid m$. There is a pairing $e_m\colon E[m]\times E[m]\to\mu_m$ with the following properties.
\begin{enumerate}[label=(\roman*)]
\item Bilinear: $e_m(S_1+S_2,T)=e_m(S_1,T)e_m(S_2,T)$ and $e_m(S,T_1+T_2)=e_m(S,T_1)e_m(S,T_2)$.
\item Alternating: $e_m(T,T)=1$; hence $e_m(S,T)=e_m(T,S)^{-1}$.
\item Nondegenerate: if $e_m(S,T)=1$ for all $S$, then $T=O$.
\item Galois-equivariant: $e_m(\sigma S,\sigma T)=\sigma\bigl(e_m(S,T)\bigr)$ for $\sigma\in G_K$.
\item Compatible: $e_{mm'}(S,T)=e_m(m'S,T)$ for $S\in E[mm']$ and $T\in E[m]$.
\item Adjoint: for an isogeny $\varphi\colon E\to E'$, $S\in E[m]$, $T\in E'[m]$: $e_m(S,\hat\varphi T)=e_m(\varphi S,T)$.
\end{enumerate}
\end{theorem}

\begin{proof}[Sketch of proof]
For $T\in E[m]$ there is $f\in\bar K(E)$ with $\operatorname{div}(f)=m(T)-m(O)$ (\Cref{cor:ellcurves:principal}). Choose $T'$ with $mT'=T$; there is $g$ with
$\operatorname{div}(g)=[m]^*(T)-[m]^*(O)=\sum_{R\in E[m]}\bigl((T'+R)-(R)\bigr)$, again by \Cref{cor:ellcurves:principal}, since $\sum_R(T'+R-R)=[m^2]T'=O$. Then $f\circ[m]$ and $g^m$ have
the same divisor, so after scaling $f\circ[m]=g^m$. For $S\in E[m]$ and any point $X$, $g(X+S)^m=f(mX+mS)=f(mX)=g(X)^m$, so $g(X+S)/g(X)$ is an $m$-th root of unity; as a
function of $X$ it is a morphism $E\to\PP^1$ that is not surjective, hence constant (\Cref{ch:curves}). Define $e_m(S,T)=g(X+S)/g(X)$. Properties (i) to (vi) are
proved in \cite[Section III.8]{SilvermanAEC}; (iv) is immediate because $\sigma$ carries the functions $f,g$ for $T$ to functions with the same properties for $\sigma T$.
Over $\C$, \Cref{prop:ellcurves:cplxisog}(iv) gives an explicit formula.
\end{proof}

\begin{corollary}[the determinant]\label{cor:ellcurves:det}
Let $m\ge1$ with $\operatorname{char}K\nmid m$.
\begin{enumerate}[label=(\roman*)]
\item For $\sigma\in G_K$ acting on $E[m]$ by a matrix $g_\sigma\in\GL_2(\Z/m)$ in some basis, $\det g_\sigma=\chi_m^{\mathrm{cyc}}(\sigma)$.
\item If $E[m]\subseteq E(K)$ then $\mu_m\subseteq K$.
\item Let $\ell$ be prime and $L\subset E[\ell]$ a $G_K$-stable line. Then $E[\ell]/L\simeq\Hom(L,\mu_\ell)$ as $G_K$-modules. If $G_K$ acts on $L$ by $\eta_1$ and on $E[\ell]/L$ by
$\eta_2$, then $\eta_1\eta_2=\bar\chi_\ell$. In particular a $K$-rational point of order $\ell$ gives an exact sequence of $G_K$-modules $0\to\Z/\ell\to E[\ell]\to\mu_\ell\to0$.
\end{enumerate}
\end{corollary}

\begin{proof}
(i) Let $S,T$ be the basis and $\zeta=e_m(S,T)$. It is a primitive $m$-th root of unity: if $\zeta^k=1$ then $e_m(kS,T)=1=e_m(kS,S)$, so $e_m(kS,\cdot)=1$ on $E[m]$ and
$kS=O$ by (iii) and (ii) of \Cref{thm:ellcurves:weil}, so $m\mid k$. If $\sigma S=aS+cT$ and $\sigma T=bS+dT$, bilinearity and the alternating property give
$e_m(\sigma S,\sigma T)=e_m(S,T)^{ad}e_m(T,S)^{cb}=\zeta^{ad-bc}$, while Galois-equivariance gives $e_m(\sigma S,\sigma T)=\sigma(\zeta)=\zeta^{\chi_m^{\mathrm{cyc}}(\sigma)}$. So
$\det g_\sigma\equiv\chi_m^{\mathrm{cyc}}(\sigma)\pmod m$. (ii) If $G_K$ fixes $E[m]$ pointwise then $\chi_m^{\mathrm{cyc}}(\sigma)=\det1=1$ for all $\sigma$, so $G_K$ fixes $\mu_m$.
(iii) Consider $\theta\colon E[\ell]\to\Hom(L,\mu_\ell)$, $P\mapsto(Q\mapsto e_\ell(Q,P))$. Its kernel is $L^\perp=\{P:e_\ell(L,P)=1\}$, which contains $L$ (alternating) and has
dimension $1$ (nondegenerate), so $\ker\theta=L$ and $\theta$ is onto by counting. With the action $(\sigma h)(Q)=\sigma(h(\sigma^{-1}Q))$ on $\Hom(L,\mu_\ell)$,
$(\sigma\theta(P))(Q)=\sigma e_\ell(\sigma^{-1}Q,P)=e_\ell(Q,\sigma P)$, so $\theta$ is $G_K$-equivariant. $G_K$ acts on $\Hom(L,\mu_\ell)$ by $\bar\chi_\ell\eta_1^{-1}$, so $\eta_2=\bar\chi_\ell\eta_1^{-1}$. If
$L$ is spanned by a $K$-rational point, $\eta_1=1$ and $\Hom(L,\mu_\ell)\simeq\mu_\ell$.
\end{proof}

\begin{definition}[Tate modules]\label{def:ellcurves:tate}
For a prime $\ell\ne\operatorname{char}K$, the \emph{$\ell$-adic Tate module} of $E$ is $T_\ell E=\varprojlim_nE[\ell^n]$, the inverse limit along the maps
$[\ell]\colon E[\ell^{n+1}]\to E[\ell^n]$, and $V_\ell E=T_\ell E\otimes_{\Z_\ell}\Q_\ell$. Also $E[\ell^\infty]=\bigcup_nE[\ell^n]$. These are the $T_\ell J$ and $V_\ell J$ of \Cref{ch:curves} for
$J=E$, an elliptic curve being its own Jacobian.
\end{definition}

\begin{proposition}[structure of the Tate module]\label{prop:ellcurves:tate}
\leavevmode
\begin{enumerate}[label=(\roman*)]
\item $T_\ell E$ is a free $\Z_\ell$-module of rank $2$, and the projection $T_\ell E\to E[\ell^n]$ is surjective with kernel $\ell^nT_\ell E$; so $T_\ell E/\ell^nT_\ell E\simeq E[\ell^n]$.
\item The map $V_\ell E\to E[\ell^\infty]$, $\ell^{-n}t\mapsto t_n$ (the $n$-th component of $t\in T_\ell E$), is well defined, surjective, and has kernel $T_\ell E$; so
$V_\ell E/T_\ell E\simeq E[\ell^\infty]\simeq(\Q_\ell/\Z_\ell)^2$.
\item The pairings $e_{\ell^n}$ assemble into a pairing $e\colon T_\ell E\times T_\ell E\to\Z_\ell(1)=\varprojlim\mu_{\ell^n}$, $e\bigl((S_n),(T_n)\bigr)=\bigl(e_{\ell^n}(S_n,T_n)\bigr)_n$, which
is bilinear, alternating, perfect and $G_K$-equivariant.
\item Let $\operatorname{char}K=p>0$ and define $T_pE=\varprojlim_nE[p^n]$ by the same formula ($E[p^n]$ being the group of $\bar K$-points killed by $p^n$). Then
$T_pE$ is $0$ or $\Z_p$; this rests on \Cref{cor:ellcurves:degm}(iii), which we do not prove.
\end{enumerate}
\end{proposition}

\begin{proof}
(i) Choose a basis $(P_1,Q_1)$ of $E[\ell]\simeq(\Z/\ell)^2$ and, inductively, $P_{n+1},Q_{n+1}$ with $\ell P_{n+1}=P_n$, $\ell Q_{n+1}=Q_n$ ($[\ell]$ is surjective). The subgroup
$H=\langle P_{n+1},Q_{n+1}\rangle$ of $E[\ell^{n+1}]\simeq(\Z/\ell^{n+1})^2$ satisfies $\ell H=\langle P_n,Q_n\rangle=E[\ell^n]$ (induction) and $H\supseteq\ell^nH=E[\ell]$, so
$\#H=\#\ell H\cdot\#H[\ell]=\ell^{2n}\cdot\ell^2$ and $H=E[\ell^{n+1}]$. In these compatible bases each $E[\ell^n]$ is $(\Z/\ell^n)^2$ and the transition maps are the reductions, so
$T_\ell E\simeq\varprojlim(\Z/\ell^n)^2=\Z_\ell^2$ (\Cref{ch:groups}), and the projection to the $n$-th term is reduction modulo $\ell^n$. (ii) If $\ell^{-n}t=\ell^{-n'}t'$ with $n'\ge n$ then
$t'=\ell^{n'-n}t$, whose $n'$-th component is $\ell^{n'-n}t_{n'}=t_n$; so the map is well defined. It is surjective by (i), and $\ell^{-n}t\mapsto0$ iff $t_n=0$ iff
$t\in\ell^nT_\ell E$ iff $\ell^{-n}t\in T_\ell E$. (iii) Compatibility of the components: by \Cref{thm:ellcurves:weil}(v),
$e_{\ell^{n+1}}(S_{n+1},T_{n+1})^\ell=e_{\ell^{n+1}}(S_{n+1},\ell T_{n+1})=e_{\ell^{n+1}}(S_{n+1},T_n)=e_{\ell^n}(\ell S_{n+1},T_n)=e_{\ell^n}(S_n,T_n)$, which is the transition map
$\zeta\mapsto\zeta^\ell$ of $\Z_\ell(1)$. Bilinearity, the alternating property and equivariance pass to the limit. Modulo $\ell$ the pairing is $e_\ell$, which is
nondegenerate, so the Gram matrix in a $\Z_\ell$-basis has unit determinant and $e$ is perfect. (iv) Given \Cref{cor:ellcurves:degm}(iii), the same argument as
(i) with one generator: either all $E[p^n]$ vanish, or $E[p^n]\simeq\Z/p^n$ with surjective transition maps $[p]$, and the limit is $\Z_p$.
\end{proof}

\begin{definition}[the $\ell$-adic and mod $\ell$ representations]\label{def:ellcurves:rho}
$G_K$ acts on each $E[\ell^n]$ by group automorphisms (\Cref{thm:ellcurves:grouplaw}(ii)), compatibly with $[\ell]$, hence on $T_\ell E$ and $V_\ell E$. After a choice of
$\Z_\ell$-basis of $T_\ell E$ this is a homomorphism
\[
\rho_{E,\ell}\colon G_K\to\Aut_{\Z_\ell}(T_\ell E)\simeq\GL_2(\Z_\ell),
\]
and its reduction $\bar\rho_{E,\ell}\colon G_K\to\Aut(E[\ell])\simeq\GL_2(\F_\ell)$ is the action on $E[\ell]=T_\ell E/\ell T_\ell E$; more generally $\rho_{E,\ell}\bmod\ell^k$ is the action on
$E[\ell^k]$. Different bases give conjugate homomorphisms.
\end{definition}

\begin{proposition}[properties of $\rho_{E,\ell}$]\label{prop:ellcurves:rho}
\leavevmode
\begin{enumerate}[label=(\roman*)]
\item $\rho_{E,\ell}$ is continuous. The field $K(E[\ell^n])$ generated by the coordinates of the points of $E[\ell^n]$ is a finite Galois extension of $K$, and $\rho_{E,\ell}\bmod\ell^n$
factors through $\Gal(K(E[\ell^n])/K)$, on which it is injective.
\item $\det\rho_{E,\ell}=\chi_\ell$, the $\ell$-adic cyclotomic character, and $\det\bar\rho_{E,\ell}=\bar\chi_\ell$.
\item An isogeny $\varphi\colon E\to E'$ defined over $K$ induces an injective $G_K$-equivariant map $\varphi_\ell\colon T_\ell E\to T_\ell E'$ with
$\coker\varphi_\ell\simeq(\ker\varphi)[\ell^\infty]$. Hence $V_\ell E\simeq V_\ell E'$, and $T_\ell E\simeq T_\ell E'$ as $G_K$-modules when $\ell\nmid\deg\varphi$.
\end{enumerate}
\end{proposition}

\begin{proof}
(i) $E[\ell^n]$ is finite and $G_K$-stable (\Cref{thm:ellcurves:grouplaw}(ii)), so its coordinates generate a finite extension $L_n$ stable under $G_K$, hence Galois, provided it is
separable over $K$. Separability: let $m=\ell^n$ and $P\in E[m]\smallsetminus\{O\}$. If $m$ is odd, $x(P)$ is a root of $\psi_m$, a polynomial over $K$ whose roots are all
simple (\Cref{cor:ellcurves:divpol-roots}). If $m$ is even ($\ell=2\ne\operatorname{char}K$), either $2P=O$ and $x(P)$ is a root of the cubic $F$, whose roots are simple
(\Cref{prop:ellcurves:addition}(iv)), or $\psi_2(P)\ne0$, and then $x(P)$ is a root of $f_m=\psi_m/\psi_2$ (\Cref{thm:ellcurves:multm}); the
$(m^2-4)/2$ distinct values $x(P)$ for $P\in E[m]\smallsetminus E[2]$ (points paired with their negatives) are roots of $f_m$, which has that degree and leading
coefficient $m/2\ne0$ in $K$ (\Cref{prop:ellcurves:divpol-deg}), so they are simple roots. A simple root of a polynomial over $K$ is separable over $K$. Then $y(P)$ is a root of
$g(Y)=Y^2+(a_1x(P)+a_3)Y-(x(P)^3+a_2x(P)^2+a_4x(P)+a_6)$, whose derivative at $y(P)$ is $\psi_2(P)$; if $\psi_2(P)\ne0$ it is a simple root, hence separable over
$K(x(P))$, and if $\psi_2(P)=0$ (so $\operatorname{char}K\ne2$, since in characteristic $2$ the prime $\ell$ is odd and $2P\ne O$) then $y(P)=-(a_1x(P)+a_3)/2\in K(x(P))$. So the coordinates of the points of $E[\ell^n]$ lie in
$K^{\mathrm{sep}}$; this argument does not assume $K$ perfect, as \cite[Section III.7]{SilvermanAEC} does. An element of $G_K$ acts trivially on $E[\ell^n]$ iff it fixes $L_n$. So the kernel of $\rho_{E,\ell}\bmod\ell^n$ is the open subgroup $G_{L_n}$, and the
inverse images of the basic open subgroups $\ker(\GL_2(\Z_\ell)\to\GL_2(\Z/\ell^n))$ are open: $\rho_{E,\ell}$ is continuous (\Cref{ch:galois}). (ii) By \Cref{cor:ellcurves:det}(i)
with $m=\ell^n$ for all $n$, $\det\rho_{E,\ell}(\sigma)\equiv\chi_\ell(\sigma)\pmod{\ell^n}$ for every $n$. (iii) $\varphi$ maps $E[\ell^n]$ to $E'[\ell^n]$ compatibly with $[\ell]$ and
with $G_K$ (\Cref{thm:ellcurves:isog-hom}(i), \Cref{def:ellcurves:isogeny}), so it induces $\varphi_\ell$, and $\varphi_V=\varphi_\ell\otimes\Q_\ell$. As $\hat\varphi\circ\varphi=[\deg\varphi]$ is
invertible on $V_\ell E$, $\varphi_V$ is injective, hence an isomorphism by dimensions, and $\varphi_\ell$ is injective. Under the identifications $V/T\simeq E[\ell^\infty]$ of
\Cref{prop:ellcurves:tate}(ii), $\varphi_V$ induces $\varphi\colon E[\ell^\infty]\to E'[\ell^\infty]$, whose kernel is $\varphi_V^{-1}(T_\ell E')/T_\ell E\simeq T_\ell E'/\varphi_\ell(T_\ell E)=\coker\varphi_\ell$
and also equals $(\ker\varphi)[\ell^\infty]$. Finally $\#\ker\varphi$ divides $\deg\varphi$ (\Cref{thm:ellcurves:isog-hom}(ii)), so if $\ell\nmid\deg\varphi$ the cokernel is $0$.
\end{proof}

\begin{corollary}[$T_\ell E$ is self-dual up to a twist]\label{cor:ellcurves:dual-T}
The Weil pairing gives $G_K$-isomorphisms
\[
T_\ell E\xrightarrow{\ \sim\ }\Hom_{\Z_\ell}(T_\ell E,\Z_\ell(1))=(T_\ell E)^\vee(1),\quad t\mapsto e(\cdot,t),\qquad E[\ell^k]\simeq\Hom(E[\ell^k],\mu_{\ell^k}).
\]
Equivalently $(T_\ell E)^\vee\simeq T_\ell E(-1)$ and $\rho_{E,\ell}^\vee\simeq\rho_{E,\ell}\otimes\chi_\ell^{-1}$, which is the identity $V^\vee\simeq V\otimes\det^{-1}$ for two-dimensional $V$
(\Cref{ch:reps}). Consequently $\ad\bigl((T_\ell E)^\vee\bigr)\simeq\ad(T_\ell E)$ as $G_K$-modules (\Cref{ch:reps}): the trace-zero adjoint does not distinguish $T_\ell E$ from its dual,
nor from any lattice homothetic to either.
\end{corollary}

\begin{proof}
The map is $\Z_\ell$-linear, bijective because $e$ is perfect (\Cref{prop:ellcurves:tate}(iii)), and equivariant by the computation in the proof of
\Cref{cor:ellcurves:det}(iii) with $e$ in place of $e_\ell$. The finite-level statement is the same argument with $e_{\ell^k}$. Tensoring with $\Z_\ell(-1)$ gives the second form.
Transposition $g\mapsto g^{\mathsf T}$ is a $\Z_\ell$-linear bijection $\End(T)\to\End(T^\vee)$ preserving traces, and it is $G_K$-equivariant because
$\rho^\vee(\sigma)g^{\mathsf T}\rho^\vee(\sigma)^{-1}=(\rho(\sigma)g\rho(\sigma)^{-1})^{\mathsf T}$ for $\rho^\vee(\sigma)=(\rho(\sigma)^{-1})^{\mathsf T}$; so $\ad(T)\simeq\ad(T^\vee)$. Twisting by a
character does not change $\End$, hence not $\ad$ (\Cref{ch:reps}).
\end{proof}

\begin{corollary}[complex conjugation]\label{cor:ellcurves:complex-conj}
Let $E/\Q$, $\ell$ odd, and $c\in G_\Q$ a complex conjugation. Then $\bar\rho_{E,\ell}(c)$ is conjugate in $\GL_2(\F_\ell)$ to $\operatorname{diag}(1,-1)$: it has eigenvalues
$1,-1$, trace $0$ and determinant $-1$. The same holds for $\rho_{E,\ell}(c)$ in $\GL_2(\Z_\ell)$.
\end{corollary}

\begin{proof}
$c^2=1$ and $\chi_\ell(c)=-1$, as $c$ inverts roots of unity. A matrix $g$ with $g^2=1$ over a ring in which $2$ is invertible satisfies $M=M^+\oplus M^-$ with
$M^\pm=\frac{1\pm g}2M$, the $\pm1$-eigenspaces; over $\Z_\ell$ (or $\F_\ell$) both are free, of ranks adding to $2$. Since $\det g=-1$, $g\ne\pm1$, so both ranks are $1$. (Directly: $E(\R)$ is
a compact real Lie group of dimension one, $S^1$ or $S^1\times\Z/2$, so its $\ell$-torsion $E[\ell]^{c}$ is cyclic of order $\ell$.)
\end{proof}

\begin{convention}[Frobenius]\label{conv:ellcurves:frobenius}
\leavevmode
\begin{enumerate}[label=(\alph*)]
\item For $E$ over a finite field $\F_{\mathfrak q}$, the \emph{Frobenius endomorphism} $\pi=\pi_{\mathfrak q}\colon(x,y)\mapsto(x^{\mathfrak q},y^{\mathfrak q})$ acts on $E(\Fbar_p)$ exactly as the
arithmetic Frobenius $\Frob_{\mathfrak q}\in G_{\F_{\mathfrak q}}$, $t\mapsto t^{\mathfrak q}$.
\item For $E/\Q$ with good reduction at $p\ne\ell$, $\Frob_p$ in a group-theoretic statement is an \emph{arithmetic} Frobenius at $p$ (\Cref{ch:ant}), and $\Frob_p^{-1}$ is a
geometric one. $\rho_{E,\ell}(\Frob_p)$ has characteristic polynomial $X^2-a_pX+p$ (\Cref{thm:ellcurves:frobQ}); its roots $\alpha_p,\beta_p$ satisfy $\alpha_p\beta_p=p$.
\item On the dual $V_\ell E^\vee=H^1$ the geometric Frobenius has the eigenvalues $\alpha_p,\beta_p$ (the inverse transpose of $\rho(\Frob_p)$ has eigenvalues
$\alpha_p^{-1},\beta_p^{-1}$, and $\Frob_p^{-1}$ inverts them again), while on $V_\ell E$ itself the geometric Frobenius has the eigenvalues $\alpha_p^{-1}=\beta_p/p$ and
$\beta_p^{-1}=\alpha_p/p$. Local $L$-factors are defined with the geometric Frobenius (Appendix B, \Cref{ch:notation}): $L_p(M,s)=\det(1-\Frob_p^{-1}p^{-s}\mid M^{I_p})^{-1}$.
Hence, at a prime $p$ of good reduction,
\[
L_p(E,s)=\det(1-\Frob_p\,p^{-s}\mid V_\ell E)^{-1}=L_p(V_\ell E^\vee,s)=(1-a_pp^{-s}+p^{1-2s})^{-1},
\]
\[
L_p(V_\ell E,s)=L_p(E,s+1),
\]
the last because $\det(1-\Frob_p^{-1}X\mid V_\ell E)=(1-\beta_pX/p)(1-\alpha_pX/p)$, or because $V_\ell E\simeq V_\ell E^\vee(1)$ (\Cref{cor:ellcurves:dual-T}).
\item On $\ad V_\ell E$ both Frobenii have the eigenvalues $\alpha_p/\beta_p,1,\beta_p/\alpha_p$, so the good factor of the adjoint does not depend on the convention; at a
multiplicative prime it does (on the invariant line the arithmetic Frobenius acts by $q$, the geometric one by $q^{-1}$), and the paper's factor $1-q^{-1-s}$ is
the geometric one (\Cref{ch:local,ch:lfunctions}).
\end{enumerate}
\end{convention}

\begin{paperlink}\label{pl:ellcurves:section11}
Section 1.1 of the paper (``Let $V=V_\ell E$, $T=T_\ell E$, $\rho=\rho_{E,\ell}$, $\bar\rho$ its reduction'') uses \Cref{def:ellcurves:tate,def:ellcurves:rho} and
\Cref{prop:ellcurves:tate,prop:ellcurves:rho}; the paper states the same Frobenius conventions as \Cref{conv:ellcurves:frobenius}. In Section 4.2 and after
Proposition prop:pullback the paper identifies the lattice $M_{f,\lambda}$ of Diamond, Flach and Guo with $H^1(E,\Z_\ell)$. That module is
$\Hom(T_\ell E,\Z_\ell)=(T_\ell E)^\vee$, which by \Cref{cor:ellcurves:dual-T} is $T_\ell E(-1)$: the dual of $T_\ell E$, not $T_\ell E$ itself. Since the adjoint is insensitive to
duals and twists, $\ad M_{f,\lambda}=\ad T_\ell E$, which is all the paper uses. Lemma lem:h0 uses that complex conjugation has distinct eigenvalues $\pm1$
(\Cref{cor:ellcurves:complex-conj}), which makes an irreducible $\bar\rho$ absolutely irreducible (\Cref{ch:galrep}), and $\det\rho=\chi_\ell$ (\Cref{prop:ellcurves:rho}(ii)) in the
determinant argument; Proposition prop:ns uses $\det\bar\rho_{E,3}=\bar\chi_3=\chi_{-3}$.
\end{paperlink}

\section{Endomorphisms and complex multiplication}\label{sec:ellcurves:cm}

Curves with complex multiplication are the exceptions of the paper: Conjecture conj:eis and the data sections exclude them, Proposition prop:no9 reaches a
contradiction by producing complex multiplication, and the theorem of Calegari and Emerton quoted in Section sec:two has a CM alternative. This section
records the structure of endomorphism rings, the Galois action on the endomorphisms of a CM curve over $\Q$, the thirteen rational CM $j$-invariants, and the
facts about $a_p$ that recognise CM. The number-theoretic side (orders, class number one, splitting of primes) is \Cref{ch:ant}.

\begin{theorem}[endomorphism rings]\label{thm:ellcurves:end}
$\End_{\bar K}(E)$ has no zero divisors, and $\deg$ is a positive definite quadratic form on it. As a ring it is isomorphic to $\Z$, to an order in an
imaginary quadratic field, or, only when $\operatorname{char}K=p>0$, to a maximal order in the quaternion algebra over $\Q$ ramified exactly at $p$ and $\infty$. In
characteristic $0$ only the first two occur.
\end{theorem}

\noindent The first sentence follows from \Cref{thm:ellcurves:dual}(iv) and the multiplicativity of degrees. We do not prove the classification; see
\cite[Sections III.9 and V.3]{SilvermanAEC}. Over $\C$ the dichotomy $\Z$ or an imaginary quadratic order is proved in \Cref{prop:ellcurves:cplxisog}(iii), and every
field of characteristic $0$ over which a given $E$ and its endomorphisms are defined embeds into $\C$.

\begin{proposition}[normalised endomorphisms and the Galois action]\label{prop:ellcurves:normalised}
Let $\operatorname{char}K=0$ and $\omega$ the invariant differential of $E$.
\begin{enumerate}[label=(\roman*)]
\item The map $\End_{\bar K}(E)\to\bar K$, $\alpha\mapsto c_\alpha$ where $\alpha^*\omega=c_\alpha\omega$, is an injective ring homomorphism. Its image $\mathcal O$ is $\Z$ or an order in an
imaginary quadratic field $F\subset\bar K$; we write $[c]$ for the endomorphism with $c_{[c]}=c$, consistently with $[m]$ for $m\in\Z$.
\item An endomorphism defined over $K$ has $c_\alpha\in K$. In particular $\End_\Q(E)=\Z$ for every $E/\Q$.
\item Let $E/\Q$ and $\End_{\Qbar}(E)=\mathcal O\ne\Z$, with fraction field $F$. For $\sigma\in G_\Q$ and $c\in\mathcal O$, $\sigma\circ[c]\circ\sigma^{-1}=[\sigma(c)]$ on $E(\Qbar)$. So every
endomorphism is defined over $F$, and $\sigma$ acts on $\mathcal O$ by complex conjugation when $\sigma|_F\ne1$. On the Tate module,
$\rho_{E,\ell}(\sigma)\,[c]_\ell\,\rho_{E,\ell}(\sigma)^{-1}=[\sigma(c)]_\ell$.
\end{enumerate}
\end{proposition}

\begin{proof}
(i) Additivity is \Cref{prop:ellcurves:invdiff}(iii), and $(\alpha\circ\beta)^*\omega=\beta^*(c_\alpha\omega)=c_\alpha c_\beta\omega$. Every holomorphic differential is a constant multiple of
$\omega$ (\Cref{ch:curves}), so $c_\alpha$ is defined. A nonzero $\alpha$ is separable (characteristic $0$), so $\alpha^*\omega\ne0$ (\Cref{ch:curves}): the map is injective. Its
image is a commutative ring isomorphic to $\End_{\bar K}(E)$, so by \Cref{thm:ellcurves:end} it is $\Z$ or an imaginary quadratic order.
(ii) If $\alpha$ is defined over $K$, then $\alpha^*\omega$ is a differential over $K$ and $\omega$ is one, so $c_\alpha\in K$. For $K=\Q$, $c_\alpha$ is a rational number integral over $\Z$
(it lies in an order), hence an integer $m$, and $\alpha=[m]$ by injectivity.
(iii) For $\sigma\in G_\Q$ let $\alpha^\sigma$ be the endomorphism obtained by applying $\sigma$ to the coefficients of the rational functions defining $\alpha$; on points,
$\alpha^\sigma(P)=\sigma(\alpha(\sigma^{-1}P))$. Applying $\sigma$ to $\alpha^*\omega=c_\alpha\omega$, and using that $\omega$ has rational coefficients, gives
$(\alpha^\sigma)^*\omega=\sigma(c_\alpha)\omega$; so $[c]^\sigma=[\sigma(c)]$ by (i). If $\sigma|_F=1$ this gives $[c]^\sigma=[c]$ for all $c$, so every endomorphism is defined over
$F$ (\Cref{ch:curves}). The statement on $T_\ell E$ is the same identity applied to $\ell$-power torsion points.
\end{proof}

\begin{definition}[complex multiplication]\label{def:ellcurves:cm}
An elliptic curve $E$ over a field of characteristic $0$ has \emph{complex multiplication} (\emph{CM}) if $\End_{\bar K}(E)\ne\Z$. Then
$\End_{\bar K}(E)\simeq\mathcal O$, an order in an imaginary quadratic field $F$ (\Cref{prop:ellcurves:normalised}), and $E$ has \emph{CM by $\mathcal O$}, or \emph{by $F$}.
\end{definition}

\begin{lemma}[CM is an isogeny invariant]\label{lem:ellcurves:cm-isog}
If $\varphi\colon E\to E'$ is a nonconstant isogeny of degree $n$, then $\alpha\mapsto\frac1n\varphi\circ\alpha\circ\hat\varphi$ is an isomorphism of $\Q$-algebras
$\End(E)\otimes\Q\to\End(E')\otimes\Q$. In particular $E$ has CM iff $E'$ has, and then by the same field $F$.
\end{lemma}

\begin{proof}
$\frac1{n^2}(\varphi\alpha\hat\varphi)(\varphi\beta\hat\varphi)=\frac1{n^2}\varphi\alpha[n]\beta\hat\varphi=\frac1n\varphi\alpha\beta\hat\varphi$ (\Cref{thm:ellcurves:dual}), additivity is clear, and $1\mapsto\frac1n\varphi\hat\varphi=1$. The map
$\beta\mapsto\frac1n\hat\varphi\beta\varphi$ in the other direction is inverse to it, by the same computation.
\end{proof}

\begin{theorem}[the rational CM $j$-invariants]\label{thm:ellcurves:13j}
If $E/\Q$ has CM by the order $\mathcal O$, then $\mathcal O$ has class number one, so its discriminant $D_{\mathcal O}$ is one of the thirteen values in the table below, and $j(E)$ is the
corresponding value; conversely each value is the $j$-invariant of a curve over $\Q$ with CM by that order. Here $j(E)=j(\tau)$ for $\tau=\frac{D_{\mathcal O}+\sqrt{D_{\mathcal O}}}2$ (\Cref{sec:ellcurves:complex}).
\begin{center}\footnotesize
\begin{tabular}{rrl|rrl}
\toprule
$D_{\mathcal O}$ & $j$ & example & $D_{\mathcal O}$ & $j$ & example\\
\midrule
$-3$ & $0$ & \curve{27a1}, \curve{36a1} & $-19$ & $-884736=-2^{15}3^3$ & \\
$-4$ & $1728$ & \curve{32a1} & $-27$ & $-12288000=-2^{15}\cdot3\cdot5^3$ & \curve{27a2}, \curve{27a4}\\
$-7$ & $-3375=-3^35^3$ & & $-28$ & $16581375=3^35^317^3$ & \\
$-8$ & $8000=2^65^3$ & & $-43$ & $-884736000=-2^{18}3^35^3$ & \\
$-11$ & $-32768=-2^{15}$ & \curve{121b1} & $-67$ & $-147197952000=-2^{15}3^35^311^3$ & \\
$-12$ & $54000=2^43^35^3$ & & $-163$ & $-262537412640768000=-2^{18}3^35^323^329^3$ & \\
$-16$ & $287496=2^33^311^3$ & & & & \\
\bottomrule
\end{tabular}
\end{center}
\end{theorem}

\noindent We do not prove this; see \cite[Appendix C, Section 11]{SilvermanAEC}, \cite[Chapter II]{SilvermanATAEC} and \cite[Sections 7 and 10 to 12]{Cox}. The
ingredients are: $j(E)$ generates the ring class field of $\mathcal O$ over $F$, of degree $h(\mathcal O)$ (so $j(E)\in\Q$ forces $h(\mathcal O)=1$), and the
theorem of Heegner, Baker and Stark (\Cref{ch:ant}) lists the orders of class number one. The thirteen values were checked numerically as $j(\tau)$ with PARI's
\texttt{ellj}, and the examples' $j$-invariants with \texttt{ellinit}.

\begin{theorem}[CM and the numbers $a_p$]\label{thm:ellcurves:cm-ap}
Let $E/\Q$ have CM by an order in the imaginary quadratic field $F$, and let $\chi_F$ be the quadratic character of $F$.
\begin{enumerate}[label=(\roman*)]
\item $\rho_{E,\ell}\simeq\rho_{E,\ell}\otimes\chi_F$ for every prime $\ell$.
\item $a_p(E)=0$ for every prime $p$ of good reduction that is inert in $F$ (the reduction is supersingular, \Cref{def:ellcurves:supersingular}).
\item For a prime $p$ of good reduction that splits in $F$, there is $\pi_p$ in the ring of integers of $F$ with $\pi_p\bar\pi_p=p$ and $a_p(E)=\pi_p+\bar\pi_p\ne0$.
\item Every prime ramified in $F$ divides the conductor $N$ of $E$. In particular, if $E$ has CM by $\Q(\sqrt{\ell^*})$ for an odd prime $\ell$, then $\ell\mid N$.
\item Conversely, if $a_p(E)=0$ for all primes $p$ inert in an imaginary quadratic field $F$ outside a set of density $0$, then $E$ has CM by $F$.
\end{enumerate}
\end{theorem}

\begin{proof}[Sketch of proof]
(i), (ii) and (iv) are proved in \Cref{ch:ant} from \Cref{prop:ellcurves:normalised}(iii) and \Cref{thm:ellcurves:frobQ}; we recall the argument. Choose $0\ne c\in\mathcal O$ with
$\bar c=-c$ (a suitable multiple of $\sqrt{D_{\mathcal O}}$). By \Cref{prop:ellcurves:normalised}(iii), $\rho(\sigma)[c]_\ell\rho(\sigma)^{-1}=\chi_F(\sigma)[c]_\ell$, so the invertible map $[c]_\ell$ on
$V_\ell E$ is an isomorphism $\rho\to\rho\otimes\chi_F$. Taking traces, $\tr\rho(\sigma)=\chi_F(\sigma)\tr\rho(\sigma)$, and at a good $p\ne\ell$ inert in $F$, $\chi_F(\Frob_p)=-1$ gives
$a_p=\tr\rho(\Frob_p)=0$ (\Cref{thm:ellcurves:frobQ}). For (iv): if $p$ ramified in $F$ were a prime of good reduction, take $\ell\ne p$; $\rho$ is unramified at $p$
(\Cref{thm:ellcurves:frobQ}), but some $\tau\in I_p$ has $\chi_F(\tau)=-1$, and then $[c]_\ell=\rho(\tau)[c]_\ell\rho(\tau)^{-1}=-[c]_\ell$, a contradiction. (iii) is Deuring's theorem
\cite{Deuring41}: the reduction is ordinary, its endomorphism ring contains the reduction of $\mathcal O$ and is an order of $F$, and the Frobenius endomorphism $\pi_p$
(\Cref{def:ellcurves:frobendo}) is an element of norm $p$ with trace $a_p$; see \cite[Chapter II, Section 10]{SilvermanATAEC}. (v) follows from Serre's open image theorem
(\Cref{ch:galrep}); see \cite{SerreAbelian}.
\end{proof}

The checks: $a_p=0$ at every good inert $p<500$ and $a_p\ne0$ at every good split $p<500$ for \curve{27a1}, \curve{36a1} ($F=\Q(\sqrt{-3})$), \curve{32a1} ($F=\Q(i)$)
and \curve{121b1} ($F=\Q(\sqrt{-11})$); and among the $462$ curves of conductor below $10{,}000$ with one of the thirteen $j$-invariants, every prime dividing the
discriminant of $F$ divides the conductor, as (iv) says.

\begin{warning}[the cyclotomic character is not a quadratic character]\label{warn:ellcurves:chi-l}
The CM condition of \Cref{thm:ellcurves:cm-ap}(i) involves the quadratic character $\chi_F$; for $F=\Q(\sqrt{\ell^*})$ this is $\chi_{\ell^*}=\bar\chi_\ell^{(\ell-1)/2}$. The cyclotomic
character $\chi_\ell$ is not quadratic, and $\rho_{E,\ell}\simeq\rho_{E,\ell}\otimes\chi_\ell$ is impossible for every $E$ (determinants give $\chi_\ell^2=1$). Only for $\ell=3$ do $\bar\chi_\ell$
and $\chi_{\ell^*}$ coincide: $\bar\chi_3=\chi_{-3}$. The original version of the paper's Lemma lem:h0 confused the two; the revised version works with $\bar\chi_\ell$ and needs no
CM argument. By \Cref{thm:ellcurves:cm-ap}(iv) the CM alternative ``$E$ has CM by $\Q(\sqrt{\ell^*})$'' cannot occur when $\ell\nmid N$.
\end{warning}

\begin{paperlink}\label{pl:ellcurves:cm}
The paper excludes CM curves in Conjecture conj:eis and in its data (``not CM''). Proposition prop:no9 (\Cref{prop:ellcurves:no27} below) ends with ``so $E'$, hence
$E$, has complex multiplication'': the passage from $E'$ to $E$ is \Cref{lem:ellcurves:cm-isog}. The CM alternative in the theorem of Calegari and Emerton (Section sec:two:
``it has complex multiplication and $N\in\{27,32,49,243\}$'') concerns CM curves of these four conductors; the optimal curves \curve{27a1}, \curve{32a1},
\curve{49a1} and \curve{243b1}, of modular degrees $1$, $1$, $1$ and $9$ (\texttt{degphi}), have $j=0$, $1728$, $-3375$ and $0$, that is, CM by the orders of
discriminant $-3$, $-4$, $-7$, $-3$ (\Cref{thm:ellcurves:13j}). The paper's examples are \curve{27a1} and \curve{243b1}.
\end{paperlink}
\section{Elliptic curves over the complex numbers}\label{sec:ellcurves:complex}

Section 1 of the paper describes the modular degree as ``$\langle f,f\rangle$ divided by the area of the period lattice of $E$'' (Zagier, Watkins; the exact
formula, with the Manin constant, is in \Cref{ch:moddeg}). This section constructs the period lattice: it proves that every elliptic curve over $\C$ is a torus
$\C/\Lambda$, through the Weierstrass function $\wp$, that isomorphism classes of curves over $\C$ correspond to lattices up to homothety and to points of
$\SL_2(\Z)\backslash\HH$, and computes the periods of \curve{11a1}. The analysis is that of \Cref{ch:prereq}: a doubly periodic function is, in the language of
physics, a function on a torus, and the Weierstrass function plays for the torus the role that a Fourier series plays for the circle.

\begin{definition}[lattices]\label{def:ellcurves:lattice}
A \emph{lattice} is a subgroup $\Lambda=\Z\omega_1+\Z\omega_2\subset\C$ with $\omega_1,\omega_2$ linearly independent over $\R$; the basis is \emph{oriented} if
$\operatorname{Im}(\omega_2/\omega_1)>0$. Lattices $\Lambda$, $\Lambda'$ are \emph{homothetic} if $\Lambda'=c\Lambda$ for some $c\in\C^\times$. The \emph{covolume} $\vol(\Lambda)=\abs{\operatorname{Im}(\bar\omega_1\omega_2)}$ is
the area of a fundamental parallelogram. For $k\ge2$ the \emph{Eisenstein series} of $\Lambda$ is $G_{2k}(\Lambda)=\sum_{\omega\in\Lambda\smallsetminus\{0\}}\omega^{-2k}$, and
\[
g_2(\Lambda)=60G_4(\Lambda),\qquad g_3(\Lambda)=140G_6(\Lambda),\qquad\Delta(\Lambda)=g_2^3-27g_3^2,\qquad j(\Lambda)=1728\frac{g_2^3}{\Delta(\Lambda)} .
\]
\end{definition}

\begin{lemma}[lattice sums]\label{lem:ellcurves:latticesum}
$\sum_{\omega\in\Lambda\smallsetminus\{0\}}\abs\omega^{-s}$ converges for every real $s>2$. Hence $G_{2k}(\Lambda)$ converges absolutely for $k\ge2$, and
$G_{2k}(c\Lambda)=c^{-2k}G_{2k}(\Lambda)$, $g_2(c\Lambda)=c^{-4}g_2(\Lambda)$, $g_3(c\Lambda)=c^{-6}g_3(\Lambda)$, $\Delta(c\Lambda)=c^{-12}\Delta(\Lambda)$, and $j(c\Lambda)=j(\Lambda)$ (once
$\Delta(\Lambda)\ne0$ is known, \Cref{thm:ellcurves:wp}).
\end{lemma}

\begin{proof}
The function $(a,b)\mapsto\abs{a\omega_1+b\omega_2}$ is continuous and positive on the compact subset of $\R^2$ where $\max(\abs a,\abs b)=1$, so it has a minimum $\mu>0$ there; by
homogeneity $\abs{a\omega_1+b\omega_2}\ge\mu n$ when $\max(\abs a,\abs b)=n$. There are $8n$ pairs of integers with $\max(\abs a,\abs b)=n$, so the sum is at most
$\sum_{n\ge1}8n(\mu n)^{-s}$, which converges for $s>2$. The homogeneity is clear from the definitions.
\end{proof}

\begin{theorem}[the Weierstrass function]\label{thm:ellcurves:wp}
Let $\Lambda$ be a lattice and
\[
\wp(z)=\wp(z;\Lambda)=\frac1{z^2}+\sum_{\omega\in\Lambda\smallsetminus\{0\}}\Bigl(\frac1{(z-\omega)^2}-\frac1{\omega^2}\Bigr).
\]
\begin{enumerate}[label=(\roman*)]
\item The series converges absolutely and uniformly on compact subsets of $\C\smallsetminus\Lambda$; $\wp$ is meromorphic on $\C$, even, $\Lambda$-periodic, with a double pole at each
point of $\Lambda$ and no other pole.
\item Near $0$, $\wp(z)=z^{-2}+\sum_{k\ge1}(2k+1)G_{2k+2}(\Lambda)z^{2k}$.
\item $\wp'(z)^2=4\wp(z)^3-g_2(\Lambda)\wp(z)-g_3(\Lambda)$.
\item The values $e_1=\wp(\omega_1/2)$, $e_2=\wp(\omega_2/2)$, $e_3=\wp((\omega_1+\omega_2)/2)$ are the three roots of $4x^3-g_2x-g_3$, and they are distinct; hence
$\Delta(\Lambda)=16\prod_{i<j}(e_i-e_j)^2\ne0$.
\item A $\Lambda$-periodic meromorphic function whose poles lie in $\Lambda$ and have order at most $2$ is $a\wp+b$ with $a,b\in\C$; if the poles have order at most $3$, it is
$a\wp'+b\wp+c$.
\end{enumerate}
\end{theorem}

\begin{proof}
(i) For $\abs z\le R$ and $\abs\omega\ge2R$, $\abs{(z-\omega)^{-2}-\omega^{-2}}=\abs{z(2\omega-z)}/\abs{\omega^2(z-\omega)^2}\le R\cdot\frac52\abs\omega\cdot\abs\omega^{-2}\cdot4\abs\omega^{-2}=10R\abs\omega^{-3}$, and
$\sum\abs\omega^{-3}<\infty$ (\Cref{lem:ellcurves:latticesum}). The finitely many terms with $\abs\omega<2R$ are meromorphic. So the series converges normally on compact sets
avoiding $\Lambda$ and $\wp$ is meromorphic (\Cref{ch:prereq}, Weierstrass's theorem on uniform limits), with poles exactly at $\Lambda$, double. Replacing $\omega$ by $-\omega$ shows
$\wp(-z)=\wp(z)$. Differentiating termwise, $\wp'(z)=-2\sum_{\omega\in\Lambda}(z-\omega)^{-3}$, which is visibly $\Lambda$-periodic; so $\wp(z+\omega_i)-\wp(z)$ is a constant $C_i$, and
evaluating at $z=-\omega_i/2$ (not a pole) gives $C_i=\wp(\omega_i/2)-\wp(-\omega_i/2)=0$.
(ii) For $\abs z<\min\abs\omega$, $(z-\omega)^{-2}-\omega^{-2}=\sum_{n\ge1}(n+1)z^n\omega^{-n-2}$; summing over $\omega$ (absolutely convergent) the odd powers cancel in pairs $\pm\omega$, and
the coefficient of $z^{2k}$ is $(2k+1)G_{2k+2}$.
(iii) From (ii), $\wp=z^{-2}+3G_4z^2+5G_6z^4+O(z^6)$, so
$\wp'^2=4z^{-6}-24G_4z^{-2}-80G_6+O(z^2)$ and $4\wp^3=4z^{-6}+36G_4z^{-2}+60G_6+O(z^2)$. Hence $\wp'^2-4\wp^3+60G_4\wp+140G_6=O(z^2)$. The left side is $\Lambda$-periodic and
holomorphic away from $\Lambda$, and by this expansion holomorphic at $0$ and vanishing there; a holomorphic doubly periodic function is bounded, hence constant (Liouville), and the
constant is $0$.
(iv) $\wp'$ is odd and periodic, so $\wp'(\omega_i/2)=-\wp'(-\omega_i/2)=-\wp'(\omega_i/2)$ vanishes at each half-period, and (iii) shows that the $e_i$ are roots of
$4x^3-g_2x-g_3$. On the compact Riemann surface $\C/\Lambda$, $\wp-e_i$ has a single pole, of order $2$, so it has exactly two zeros counted with multiplicity (\Cref{ch:curves}); the
half-period $h_i$ with $\wp(h_i)=e_i$ is a zero of order at least $2$ since $\wp'(h_i)=0$, so it is the only zero. As the three half-periods are distinct modulo $\Lambda$, the $e_i$ are
distinct. The discriminant of $4x^3-g_2x-g_3$ is $16(g_2^3-27g_3^2)$ and also $4^4\prod_{i<j}(e_i-e_j)^2$.
(v) If $f$ has poles of order at most $2$ at the points of $\Lambda$, choose $a$ so that $f-a\wp$ has at most a simple pole at $0$. The residues of a periodic meromorphic function
on $\C/\Lambda$ sum to $0$ (\Cref{ch:curves}, residue theorem), so $f-a\wp$ has no pole; it is holomorphic and doubly periodic, hence constant. The case of order $3$ is the same,
subtracting first a multiple of $\wp'$, which has a triple pole at $0$.
\end{proof}

\begin{theorem}[complex uniformisation]\label{thm:ellcurves:cplxuniform}
\leavevmode
\begin{enumerate}[label=(\roman*)]
\item Let $\Lambda$ be a lattice and $E_\Lambda$ the curve $y^2=4x^3-g_2(\Lambda)x-g_3(\Lambda)$ (a Weierstrass curve after $y\mapsto2y$; it is smooth by \Cref{thm:ellcurves:wp}(iv)). The
map $\Theta\colon\C/\Lambda\to E_\Lambda(\C)$, $z\mapsto(\wp(z),\wp'(z))$ for $z\notin\Lambda$ and $0\mapsto O$, is a biholomorphic group isomorphism, and $\Theta^*(dx/y)=dz$.
\item Let $E$ be an elliptic curve over $\C$ with invariant differential $\omega$, and $\Lambda_E=\{\int_\gamma\omega:\gamma\in H_1(E(\C),\Z)\}$ its lattice of periods. Then $E\simeq E_{\Lambda_E}$, by
an isomorphism that matches $\omega$ with $dx/y$, so $E(\C)\simeq\C/\Lambda_E$ with $\omega\leftrightarrow dz$. Equivalently: for $A,B\in\C$ with $4A^3+27B^2\ne0$ there is a unique lattice $\Lambda$ with
$g_2(\Lambda)=-4A$ and $g_3(\Lambda)=-4B$.
\end{enumerate}
\end{theorem}

\begin{proof}
(i) $\Theta$ is holomorphic (in the chart $(x/y,1/y)$ near $O$ it sends $0$ to $(0,0)$ holomorphically) and lands in $E_\Lambda$ by \Cref{thm:ellcurves:wp}(iii). It is bijective: $O$ has the single
preimage $0$, and for an affine point $(x_0,y_0)$ the function $\wp-x_0$ has exactly two zeros $\pm z_0$ modulo $\Lambda$ (counted with multiplicity, as in the proof of
\Cref{thm:ellcurves:wp}(iv)); since $\wp'(-z_0)=-\wp'(z_0)$ and $\wp'(z_0)^2=y_0^2$, exactly one of $z_0,-z_0$ maps to $(x_0,y_0)$, and if $z_0\equiv-z_0$ then $y_0=\wp'(z_0)=0$. A holomorphic
bijection of compact Riemann surfaces is biholomorphic (\Cref{ch:curves}). Group law: let $z_1,z_2\in\C$ with $\Theta(z_1)$, $\Theta(z_2)$ affine, distinct and not opposite, and let $y=\lambda x+\nu$
be the line through them. Then $f=\wp'-\lambda\wp-\nu$ is periodic with a single pole, of order $3$, at $0$ modulo $\Lambda$, so it has three zeros $z_1,z_2,z_3$ modulo $\Lambda$, and
$z_1+z_2+z_3\in\Lambda$ (the ``only if'' half of Abel's theorem, \Cref{ch:curves}, for the torus: the sum of the zeros minus the sum of the poles of a periodic function lies in $\Lambda$). The
three intersections of the line with $E_\Lambda$ are $\Theta(z_1),\Theta(z_2),\Theta(z_3)$, so $\Theta(z_1)+\Theta(z_2)=-\Theta(z_3)=\Theta(-z_3)=\Theta(z_1+z_2)$ (\Cref{con:ellcurves:chordtangent}; on $E_\Lambda$,
$-(x,y)=(x,-y)$, and $\wp$ is even, $\wp'$ odd). The remaining cases follow by continuity, since $\Theta(z_1+z_2)-\Theta(z_1)-\Theta(z_2)$ is continuous in $(z_1,z_2)$ and vanishes on a dense set.
Finally $\Theta^*(dx/y)=\wp'(z)dz/\wp'(z)=dz$.

(ii) We use the Abel-Jacobi theorem of \Cref{ch:curves}: $E(\C)$ is a compact Riemann surface of genus one, $\omega$ spans its holomorphic differentials, $\Lambda_E$ is a lattice, and
$\iota\colon E(\C)\to\C/\Lambda_E$, $P\mapsto\int_O^P\omega$, is biholomorphic. It is a homomorphism: $\iota(P)$ is the Abel-Jacobi image of $(P)-(O)$, the Abel-Jacobi map is additive on
divisors of degree $0$ and kills principal divisors, and $(P+Q)-(O)\sim(P)-(O)+(Q)-(O)$ (\Cref{prop:ellcurves:pic0}). Since $d\iota=\omega$, $(\iota^{-1})^*\omega=dz$. Put
$\eta=2y+a_1x+a_3$, so that $\eta^2=F(x)=4x^3+b_2x^2+2b_4x+b_6$ (\Cref{lem:ellcurves:identities}), and let $X(z)=x(\iota^{-1}(z))$, $H(z)=\eta(\iota^{-1}(z))$. These are $\Lambda_E$-periodic
meromorphic functions with poles only at $\Lambda_E$, of orders $2$ and $3$ (as $x$ and $\eta$ have poles of orders $2$ and $3$ at $O$ only), and $dX/dz=H$ because $\omega=dx/\eta$
corresponds to $dz$. If $X=cz^{-2}+\cdots$ then $H=-2cz^{-3}+\cdots$, and comparing the $z^{-6}$ terms of $H^2=F(X)$ gives $4c^2=4c^3$, so $c=1$. By \Cref{thm:ellcurves:wp}(v),
$X=\wp+r$ for a constant $r$ (with $\wp=\wp(\cdot;\Lambda_E)$), and then $H=\wp'$. So $\wp'^2=F(\wp+r)$ as functions; since $\wp$ takes infinitely many values, $F(x+r)=4x^3-g_2x-g_3$ as
polynomials. Hence the change of coordinates $x\mapsto x-r$, $\eta\mapsto\eta$ (of the form of \Cref{prop:ellcurves:coordchange}, with $u=1$ after the normalisation $\eta=2y$) maps $E$ to
$E_{\Lambda_E}$, and $\Theta^{-1}$ composed with it is $\iota$; it matches $\omega=dx/\eta$ with $dx/y$ on $E_{\Lambda_E}$. For the second form, apply this to $E\colon y^2=x^3+Ax+B$, for which
$\eta=2y$ and $F(x)=4x^3+4Ax+4B$ has no $x^2$ term, so $r=0$ and $g_2=-4A$, $g_3=-4B$. Uniqueness: if $\Lambda$, $\Lambda'$ have the same $g_2,g_3$, then $E_\Lambda=E_{\Lambda'}$, and $\Theta'^{-1}\circ\Theta$ is
a holomorphic group isomorphism $\C/\Lambda\to\C/\Lambda'$ pulling $dz$ back to $dz$; by \Cref{prop:ellcurves:cplxisog}(i) it is $z\mapsto z$, so $\Lambda=\Lambda'$.
\end{proof}

\noindent A second route to (ii) avoids Abel-Jacobi: the function $j\colon\HH\to\C$ is surjective (\Cref{ch:modforms}, valence formula), so for given $A,B$ there is a lattice with
$j(\Lambda)=j$ of $y^2=x^3+Ax+B$, and a homothety adjusts $g_2,g_3$ \cite[Section VI.5]{SilvermanAEC}. The book uses the route above, so that \Cref{ch:modforms} may cite this
theorem without circularity.

\begin{remark}[translation invariance over $\C$]\label{rem:ellcurves:invdiff-C}
\Cref{thm:ellcurves:cplxuniform}(ii) gives a second proof of \Cref{prop:ellcurves:invdiff}(ii) and (iii) in characteristic $0$. Under $E(\C)\simeq\C/\Lambda_E$ the group law
is addition, $\tau_Q$ is $z\mapsto z+z_Q$, and $\omega$ is $dz$, which translations preserve. For isogenies $\varphi,\psi\colon E'\to E$, \Cref{prop:ellcurves:cplxisog}(i) below
writes them as $z\mapsto\alpha z$ and $z\mapsto\beta z$; then $\varphi+\psi$ is $z\mapsto(\alpha+\beta)z$, which pulls $dz$ back to $(\alpha+\beta)\,dz$, the sum of the
pullbacks. For a curve over any field of characteristic $0$, each identity involves finitely many elements, which generate a field that embeds into $\C$. The proof of \Cref{thm:ellcurves:cplxuniform} used only \Cref{prop:ellcurves:pic0} and \Cref{ch:curves}, not \Cref{prop:ellcurves:invdiff}.
\end{remark}

\begin{proposition}[maps between complex tori]\label{prop:ellcurves:cplxisog}
Let $\Lambda_1,\Lambda_2,\Lambda$ be lattices.
\begin{enumerate}[label=(\roman*)]
\item The holomorphic maps $\C/\Lambda_1\to\C/\Lambda_2$ sending $0$ to $0$ are the maps $z\mapsto\alpha z$ with $\alpha\in\C$, $\alpha\Lambda_1\subseteq\Lambda_2$. They are group homomorphisms, and
isomorphisms iff $\alpha\Lambda_1=\Lambda_2$. Hence $E_{\Lambda_1}\simeq E_{\Lambda_2}$ over $\C$ iff $\Lambda_1$ and $\Lambda_2$ are homothetic, and isogenies $E_{\Lambda_1}\to E_{\Lambda_2}$ correspond to such
$\alpha$.
\item For $\alpha\ne0$ the map has degree $[\Lambda_2:\alpha\Lambda_1]$ and kernel $\alpha^{-1}\Lambda_2/\Lambda_1$. For $\alpha=m$: $\deg[m]=m^2$ and $E[m]=\frac1m\Lambda/\Lambda\simeq(\Z/m)^2$.
\item $\End(\C/\Lambda)=\{\alpha\in\C:\alpha\Lambda\subseteq\Lambda\}$ is $\Z$ or an order in an imaginary quadratic field; the latter happens iff $\tau=\omega_2/\omega_1$ is an imaginary
quadratic number, and then the order lies in $\Q(\tau)$.
\item For a basis $(\omega_1,\omega_2)$ of $\Lambda$, $e_m(a\frac{\omega_1}m+b\frac{\omega_2}m,c\frac{\omega_1}m+d\frac{\omega_2}m)=\zeta^{ad-bc}$, where $\zeta=e_m(\frac{\omega_1}m,\frac{\omega_2}m)$ is a primitive
$m$-th root of unity. For an oriented basis $\zeta=e^{\pm2\pi i/m}$, the sign depending only on the normalisation of $e_m$ (\Cref{conv:ellcurves:conventions}(f)); nothing in the
book depends on it.
\end{enumerate}
\end{proposition}

\begin{proof}
(i) A holomorphic $f\colon\C/\Lambda_1\to\C/\Lambda_2$ with $f(0)=0$ lifts to a holomorphic $F\colon\C\to\C$ with $F(0)=0$ ($\C$ is simply connected and $\C\to\C/\Lambda_2$ is a covering). For
$\omega\in\Lambda_1$, $F(z+\omega)-F(z)\in\Lambda_2$ is continuous in $z$, hence constant; so $F'$ is $\Lambda_1$-periodic and entire, hence constant by Liouville, and $F(z)=\alpha z$ with
$\alpha\Lambda_1\subseteq\Lambda_2$. Conversely such $\alpha$ define holomorphic homomorphisms, with inverse of the same kind iff $\alpha\Lambda_1=\Lambda_2$. By \Cref{thm:ellcurves:cplxuniform} and the fact
that holomorphic maps of compact Riemann surfaces of smooth projective curves are algebraic (\Cref{ch:curves}, Riemann existence and GAGA), isomorphisms and isogenies of the
curves are the same as these maps. (ii) The kernel is $\{z:\alpha z\in\Lambda_2\}/\Lambda_1=\alpha^{-1}\Lambda_2/\Lambda_1$, of order $[\alpha^{-1}\Lambda_2:\Lambda_1]=[\Lambda_2:\alpha\Lambda_1]$; a nonconstant holomorphic
homomorphism of tori is an unramified covering whose fibres are cosets of the kernel, so its degree is the order of the kernel. For $\alpha=m$, $[\Lambda:m\Lambda]=m^2$.
(iii) If $\alpha\Lambda\subseteq\Lambda$, write $\alpha\omega_1=a\omega_1+b\omega_2$ and $\alpha\omega_2=c\omega_1+d\omega_2$ with $a,b,c,d\in\Z$. Then $\alpha$ is an eigenvalue of
$M=\sm abcd$ with eigenvector $(\omega_1,\omega_2)$, so it is a root of the monic integral polynomial $X^2-(a+d)X+(ad-bc)$. If $\alpha$ were real, the real and imaginary parts of the
eigenvector, which span $\R^2$ because $\omega_1,\omega_2$ are $\R$-independent, would both be eigenvectors, so $M=\alpha\cdot1$ and $\alpha\in\Z$. So either $\alpha\in\Z$, or $\alpha$ is a non-real
quadratic integer, and then $b\ne0$ and $\tau=(\alpha-a)/b\in\Q(\alpha)$ is imaginary quadratic. Conversely if $A\tau^2+B\tau+C=0$ with integers $A\ne0$, $B$, $C$, then
$A\tau\Lambda_\tau\subseteq\Lambda_\tau$ for $\Lambda_\tau=\Z+\Z\tau$ ($A\tau\cdot\tau=-B\tau-C$), and $A\tau\notin\Z$. The ring $\End(\C/\Lambda)$ is a subring of $\C$ containing $\Z$, of rank at most $2$ over
$\Z$ (it embeds into $\Lambda\simeq\Z^2$ by $\alpha\mapsto\alpha\omega_1$), so it is $\Z$ or an order in $\Q(\tau)$. (iv) By \Cref{thm:ellcurves:weil}(i), (ii) and the computation in the
proof of \Cref{cor:ellcurves:det}(i). The value of $\zeta$ is computed in \cite[Chapter VI, exercises]{SilvermanAEC}; it matches the intersection pairing on $H_1$ up to a universal sign
(\Cref{ch:curves}).
\end{proof}

\begin{theorem}[moduli of curves over $\C$]\label{thm:ellcurves:moduli}
\leavevmode
\begin{enumerate}[label=(\roman*)]
\item Every lattice is homothetic to $\Lambda_\tau=\Z+\Z\tau$ for some $\tau\in\HH$, and $\Lambda_\tau$, $\Lambda_{\tau'}$ are homothetic iff $\tau'=\frac{a\tau+b}{c\tau+d}$ for some $\sm abcd\in\SL_2(\Z)$. Hence the
isomorphism classes of elliptic curves over $\C$, the homothety classes of lattices, and the points of $\SL_2(\Z)\backslash\HH$ are in bijection.
\item $j(\tau):=j(\Lambda_\tau)$ is holomorphic on $\HH$ and $\SL_2(\Z)$-invariant, and $j(i)=1728$, $j(e^{2\pi i/3})=0$.
\item $j\colon\SL_2(\Z)\backslash\HH\to\C$ is a bijection, and $j(\tau)=q^{-1}+744+196884q+\cdots$ with $q=e^{2\pi i\tau}$.
\end{enumerate}
\end{theorem}

\begin{proof}
(i) Given a lattice, choose an oriented basis (swap $\omega_1$, $\omega_2$ if necessary); then $\omega_1^{-1}\Lambda=\Lambda_\tau$ with $\tau=\omega_2/\omega_1\in\HH$. If $\Lambda_{\tau'}=c\Lambda_\tau$, then $(c,c\tau)$ and
$(1,\tau')$ are two oriented bases of $\Lambda_{\tau'}$ (as $\operatorname{Im}(c\tau/c)=\operatorname{Im}\tau>0$), so $c\tau=a\tau'+b$, $c=c'\tau'+d$ with $\sm ab{c'}d\in\GL_2(\Z)$; then
$\tau=\frac{a\tau'+b}{c'\tau'+d}$, and $\operatorname{Im}\tau=(ad-bc')\operatorname{Im}\tau'/\abs{c'\tau'+d}^2$ forces $ad-bc'=+1$; the inverse matrix expresses $\tau'$ in terms of $\tau$. Conversely $\Lambda_{\gamma\tau}=(c\tau+d)^{-1}\Lambda_\tau$ for $\gamma=\sm abcd$. The bijection with curves is
\Cref{thm:ellcurves:cplxuniform} and \Cref{prop:ellcurves:cplxisog}(i). (ii) $G_{2k}(\Lambda_\tau)=\sum_{(m,n)\ne(0,0)}(m+n\tau)^{-2k}$ converges uniformly on compact subsets of $\HH$ (the bound of
\Cref{lem:ellcurves:latticesum} is locally uniform in $\tau$), so it is holomorphic (\Cref{ch:prereq}); $\Delta(\Lambda_\tau)\ne0$, so $j$ is holomorphic, and invariant by (i) and
\Cref{lem:ellcurves:latticesum}. For $\tau=i$, $i\Lambda_i=\Lambda_i$, so $G_6=i^{-6}G_6=-G_6$, $g_3=0$ and $j=1728$; for $\tau=\zeta_3=e^{2\pi i/3}$, $\zeta_3\Lambda=\Lambda$, so $G_4=\zeta_3^{-4}G_4$,
$g_2=0$ and $j=0$. (iii) We do not prove this here; it is proved in \Cref{ch:modforms} (fundamental domain, valence formula, $q$-expansions). Injectivity also follows from (i) and
\Cref{thm:ellcurves:jinv}(i), surjectivity from \Cref{thm:ellcurves:jinv}(ii) and \Cref{thm:ellcurves:cplxuniform}(ii).
\end{proof}

\begin{definition}[the period lattice of a curve over $\Q$]\label{def:ellcurves:periodlattice}
Let $E/\Q$ be given by its global minimal model (\Cref{ch:reduction}) and $\omega=dx/(2y+a_1x+a_3)$ its N\'eron differential. The \emph{period lattice} of $E$ is
$\Lambda_E=\{\int_\gamma\omega:\gamma\in H_1(E(\C),\Z)\}$, so that $E(\C)\simeq\C/\Lambda_E$ with $\omega\leftrightarrow dz$ (\Cref{thm:ellcurves:cplxuniform}(ii)). Its \emph{area} is
$\vol(\Lambda_E)$, and $\omega_1(E)$ is the least positive real element of $\Lambda_E$ (the \emph{real period}).
\end{definition}

\begin{proposition}[periods of curves over $\Q$]\label{prop:ellcurves:periods}
Let $E/\Q$ be as in \Cref{def:ellcurves:periodlattice}.
\begin{enumerate}[label=(\roman*)]
\item $\Lambda_E$ is stable under complex conjugation. If $\Delta>0$, $\Lambda_E=\Z\omega_1+\Z\omega_2$ with $\omega_1>0$ and $\omega_2\in i\R$ (a rectangular lattice; $E(\R)$ has two components). If
$\Delta<0$, $\Lambda_E=\Z\omega_1+\Z\omega_2$ with $\omega_1>0$ and $\operatorname{Re}\omega_2=\omega_1/2$ ($E(\R)$ is connected).
\item $\vol(\Lambda_E)=\omega_1\operatorname{Im}\omega_2$ for such a basis with $\operatorname{Im}\omega_2>0$, and $\vol(\Lambda_E)=\frac i2\int_{E(\C)}\omega\wedge\bar\omega$.
\item The periods are computed by the arithmetic-geometric mean $\operatorname{AGM}(a,b)$, the common limit of $a_{n+1}=(a_n+b_n)/2$, $b_{n+1}=\sqrt{a_nb_n}$ for $a_0=a>0$,
$b_0=b>0$. If $\Delta<0$, let $e_1$ be the real root of $F(x)=4x^3+b_2x^2+2b_4x+b_6$, $\gamma=3e_1+b_2/4$ and $\delta=\sqrt{3e_1^2+b_2e_1/2+b_4/2}$; then
\[
\omega_1=\frac{2\pi}{\operatorname{AGM}\bigl(2\sqrt\delta,\sqrt{2\delta+\gamma}\bigr)},\qquad\operatorname{Im}\omega_2=\frac{\pi}{\operatorname{AGM}\bigl(2\sqrt\delta,\sqrt{2\delta-\gamma}\bigr)}
\]
for the basis of (i) with $\operatorname{Im}\omega_2>0$ \cite[Section 7.4]{Cohen93}, \cite[Section 3.7]{Cremona97}. In PARI the periods are \texttt{E.omega}.
\item If $\varphi\colon E\to E'$ is an isogeny over $\Q$ and $\omega'$ is the N\'eron differential of $E'$, then $\varphi^*\omega'=c_\varphi\,\omega$ with $c_\varphi\in\Q^\times$, and $z\mapsto c_\varphi z$ maps $\Lambda_E$ into
$\Lambda_{E'}$; so the period lattices of isogenous curves are commensurable.
\end{enumerate}
\end{proposition}

\begin{proof}[Sketch of proof]
(i) Complex conjugation acts on $E(\C)$ and on $H_1(E(\C),\Z)$, and $\int_{\bar\gamma}\omega=\overline{\int_\gamma\omega}$ because $\omega$ has real coefficients; so $\bar\Lambda_E=\Lambda_E$. For a conjugation-stable
lattice, $\Lambda\cap\R=\Z\omega_1$ and $\Lambda\cap i\R=\Z i\mu$ with $\omega_1,\mu>0$, and $2\lambda=(\lambda+\bar\lambda)+(\lambda-\bar\lambda)\in\Z\omega_1+\Z i\mu$ for every $\lambda\in\Lambda$, so $\Lambda$ contains $\Z\omega_1+\Z i\mu$
with index $1$ or $2$: rectangular, or with basis $\omega_1$, $\frac{\omega_1}2+\frac{i\mu}2$. Which case occurs is decided by the number of real roots of $F$, that is, by the sign of
$\Delta$ \cite[Section 3.7]{Cremona97}. (ii) The area of the parallelogram; and $\frac i2dz\wedge d\bar z=dx\wedge dy$ on $\C$. (iii) We do not prove the AGM formula; see the references given, and
\Cref{ex:ellcurves:ex-agm} for the numbers of \curve{11a1}. (iv) $\varphi^*\omega'$ is a holomorphic differential defined over $\Q$,
hence a rational multiple of $\omega$, nonzero because $\varphi$ is separable; under the uniformisations $\varphi$ becomes $z\mapsto c_\varphi z$ (\Cref{prop:ellcurves:cplxisog}(i)). For a modular parametrisation in place of $\varphi$ the
analogous constant is the subject of the Manin constant (\Cref{ch:moddeg}).
\end{proof}

\begin{example}[the periods of \curve{11a1}]\label{ex:ellcurves:11a1-periods}
For \curve{11a1}, $\Delta=-11^5<0$, and PARI gives (to $20$ digits)
\[
\omega_1=1.2692093042795534217\ldots,
\]
\[
\omega_2=\tfrac12\omega_1+1.4588166169384952294\ldots\,i,
\]
so that
\[
\vol(\Lambda_E)=\omega_1\operatorname{Im}\omega_2=1.8515436234559593178\ldots .
\]
The value $\omega_1$ agrees with the column $\Omega=1.26920930427955$ of Cremona's \texttt{allbsd} file. The point
$\tau=\omega_2/\omega_1=\frac12+1.1493901061232523807\ldots\,i$ is equivalent under $\tau\mapsto\tau-1$ to $-\frac12+1.14939\ldots\,i$, which lies in the standard fundamental domain
$\abs{\operatorname{Re}\tau}\le\frac12$, $\abs\tau\ge1$, and $j(\tau)=-757.672637860056752\ldots=-122023936/161051=j(\curve{11a1})$ to $18$ significant digits, the precision to which $\tau$ is given. (PARI returns the
basis $(\omega_1,\frac12\omega_1-1.4588\ldots i)$, which is not oriented.) The area enters the formula $\degphi=4\pi^2c_E^2\langle f,f\rangle/\vol(\Lambda_E)$ of \Cref{ch:moddeg}, where for
\curve{11a1} the modular degree is $1$.
\end{example}

\section{Elliptic curves over finite fields}\label{sec:ellcurves:finite}

The numbers $a_q$ of the paper are point counts over finite fields, and its local factors at good primes and at the type $\kod{I}_0^*$ primes are built from the roots $\alpha,\beta$ of
$X^2-a_qX+q$. This section proves the Hasse bound, the characteristic polynomial of Frobenius on the Tate module, and the identity
$(\mathfrak q+1)^2-a^2=\#E(\F_{\mathfrak q})\cdot\#E^{\mathrm{tw}}(\F_{\mathfrak q})$ behind the type $\kod{I}_0^*$ row of the paper's table, and it connects $a_p$ with $\rho_{E,\ell}(\Frob_p)$ for a
curve over $\Q$ with good reduction at $p$.

\begin{definition}[the Frobenius endomorphism and $a_{\mathfrak q}$]\label{def:ellcurves:frobendo}
Let $E$ be an elliptic curve over $\F_{\mathfrak q}$, $\mathfrak q=p^r$. The \emph{Frobenius endomorphism} is $\pi=\pi_{\mathfrak q}\colon(x,y)\mapsto(x^{\mathfrak q},y^{\mathfrak q})$, and
\[
a_{\mathfrak q}=a_{\mathfrak q}(E)=\mathfrak q+1-\#E(\F_{\mathfrak q}).
\]
\end{definition}

\begin{proposition}[properties of $\pi$]\label{prop:ellcurves:frobendo}
\leavevmode
\begin{enumerate}[label=(\roman*)]
\item $\pi$ is an endomorphism of $E$ of degree $\mathfrak q$, purely inseparable, and it acts on $E(\Fbar_p)$ as the arithmetic Frobenius $\Frob_{\mathfrak q}$.
\item $E(\F_{\mathfrak q^n})=\ker(1-\pi^n)$ for $n\ge1$.
\item $1-\pi$ is separable, so $\#E(\F_{\mathfrak q})=\deg(1-\pi)$.
\end{enumerate}
\end{proposition}

\begin{proof}
(i) $\pi$ maps $E$ to itself because the coefficients of $E$ lie in $\F_{\mathfrak q}$ and are fixed by $t\mapsto t^{\mathfrak q}$; it fixes $O$. Its degree, inseparability and its action on points are
\Cref{ch:curves} (the Frobenius morphism). (ii) A point of $E(\Fbar_p)$ is fixed by $\pi^n$ iff its coordinates are fixed by $t\mapsto t^{\mathfrak q^n}$, iff they lie in $\F_{\mathfrak q^n}$. (iii)
$\pi^*(dx)=d(x^{\mathfrak q})=\mathfrak qx^{\mathfrak q-1}dx=0$, so $\pi^*\omega=0$ and, by \Cref{prop:ellcurves:invdiff}(iii), $(1-\pi)^*\omega=\omega\ne0$; so $1-\pi$ is separable (\Cref{ch:curves}) and
$\#\ker(1-\pi)=\deg(1-\pi)$ (\Cref{thm:ellcurves:isog-hom}(iii)).
\end{proof}

\begin{theorem}[Hasse]\label{thm:ellcurves:hasse}
$\abs{a_{\mathfrak q}}\le2\sqrt{\mathfrak q}$.
\end{theorem}

\begin{proof}
By \Cref{thm:ellcurves:dual}(iv) the pairing $\langle\varphi,\psi\rangle=\deg(\varphi+\psi)-\deg\varphi-\deg\psi$ is bilinear on $\End(E)$, and $\deg(1-\pi)=1+\mathfrak q-\langle1,\pi\rangle$, so
$\langle1,\pi\rangle=1+\mathfrak q-\#E(\F_{\mathfrak q})=a_{\mathfrak q}$. For $m,n\in\Z$,
\[
0\le\deg(m-n\pi)=m^2\deg1+n^2\deg\pi-mn\langle1,\pi\rangle=m^2-a_{\mathfrak q}mn+\mathfrak qn^2 .
\]
A binary quadratic form that is nonnegative on $\Z^2$ is nonnegative on $\Q^2$ (by homogeneity) and on $\R^2$ (by continuity), so its discriminant $a_{\mathfrak q}^2-4\mathfrak q$ is at most $0$.
\end{proof}

\begin{theorem}[the characteristic polynomial of Frobenius]\label{thm:ellcurves:charpoly}
Let $\ell\ne p$.
\begin{enumerate}[label=(\roman*)]
\item For every $\varphi\in\End(E)$, the induced endomorphism $\varphi_\ell$ of $T_\ell E$ has $\det\varphi_\ell=\deg\varphi$ and $\tr\varphi_\ell=1+\deg\varphi-\deg(1-\varphi)$; both lie in $\Z$.
\item $\pi_\ell$ has characteristic polynomial $X^2-a_{\mathfrak q}X+\mathfrak q$, and $\pi^2-a_{\mathfrak q}\pi+\mathfrak q=0$ in $\End(E)$.
\item Let $\alpha,\beta\in\C$ be the roots of $X^2-a_{\mathfrak q}X+\mathfrak q$. Then $\alpha+\beta=a_{\mathfrak q}$, $\alpha\beta=\mathfrak q$, $\abs\alpha=\abs\beta=\sqrt{\mathfrak q}$, and for every $n\ge1$
\[
\#E(\F_{\mathfrak q^n})=\mathfrak q^n+1-\alpha^n-\beta^n;\quad\text{in particular}\quad\#E(\F_{\mathfrak q^2})=\mathfrak q^2+1-(a_{\mathfrak q}^2-2\mathfrak q).
\]
Equivalently, the zeta function of $E$ is
\[
Z(E/\F_{\mathfrak q},T)=\exp\Bigl(\sum_{n\ge1}\#E(\F_{\mathfrak q^n})\frac{T^n}n\Bigr)=\frac{1-a_{\mathfrak q}T+\mathfrak qT^2}{(1-T)(1-\mathfrak qT)} .
\]
\end{enumerate}
\end{theorem}

\begin{proof}
(i) The Weil pairing $e$ on $T_\ell E$ (\Cref{prop:ellcurves:tate}(iii)) satisfies $e(\varphi s,\varphi t)=e(s,\hat\varphi\varphi t)=e(s,t)^{\deg\varphi}$ by \Cref{thm:ellcurves:weil}(vi) and
\Cref{thm:ellcurves:dual}. For any $\Z_\ell$-linear $g$ on a free module of rank $2$ and an alternating bilinear form $e$, $e(gs,gt)=e(s,t)^{\det g}$ (check on a basis, as in
\Cref{cor:ellcurves:det}). Since $e$ is perfect, some $e(s,t)$ generates $\Z_\ell(1)$, so $\det\varphi_\ell=\deg\varphi$. For a $2\times2$ matrix $g$, $\det(1-g)=1-\tr g+\det g$, and $(1-\varphi)_\ell=1-\varphi_\ell$.
(ii) By (i) and \Cref{prop:ellcurves:frobendo}(iii), $\det\pi_\ell=\deg\pi=\mathfrak q$ and $\tr\pi_\ell=1+\mathfrak q-\#E(\F_{\mathfrak q})=a_{\mathfrak q}$. By Cayley-Hamilton, $\pi^2-a_{\mathfrak q}\pi+\mathfrak q$ acts as $0$ on
$T_\ell E$, hence kills every $E[\ell^n]$; a nonzero endomorphism has finite kernel (\Cref{thm:ellcurves:isog-hom}(ii)), so it is $0$.
(iii) The discriminant $a_{\mathfrak q}^2-4\mathfrak q\le0$ (\Cref{thm:ellcurves:hasse}), so $\alpha,\beta$ are complex conjugate, and $\alpha\beta=\mathfrak q$ gives $\abs\alpha^2=\mathfrak q$. As $\pi^n$ is the Frobenius
of $\F_{\mathfrak q^n}$, $\#E(\F_{\mathfrak q^n})=\deg(1-\pi^n)=\det(1-\pi_\ell^n)=(1-\alpha^n)(1-\beta^n)$, the eigenvalues of $\pi_\ell^n$ being $\alpha^n,\beta^n$. For $n=2$, $\alpha^2+\beta^2=a_{\mathfrak q}^2-2\mathfrak q$. Finally
$\sum_n(1+\mathfrak q^n-\alpha^n-\beta^n)T^n/n=-\log(1-T)-\log(1-\mathfrak qT)+\log(1-\alpha T)+\log(1-\beta T)$.
\end{proof}

\noindent For \curve{11a1}, with $a_3=-1$, $a_5=1$, $a_7=-2$, $a_{13}=4$, the formula gives $\#E(\F_{p^2})=15,35,60,180$ for $p=3,5,7,13$, and a direct count over $\F_{p^2}$ with PARI
gives the same numbers.

\begin{definition}[supersingular and ordinary]\label{def:ellcurves:supersingular}
An elliptic curve $E$ over $\F_{\mathfrak q}$ is \emph{supersingular} if $E[p]=0$, and \emph{ordinary} otherwise.
\end{definition}

\begin{proposition}[criteria for supersingularity]\label{prop:ellcurves:supersingular}
\leavevmode
\begin{enumerate}[label=(\roman*)]
\item $E$ is supersingular iff $p\mid a_{\mathfrak q}$, iff $\End_{\Fbar_p}(E)$ is an order in a quaternion algebra.
\item If $\mathfrak q=p\ge5$, then $E$ is supersingular iff $a_p=0$.
\end{enumerate}
\end{proposition}

\begin{proof}
(i) We do not prove this; see \cite[Section V.3]{SilvermanAEC}. (ii) By (i) supersingularity means $p\mid a_p$, and $\abs{a_p}\le2\sqrt p<p$ for $p\ge5$ (\Cref{thm:ellcurves:hasse}), so $p\mid a_p$ iff $a_p=0$.
\end{proof}

\begin{proposition}[point counts of a quadratic twist over $\F_{\mathfrak q}$]\label{prop:ellcurves:twistFq}
Let $\mathfrak q$ be odd, $E\colon y^2=f(x)$ with $f$ a cubic over $\F_{\mathfrak q}$ with distinct roots (so that $E$ is smooth), $\theta\in\F_{\mathfrak q}^\times$ a non-square, and $E^{\mathrm{tw}}\colon\theta y^2=f(x)$, which is isomorphic to
$y^2=\theta^3f(x/\theta)$, the twist of $E$ by $\theta$ in the sense of \Cref{def:ellcurves:twist}. Let $\eta$ be the quadratic character of $\F_{\mathfrak q}^\times$, extended by $\eta(0)=0$. Then
\[
\#E(\F_{\mathfrak q})=\mathfrak q+1+\sum_{x\in\F_{\mathfrak q}}\eta(f(x)),\qquad a_{\mathfrak q}(E^{\mathrm{tw}})=-a_{\mathfrak q}(E),\qquad\#E(\F_{\mathfrak q})+\#E^{\mathrm{tw}}(\F_{\mathfrak q})=2\mathfrak q+2,
\]
\[
(\mathfrak q+1)^2-a_{\mathfrak q}(E)^2=\#E(\F_{\mathfrak q})\cdot\#E^{\mathrm{tw}}(\F_{\mathfrak q}).
\]
\end{proposition}

\begin{proof}
For each $x$, the number of $y\in\F_{\mathfrak q}$ with $y^2=c$ is $1+\eta(c)$, and with $\theta y^2=c$ it is $1+\eta(c/\theta)=1-\eta(c)$. Adding the point $O$, $\#E(\F_{\mathfrak q})=1+\sum_x(1+\eta(f(x)))$ and
$\#E^{\mathrm{tw}}(\F_{\mathfrak q})=1+\sum_x(1-\eta(f(x)))$. So $a_{\mathfrak q}(E)=-\sum_x\eta(f(x))=-a_{\mathfrak q}(E^{\mathrm{tw}})$, the sum of the two counts is $2\mathfrak q+2$, and their product is
$(\mathfrak q+1-a_{\mathfrak q})(\mathfrak q+1+a_{\mathfrak q})$. The isomorphism: multiplying $\theta y^2=f(x)$ by $\theta^3$ and putting $Y=\theta^2y$, $X=\theta x$ gives $Y^2=\theta^3f(X/\theta)$.
\end{proof}

\begin{theorem}[Frobenius at a prime of good reduction]\label{thm:ellcurves:frobQ}
Let $E/\Q$ have good reduction at $p$ (its minimal model has $p\nmid\Delta$), let $\tilde E$ be the reduction, $a_p=a_p(E)=p+1-\#\tilde E(\F_p)$, and $\ell\ne p$.
\begin{enumerate}[label=(\roman*)]
\item $\rho_{E,\ell}$ is unramified at $p$, and $\rho_{E,\ell}(\Frob_p)$ has characteristic polynomial $X^2-a_pX+p$, for every choice of arithmetic Frobenius $\Frob_p$; in particular
$\tr\rho_{E,\ell}(\Frob_p)=a_p$ does not depend on $\ell$.
\item $\bar\rho_{E,\ell}(\Frob_p)$ has characteristic polynomial $X^2-a_pX+p$ modulo $\ell$, so $\tr\bar\rho_{E,\ell}(\Frob_p)\equiv a_p\pmod\ell$ and $\det\bar\rho_{E,\ell}(\Frob_p)\equiv p\pmod\ell$.
\end{enumerate}
We write $\alpha_p$, $\beta_p$ for the roots of $X^2-a_pX+p$.
\end{theorem}

\begin{proof}[Sketch of proof]
Fix an embedding $\Qbar\to\Qbar_p$; it defines a decomposition group $D_p\simeq G_{\Q_p}$ and its inertia group $I_p$ (\Cref{ch:ant,ch:local}). The points of $E[\ell^n]$ have
coordinates in the finite extension $L=\Q_p(E[\ell^n])$ of $\Q_p$ inside $\Qbar_p$, and reduction modulo the maximal ideal of the valuation ring of $L$ (equivalently of $\Qbar_p$, with
residue field $\bar\F_p$) maps $E[\ell^n]$ injectively into $\tilde E[\ell^n]$ (\Cref{ch:reduction}: the kernel of reduction has no torsion prime to $p$); both have $\ell^{2n}$ elements (\Cref{cor:ellcurves:degm}(ii)),
so it is a bijection, equivariant for $D_p\to G_{\F_p}$. Hence $I_p$ acts trivially on $E[\ell^n]$, and $\Frob_p$ acts as $\Frob_p\in G_{\F_p}$ does on $\tilde E[\ell^n]$, that is, as $\pi$
(\Cref{prop:ellcurves:frobendo}(i)). By \Cref{thm:ellcurves:charpoly}(ii) the characteristic polynomial is $X^2-a_pX+p$. Part (ii) is the reduction of (i). This is the easy half of the criterion of
N\'eron, Ogg and Shafarevich (\Cref{ch:galrep}); see \cite[Section VII.7]{SilvermanAEC}.
\end{proof}

\begin{example}[the numbers $a_p$ for the running examples]\label{ex:ellcurves:aptable}
Point counts on the global minimal models, $p<40$; at a bad prime $a_p=p+1-\#\tilde E(\F_p)$ counting the singular point, which gives $1$, $-1$ or $0$ for split multiplicative, nonsplit
multiplicative and additive reduction (\Cref{ch:reduction}).
\begin{center}\small
\begin{tabular}{lrrrrrrrrrrrr}
\toprule
curve & $a_2$ & $a_3$ & $a_5$ & $a_7$ & $a_{11}$ & $a_{13}$ & $a_{17}$ & $a_{19}$ & $a_{23}$ & $a_{29}$ & $a_{31}$ & $a_{37}$\\
\midrule
\curve{11a1} & $-2$ & $-1$ & $1$ & $-2$ & $1$ & $4$ & $-2$ & $0$ & $-1$ & $0$ & $7$ & $3$\\
\curve{20a1} & $0$ & $-2$ & $-1$ & $2$ & $0$ & $2$ & $-6$ & $-4$ & $6$ & $6$ & $-4$ & $2$\\
\curve{99a1} & $-1$ & $0$ & $-4$ & $-2$ & $-1$ & $-2$ & $2$ & $-6$ & $4$ & $-6$ & $4$ & $-6$\\
\curve{121a1} & $-1$ & $2$ & $1$ & $2$ & $0$ & $-1$ & $5$ & $-6$ & $2$ & $-9$ & $-2$ & $-3$\\
\curve{121d1} & $2$ & $-1$ & $1$ & $2$ & $0$ & $-4$ & $2$ & $0$ & $-1$ & $0$ & $7$ & $3$\\
\curve{539c1} & $1$ & $-2$ & $2$ & $0$ & $1$ & $-4$ & $-4$ & $0$ & $-4$ & $-6$ & $-10$ & $-6$\\
\curve{726a1} & $-1$ & $-1$ & $0$ & $0$ & $0$ & $-6$ & $6$ & $-6$ & $6$ & $-6$ & $4$ & $-2$\\
\curve{980i1} & $0$ & $-3$ & $1$ & $0$ & $-5$ & $3$ & $1$ & $-6$ & $6$ & $-9$ & $4$ & $2$\\
\bottomrule
\end{tabular}
\end{center}
Two patterns can be read off. For \curve{726a1}, whose mod $3$ image is $N_s(\F_3)$ and so satisfies $\bar\rho\simeq\bar\rho\otimes\chi_{-3}$, $a_p\equiv0\pmod3$ at the good $p\equiv2\pmod3$
($a_5=0$, $a_{17}=6$, $a_{23}=6$, $a_{29}=-6$), by \Cref{thm:ellcurves:frobQ}(ii) and $\chi_{-3}(\Frob_p)=-1$. For \curve{11a1}, $a_p\equiv1+p\pmod5$ at every good $p$ (\Cref{prop:ellcurves:zlmul}). The
row of \curve{121d1} is $\chi_{-11}(p)$ times that of \curve{11a1} for $p\ne11$ (\Cref{cor:ellcurves:twistap}). The coefficients of \curve{11a1} agree with those of
$q\prod_{n\ge1}(1-q^n)^2(1-q^{11n})^2$, $q=e^{2\pi i\tau}$, up to $q^{20}$ (\Cref{ch:modforms}).
\end{example}

\begin{paperlink}\label{pl:ellcurves:table}
In the table of Section 1 and in Lemma lem:euler, $E'$ is the twist of good reduction at a type $\kod{I}_0^*$ prime $q$, and $\alpha,\beta$ are the roots of $X^2-a_q(E')X+q$; the proof uses
``$\alpha\beta=q$ and $\alpha+\beta=a_q(E')$'' (\Cref{thm:ellcurves:charpoly}(iii)), and that $\ad V$ has the Frobenius eigenvalues $\alpha/\beta,1,\beta/\alpha$ for either Frobenius
(\Cref{conv:ellcurves:frobenius}(d)). By \Cref{prop:ellcurves:twistFq}, the factor $(q+1)^2-a_q(E')^2$ of $e_\ell(E,q)$ is $\#E'(\F_q)\cdot\#E'^{\mathrm{tw}}(\F_q)$, the product of the numbers of points of
the reduction of $E'$ and of its unramified quadratic twist; so $v_\ell((q+1)^2-a_q(E')^2)$ is the sum of the $\ell$-adic valuations of two point counts. Lemma lem:old uses that the characteristic
polynomial of $\bar\rho_h(\Frob_p)$ is $X^2-a_p(h)X+p$, the analogue of \Cref{thm:ellcurves:frobQ}(ii) for a newform $h$ (\Cref{ch:modgalrep}).
\end{paperlink}
\section[Curves over the rationals: torsion and isogenies]{Elliptic curves over the rationals: Mordell-Weil, torsion, isogenies}\label{sec:ellcurves:overQ}

The paper's Eisenstein case (Conjecture conj:eis, Section sec:eis, Proposition prop:no9) is about rational torsion points and rational isogenies, and its
remarks on optimality (Section 1, Section sec:oot) are about isogeny classes. This section states the theorems of Mordell and Weil, of Mazur and of Kenku,
proves that $E[\ell]$ is reducible exactly when $E$ has a rational $\ell$-isogeny, describes the case $E[\ell]\simeq\Z/\ell\oplus\mu_\ell$, proves the paper's Proposition
prop:no9, and proves that when $E[\ell]$ is irreducible nothing $\ell$-adic depends on which curve of the isogeny class one chooses.

\begin{theorem}[Mordell, Weil]\label{thm:ellcurves:mordellweil}
For $E/\Q$, the group $E(\Q)$ is finitely generated: $E(\Q)\simeq\Z^r\oplus E(\Q)_{\mathrm{tors}}$ with $E(\Q)_{\mathrm{tors}}$ finite. The integer $r$ is the \emph{rank} of $E$.
\end{theorem}

\noindent We do not prove this; see \cite[Chapter VIII]{SilvermanAEC}. The proof has two halves: the weak Mordell-Weil theorem, that $E(\Q)/2E(\Q)$ is finite, which
comes from the Kummer map into a Selmer group (\Cref{ch:galcoh}); and a descent with the canonical height, which is a quadratic form on $E(\Q)$ with finitely many
points below any bound.

\begin{proposition}[torsion and reduction]\label{prop:ellcurves:torsred}
Let $E/\Q$ be given by its minimal model and $p$ a prime of good reduction.
\begin{enumerate}[label=(\roman*)]
\item Reduction modulo $p$ is injective on $E(\Q)[m]$ for every $m$ prime to $p$, and, if $p\ge3$, on all of $E(\Q)_{\mathrm{tors}}$. Hence $\#E(\Q)_{\mathrm{tors}}$ divides $\#\tilde E(\F_p)$ for every
good prime $p\ge3$.
\item (Nagell and Lutz.) A torsion point of a model with integral coefficients and $a_1=a_3=0$ has integral coordinates.
\item For \curve{11a1}, $E(\Q)_{\mathrm{tors}}=\langle(5,5)\rangle\simeq\Z/5$; for \curve{20a1}, $E(\Q)_{\mathrm{tors}}=\langle(4,10)\rangle\simeq\Z/6$.
\end{enumerate}
\end{proposition}

\begin{proof}
(i) is \Cref{ch:reduction} (the kernel of reduction is the formal group, which has no torsion prime to $p$, and no $p$-torsion for $p\ge3$ over $\Z_p$); see
\cite[Sections VII.3 and VIII.7]{SilvermanAEC}. (ii) We do not prove this; see \cite[Section VIII.7]{SilvermanAEC}. (iii) For \curve{11a1}, $\#\tilde E(\F_3)=5$ and $\#\tilde E(\F_7)=10$ (from
$a_3=-1$, $a_7=-2$, \Cref{ex:ellcurves:aptable}), so $\#E(\Q)_{\mathrm{tors}}$ divides $5$, and $(5,5)$ has order $5$ (\Cref{ex:ellcurves:torsion-points}). For \curve{20a1}, $\#\tilde E(\F_3)=6$ and
$\#\tilde E(\F_7)=6$, so the torsion has order dividing $6$, and $(4,10)$ has order $6$.
\end{proof}

\begin{theorem}[Mazur]\label{thm:ellcurves:mazur}
For $E/\Q$, $E(\Q)_{\mathrm{tors}}$ is isomorphic to one of
\[
\Z/n\ (1\le n\le10\text{ or }n=12),\qquad\Z/2\times\Z/2n\ (1\le n\le4),
\]
and each of these fifteen groups occurs for infinitely many $j$-invariants.
\end{theorem}

\noindent We do not prove this; see \cite{Mazur77} and \cite[Section VIII.7]{SilvermanAEC}. In Cremona's tables for conductor below $10{,}000$ the torsion orders that occur are
$1,\dots,10$, $12$ and $16$ (the last only as $\Z/2\times\Z/8$). Consequences used by the paper: $v_\ell(\#E(\Q)_{\mathrm{tors}})\le1$ for $\ell=5,7$, $E(\Q)$ has no point of order $\ell\ge11$, and
$v_3(\#E(\Q)_{\mathrm{tors}})\le2$ with equality only for $\Z/9$.

\begin{definition}[isogeny classes]\label{def:ellcurves:isogclass}
Curves $E,E'$ over $\Q$ are \emph{isogenous} if there is a nonconstant isogeny $E\to E'$ defined over $\Q$. The equivalence classes (\Cref{thm:ellcurves:isogclass}(i)) are the
\emph{isogeny classes}, labelled by Cremona as $N$ followed by a letter code (\Cref{sec:ellcurves:labels}). Cremona's \emph{isogeny matrix} of a class lists, for each pair of curves, the
least degree of a cyclic isogeny over $\Q$ between them.
\end{definition}

\begin{theorem}[isogeny classes]\label{thm:ellcurves:isogclass}
\leavevmode
\begin{enumerate}[label=(\roman*)]
\item Being isogenous over $\Q$ is an equivalence relation.
\item Isogenous curves $E$, $E'$ have $V_\ell E\simeq V_\ell E'$ for every $\ell$; hence the same $a_p$ at every prime of good reduction for both, the same $L$-function
(\Cref{sec:ellcurves:Lfunction}), and the same rank. They have the same conductor. In general they have different torsion, discriminants, $j$-invariants, Tamagawa numbers and periods.
\item (Faltings.) Conversely, $E$ and $E'$ are isogenous over $\Q$ iff $a_p(E)=a_p(E')$ for almost all $p$, iff $V_\ell E\simeq V_\ell E'$ as $G_\Q$-modules for one $\ell$.
\item (Kenku.) An isogeny class over $\Q$ has at most $8$ curves (up to isomorphism over $\Q$).
\end{enumerate}
\end{theorem}

\begin{proof}
(i) Reflexive by $[1]$, transitive by composition, symmetric by the dual isogeny, defined over $\Q$ (\Cref{thm:ellcurves:dual}). (ii) $V_\ell E\simeq V_\ell E'$ is
\Cref{prop:ellcurves:rho}(iii); traces of Frobenius give the $a_p$ (\Cref{thm:ellcurves:frobQ}). For the rank: $\varphi\colon E(\Q)\to E'(\Q)$ and $\hat\varphi$ compose to $[\deg\varphi]$, which has finite
kernel, so $\rank E(\Q)\le\rank E'(\Q)$, and symmetrically. The conductor depends only on $V_\ell E$ as a representation of the local Galois groups (\Cref{ch:reduction,ch:galrep}); we do not
prove this here. The examples of \Cref{sec:ellcurves:labels} show the differences. (iii) We do not prove this; see \cite{Faltings83}. (iv) We do not prove this; see \cite[Theorem 2]{Kenku}.
\end{proof}

\noindent Examples (Cremona's \texttt{allisog} file; rows and columns ordered by curve number): \curve{11a} has matrix $\left(\begin{smallmatrix}1&5&5\\5&1&25\\5&25&1\end{smallmatrix}\right)$;
\curve{20a} has $\left(\begin{smallmatrix}1&2&3&6\\2&1&6&3\\3&6&1&2\\6&3&2&1\end{smallmatrix}\right)$; \curve{121a} and \curve{121c} consist of two curves joined by an $11$-isogeny;
\curve{121d} has $\left(\begin{smallmatrix}1&5&25\\5&1&5\\25&5&1\end{smallmatrix}\right)$; \curve{99a}, \curve{539c} and \curve{726a} consist of two curves joined by a $2$-isogeny; \curve{980i} has one curve. Below
conductor $10{,}000$ the classes have $1,2,3,4,6$ or $8$ curves.

\begin{proposition}[rational $\ell$-isogenies and reducibility]\label{prop:ellcurves:isogline}
Let $\ell$ be a prime, $\operatorname{char}K\ne\ell$, and $E/K$.
\begin{enumerate}[label=(\roman*)]
\item The isogenies of degree $\ell$ from $E$ defined over $K$, up to isomorphism of the target, correspond bijectively to the $G_K$-stable lines of $E[\ell]$, by $\varphi\mapsto\ker\varphi$.
\item $E[\ell]$ is reducible as an $\F_\ell[G_K]$-module iff $E$ has a $K$-rational $\ell$-isogeny. A $K$-rational point of order $\ell$ spans such a line; the converse fails.
\item If $E[\ell]$ is reducible, with $G_K$ acting on a stable line by $\eta_1$ and on the quotient by $\eta_2$, then $\bar\rho_{E,\ell}^{\mathrm{ss}}\simeq\eta_1\oplus\eta_2$ and $\eta_1\eta_2=\bar\chi_\ell$.
\end{enumerate}
\end{proposition}

\begin{proof}
(i) An isogeny of degree $\ell$ is separable, since its inseparable degree is a power of $\operatorname{char}K$ dividing $\ell$; so its kernel has order $\ell$ (\Cref{thm:ellcurves:isog-hom}(iii)) and
is a line in $E[\ell]$, $G_K$-stable if the isogeny is defined over $K$ (\Cref{def:ellcurves:isogeny}). Conversely a $G_K$-stable line $L$ gives $E\to E/L$ over $K$, and two separable
isogenies with the same kernel differ by an isomorphism of the targets (\Cref{thm:ellcurves:quotient}). (ii) $E[\ell]$ is two-dimensional, so it is reducible iff it has a stable line; a
rational point of order $\ell$ is fixed by $G_K$, so its span is stable. For $K=\Q$: \curve{11a2} has a rational $5$-isogeny (its class is \curve{11a}) and trivial torsion; \curve{539a1}
has a rational $3$-isogeny (image \texttt{3B}) and trivial torsion. (iii) is \Cref{cor:ellcurves:det}(iii).
\end{proof}

\noindent As a check on \Cref{prop:ellcurves:isogline} against the data: for $\ell=3,5,7$ and curve number $1$ of every class without CM below conductor $10{,}000$, Cremona's
\texttt{galrep} file lists a Borel or split Cartan image at $\ell$ (\Cref{ch:galrep}) exactly when the row of the isogeny matrix contains a degree divisible by $\ell$.

\begin{theorem}[Mazur, Kenku: rational cyclic isogenies]\label{thm:ellcurves:mazurkenku}
Let $E/\Q$.
\begin{enumerate}[label=(\roman*)]
\item If $E$ has a rational isogeny of prime degree $\ell$, then $\ell\in\{2,3,5,7,11,13,17,19,37,43,67,163\}$.
\item If $E$ has a rational cyclic isogeny of degree $n$, then $n\le19$ or $n\in\{21,25,27,37,43,67,163\}$.
\item The degrees $n=19,27,43,67,163$ occur only for curves with complex multiplication. For $n=27$ the curves have $j=-2^{15}\cdot3\cdot5^3$ and CM by the order of discriminant $-27$.
\end{enumerate}
\end{theorem}

\noindent We do not prove this. Part (i) is Mazur's theorem \cite{Mazur78}; (ii) and (iii) combine it with Kenku's determinations of the rational points of $X_0(N)$ for the
remaining $N$ \cite{Kenku79,Kenku80a,Kenku80b,Kenku81}, and the case $N=27$, where $X_0(27)$ has genus one, is Ligozat's determination of the rational points of the modular curves of genus one \cite{Ligozat75}. The complete list (ii) is \cite[Theorem 1]{Kenku}, which attributes the levels of genus one to Ligozat, Ogg and Joly; Kenku does not discuss complex multiplication, so (iii) is not in that paper. A summary table is in \cite{LozanoRobledo13} (we have not checked its numbering). In Cremona's tables the cyclic $27$-isogeny joins \curve{27a2} and \curve{27a4}, both with
$j=-12288000=-2^{15}\cdot3\cdot5^3$ (\Cref{thm:ellcurves:13j}), and the isogeny degrees occurring below conductor $10{,}000$ are $1$ to $16$, $18$, $19$, $21$, $25$, $27$, $37$, $43$ and $67$.

\begin{proposition}[the structure $\Z/\ell\oplus\mu_\ell$]\label{prop:ellcurves:zlmul}
Let $\ell$ be odd and let $E/\Q$ have a rational point $P_0$ of order $\ell$.
\begin{enumerate}[label=(\roman*)]
\item In a basis $(P_0,Q)$ of $E[\ell]$, $\bar\rho_{E,\ell}=\sm1*0{\bar\chi_\ell}$.
\item The following are equivalent: the sequence $0\to\Z/\ell\to E[\ell]\to\mu_\ell\to0$ of \Cref{cor:ellcurves:det}(iii) splits; $E[\ell]\simeq\Z/\ell\oplus\mu_\ell$ as $G_\Q$-modules; the image of
$\bar\rho_{E,\ell}$ is conjugate to $\{\operatorname{diag}(1,d):d\in\F_\ell^\times\}$ (Sutherland's label $\ell$\texttt{Cs.1.1}, \Cref{ch:groups,ch:galrep}); $E[\ell]$ has two $G_\Q$-stable lines, that is,
$E$ has two rational $\ell$-isogenies, with kernels $\Z/\ell$ and $\mu_\ell$.
\item $a_p\equiv1+p\pmod\ell$ and $\#\tilde E(\F_p)\equiv0\pmod\ell$ for every good prime $p\ne\ell$.
\end{enumerate}
\end{proposition}

\begin{proof}
(i) is \Cref{cor:ellcurves:det}(iii) with $\eta_1=1$. (ii) The sequence splits iff there is a stable complement to $\langle P_0\rangle$, iff in a suitable basis the matrices are diagonal,
$\operatorname{diag}(1,\bar\chi_\ell(\sigma))$; as $\bar\chi_\ell\colon G_\Q\to\F_\ell^\times$ is surjective, the image is then the whole group $\{\operatorname{diag}(1,d)\}$. A stable complement is a second stable line; since
$\ell$ is odd, $1\ne\bar\chi_\ell$, the two characters are distinct, and the stable lines are then exactly the two eigenlines, on which $G_\Q$ acts trivially ($\Z/\ell$) and by $\bar\chi_\ell$
($\mu_\ell$). Conversely two distinct stable lines give a splitting. Rational $\ell$-isogenies correspond to stable lines by \Cref{prop:ellcurves:isogline}. (iii) By (i) and
\Cref{thm:ellcurves:frobQ}(ii), $a_p\equiv\tr\sm1*0p=1+p$, and $\#\tilde E(\F_p)=p+1-a_p\equiv0$.
\end{proof}

\noindent Examples (all checked with PARI and the \texttt{galrep} file). The curves \curve{14a1}, \curve{19a1}, \curve{26a1} and the CM curve \curve{27a1} have $E[3]\simeq\Z/3\oplus\mu_3$ (image
\texttt{3Cs.1.1}); \curve{20a1} has a rational point of order $3$ with non-split $E[3]$ (image \texttt{3B.1.1}). The curve \curve{11a1} has $E[5]\simeq\Z/5\oplus\mu_5$ (image \texttt{5Cs.1.1}), and
$a_p\equiv1+p\pmod5$ at every good $p<500$. Its two $5$-isogenies have kernels $\Z/5=\langle(5,5)\rangle$ and $\mu_5$ (the points with $x$-coordinate a root of $5x^2+5x-29$, a factor of
$\psi_5$); V\'elu's formulas give $\curve{11a1}/\langle(5,5)\rangle\simeq\curve{11a2}$, with trivial torsion, and $\curve{11a1}/\mu_5\simeq\curve{11a3}$, which contains the rational group
$E[5]/\mu_5\simeq\Z/5$ and has torsion $\Z/5$.

\begin{proposition}[no rational point of order $9$ beside $\Z/3\oplus\mu_3$]\label{prop:ellcurves:no27}
Let $E/\Q$ be without complex multiplication, with $E[3]\simeq\Z/3\oplus\mu_3$. Then $E(\Q)$ has no point of order $9$; so $E(\Q)[3^\infty]\simeq\Z/3$. Consequently, for every $E/\Q$
without complex multiplication,
\[
v_3\bigl(\#E(\Q)_{\mathrm{tors}}\bigr)+\bigl[E[3]\simeq\Z/3\oplus\mu_3\bigr]\le2,
\]
where $[\cdot]$ is $1$ if the condition holds and $0$ otherwise; and $v_\ell(\#E(\Q)_{\mathrm{tors}})\le1$ for $\ell=5,7$, $=0$ for $\ell\ge11$.
\end{proposition}

\begin{proof}
Suppose $P\in E(\Q)$ has order $9$.

\emph{Step 1.} $3P$ is a rational point of order $3$. The rational points of $E[3]\simeq\Z/3\oplus\mu_3$ are $\Z/3\oplus\mu_3(\Q)=\Z/3\oplus0$, because $\mu_3(\Q)=\{1\}$ ($\chi_{-3}\ne1$). So
$\langle3P\rangle$ is the line $\Z/3$, and $\langle P\rangle\cap E[3]=\langle3P\rangle$ (the elements of order dividing $3$ in a cyclic group of order $9$) meets the line $\mu_3$ only in $O$.

\emph{Step 2.} The line $\mu_3$ is $G_\Q$-stable, so $\varphi\colon E\to E'=E/\mu_3$ is defined over $\Q$ (\Cref{thm:ellcurves:quotient}), of degree $3$; let $\hat\varphi\colon E'\to E$ be its dual, also
over $\Q$, of degree $3$, with $\varphi\circ\hat\varphi=[3]_{E'}$ (\Cref{thm:ellcurves:dual}). For $Q\in E'[3]$, $\varphi(\hat\varphi(Q))=3Q=O$, so $\hat\varphi(E'[3])\subseteq\ker\varphi=\mu_3$. The kernel of
$\hat\varphi$ has order $3$ and is killed by $3$, so it lies in $E'[3]$, and $\#\hat\varphi(E'[3])=9/3=3$. Hence $\hat\varphi(E'[3])=\mu_3$.

\emph{Step 3.} Let $C=\hat\varphi^{-1}(\langle P\rangle)\subset E'(\Qbar)$. It is a subgroup, $G_\Q$-stable because $\langle P\rangle$ is (it consists of rational points) and $\hat\varphi$ is defined over $\Q$,
and $\#C=\#\ker\hat\varphi\cdot\#\langle P\rangle=27$, because $\hat\varphi$ is surjective on $\Qbar$-points.

\emph{Step 4.} $C$ is cyclic. If $E'[3]\subseteq C$, then $\mu_3=\hat\varphi(E'[3])\subseteq\hat\varphi(C)=\langle P\rangle$, contradicting Step 1. By \Cref{prop:ellcurves:cyclicfactor}(i), $C\simeq\Z/a\times\Z/ab$
with $a^2b=27$ and $E'[a]\subseteq C$; so $a=1$ and $C\simeq\Z/27$.

\emph{Step 5.} $E'\to E'/C$ is a rational cyclic isogeny of degree $27$ (\Cref{thm:ellcurves:quotient}). By \Cref{thm:ellcurves:mazurkenku}(iii), $E'$ has complex multiplication, hence so
has $E$ (\Cref{lem:ellcurves:cm-isog}): a contradiction.

For the consequences: by \Cref{thm:ellcurves:mazur}, $v_3(\#E(\Q)_{\mathrm{tors}})\le2$ with equality only for $E(\Q)_{\mathrm{tors}}\simeq\Z/9$; in that case the bracket is $0$ by what was just proved,
and otherwise $v_3\le1$. The statements for $\ell\ge5$ are read off Mazur's list.
\end{proof}

\noindent Data: the $264$ curves of conductor below $10{,}000$ with image \texttt{3Cs.1.1} have torsion of order $3$ ($244$ curves) or $6$ ($20$ curves); the six curves there with a rational point of
order $9$ (\curve{54b3}, \curve{714i1}, \curve{1482l1}, \curve{1554n1}, \curve{1890r3}, \curve{6942n1}) all have image \texttt{3B.1.1}.

\begin{paperlink}\label{pl:ellcurves:eis}
\Cref{prop:ellcurves:no27} is the paper's Proposition prop:no9, with every step expanded; it turns the bound ``$D\le2$ for $\ell=3$'' of Conjecture conj:eis into a consequence of the first bound
$D\le v_\ell(\#E(\Q)_{\mathrm{tors}})+[E[3]\simeq\Z/3\oplus\mu_3]$, and gives ``$D\le1$ for $\ell\ge5$'' in the same way. The paper cites \cite{Ligozat75} and \cite{Kenku} for the fact that a rational cyclic
$27$-isogeny forces complex multiplication (\Cref{thm:ellcurves:mazurkenku}(iii)). The source is Ligozat's study of $X_0(27)$, a curve of genus one.
\cite{Kenku} is Kenku's paper on the number of curves in an isogeny class; its Theorem 1 lists $27$ among the levels $N$ with $Y_0(N)(\Q)\ne\emptyset$ and credits the levels of genus one to
Ligozat, Ogg and Joly, but it says nothing about complex multiplication, so it supports only the existence of the rational point, not the fact cited. In Section sec:eis, ``a rational $\ell$-isogeny'' and ``Borel type'' refer to \Cref{prop:ellcurves:isogline},
and the condition $E[3]\simeq\Z/3\oplus\mu_3$, which the paper identifies with the image \texttt{3Cs.1.1} (the subgroup of order $2$ of a split Cartan subgroup fixing a nonzero vector), is
\Cref{prop:ellcurves:zlmul}(ii).
The second bound of Conjecture conj:eis counts the curves of the isogeny class with a rational point of order $\ell$; in the class \curve{11a} at $\ell=5$ these are \curve{11a1} and \curve{11a3}.
\end{paperlink}

\begin{lemma}[isogenies prime to $\ell$]\label{lem:ellcurves:primetol}
Let $E/\Q$ and $\ell$ a prime with $E[\ell]$ irreducible.
\begin{enumerate}[label=(\roman*)]
\item Every isogeny $\varphi\colon E\to E'$ over $\Q$ is $\psi\circ[a]$ with $\psi\colon E\to E'$ a cyclic isogeny over $\Q$ of degree prime to $\ell$.
\item Hence $T_\ell E\simeq T_\ell E'$ as $G_\Q$-modules, $\rho_{E',\ell}$ and $\rho_{E,\ell}$ are conjugate in $\GL_2(\Z_\ell)$, $E'[\ell]\simeq E[\ell]$ is irreducible, and every invariant built from the
$G_\Q$-module $T_\ell E$ (the adjoint $\ad T_\ell E$, the length $h_\ell(E)$, the $\ell$-adic Tamagawa factors and Selmer groups of $\ad T_\ell E$) is the same for $E$ and $E'$.
\item $t_\ell(E,q)=t_\ell(E',q)$ for every prime $q$.
\end{enumerate}
\end{lemma}

\begin{proof}
(i) By \Cref{prop:ellcurves:cyclicfactor}, $\varphi=\psi\circ[a]$ with $\psi$ cyclic over $\Q$ of some degree $n$. If $\ell\mid n$, then by \Cref{prop:ellcurves:cyclicfactor}(iii) $E$ has a rational cyclic
$\ell$-isogeny, so $E[\ell]$ has a $G_\Q$-stable line (\Cref{prop:ellcurves:isogline}), contradicting irreducibility. (ii) $\psi$ induces $T_\ell E\simeq T_\ell E'$ by \Cref{prop:ellcurves:rho}(iii), and
everything listed is a function of the $G_\Q$-module $T_\ell E$ (with its restrictions to decomposition groups).
\end{proof}

\begin{proof}[Sketch of proof of (iii)]
$t_\ell(E,q)=v_\ell(-v_q(j_E))$ when $v_q(j_E)<0$. Over a finite extension of
$\Q_q$, $E$ becomes isomorphic to a Tate curve with parameter $q_E$, $v_q(q_E)=-v_q(j_E)$ (\Cref{ch:reduction}). An isogeny of prime degree $r$ between Tate curves has kernel either $\mu_r$ or
generated by an $r$-th root of the parameter, and the parameter of the target is $q_E^r$, resp.\ that $r$-th root; so $v_q(j)$ is multiplied or divided by $r$. Composing, $v_q(j_{E'})=v_q(j_E)\cdot
r_1^{\pm1}\cdots r_k^{\pm1}$ with $r_1\cdots r_k=\deg\psi$ prime to $\ell$, and $v_\ell(-v_q(j))$ is unchanged. (The multiplication map $[a]$ does not change the target.) If $v_q(j_E)\ge0$ then also
$v_q(j_{E'})\ge0$, potentially good reduction being an isogeny invariant (\Cref{ch:reduction}). The facts about Tate curves used here are proved in \Cref{ch:reduction}.
\end{proof}

\begin{corollary}[trivial isogeny classes]\label{cor:ellcurves:trivialclass}
If the isogeny class of $E/\Q$ consists of $E$ alone, then $E$ has no rational $\ell$-isogeny for any prime $\ell$, so $E[\ell]$ is irreducible for every $\ell$ (including $\ell=2$), and $E$ is the
$X_0(N)$-optimal curve of its class. Example: \curve{980i1}.
\end{corollary}

\begin{proof}
A rational $\ell$-isogeny $E\to E'$ would put $E'$ in the class, and $E'\not\simeq E$ over $\Q$, since an isogeny of degree $\ell$ from $E$ to a curve isomorphic to $E$ would be an endomorphism of
degree $\ell$ defined over $\Q$, impossible as $\End_\Q(E)=\Z$ (\Cref{prop:ellcurves:normalised}(ii)) and $\deg[m]=m^2$. Irreducibility follows from \Cref{prop:ellcurves:isogline}. The optimal curve is
a curve of the class (\Cref{ch:moddeg}).
\end{proof}

\begin{paperlink}\label{pl:ellcurves:optimal}
Section 1 of the paper: ``When $E[\ell]$ is irreducible, every curve of the class is linked to $E$ by a cyclic isogeny of degree prime to $\ell$'' is
\Cref{lem:ellcurves:primetol}(i); the consequence for the modular degree is \Cref{ch:moddeg}. By (ii) and (iii) of the lemma, the hypothesis of Theorem thm:main on $\bar\rho_{E,\ell}$, the right-hand side of (eq:identity), and
the terms $e_\ell(E,q)$ (read off $V_\ell E$ at $q$), $t_\ell(E,q)$ and $h_\ell(E)$ of (eq:rule) in Conjecture conj:euler do not depend on which curve of the class is taken. The
remaining term $w_3(E)$ of (eq:rule) is defined by the Kodaira type at $3$, and Kodaira types are not isogeny invariants in general: \curve{121c1} and \curve{121c2},
joined by an $11$-isogeny, have types $\kod{IV}$ and $\kod{IV}^*$ at $11$, and \curve{11a1} and \curve{11a3}, joined by a $5$-isogeny, have $\kod{I}_5$ and $\kod{I}_1$ at
$11$. In the situation of (eq:rule) at $\ell=3$, where $E[3]$ is irreducible and the class is linked by isogenies of degree prime to $3$, the type at $3$ was the same for all
curves of the class in each of the $1{,}295$ classes of Cremona's tables below conductor $10{,}000$ with at least two curves, $3\mid N$, no $3$-isogeny and potentially good
reduction at $3$ ($27$ of which have complex multiplication, which Conjecture conj:euler excludes); we do not prove this in general.
With a $3$-isogeny the type at $3$ does vary: among the $1{,}293$ such classes below $10{,}000$ (with $3\mid N$, at least two curves and potentially good
reduction at $3$), $w_3$ differs within $828$, for example \curve{54a} with types $\kod{IV}^*$, $\kod{II}^*$, $\kod{II}$ and $w_3=2,2,1$; Conjecture conj:euler at $\ell=3$
excludes these classes, as $E[3]$ is then reducible. The conjecture is in any case stated for the optimal curve. Section sec:oot keeps curves ``with trivial isogeny class (so that the curve is optimal and $E[\ell]$
is irreducible for every $\ell$)'': this is \Cref{cor:ellcurves:trivialclass}.
\end{paperlink}

\section{Quadratic twists}\label{sec:ellcurves:twists}

Quadratic twists run through the paper: Lemma lem:twist describes the twist $E^{(q)}$ of $E$ by $\Q(\sqrt{q^*})$, the table of Section 1 and Lemma lem:euler use the twist $E'$ of good reduction at a
type $\kod{I}_0^*$ prime, Theorems thm:twist and thm:quant twist by $\chi_S$, and the Example of Section 3 compares \curve{121a1} with the class \curve{121c}, and \curve{539c1} with \curve{77c1}. This
section defines twists for a general Weierstrass model, computes their invariants, proves $\rho_{E^{(d)},\ell}\simeq\rho_{E,\ell}\otimes\chi_d$ and $a_p(E^{(d)})=\chi_d(p)a_p(E)$, and settles which of
the two possible twists the paper means. The effect of twisting on Kodaira types and conductors (Lemma lem:twist (i), (ii)) is proved in \Cref{ch:reduction}.

\begin{definition}[quadratic twists]\label{def:ellcurves:twist}
Let $\operatorname{char}K\ne2$ and $d\in K^\times$.
\begin{enumerate}[label=(\alph*)]
\item For a model $y^2=x^3+a_2x^2+a_4x+a_6$ of $E$, the \emph{quadratic twist} of $E$ by $d$ (or by $K(\sqrt d)$, or by $\chi_d$) is
\[
E^{(d)}\colon y^2=x^3+da_2x^2+d^2a_4x+d^3a_6 ,
\]
which is isomorphic to $dy^2=x^3+a_2x^2+a_4x+a_6$ by $(x,y)\mapsto(dx,d^2y)$.
\item For a general model, complete the square first (\Cref{lem:ellcurves:identities}(ii)): from $\eta^2=4x^3+b_2x^2+2b_4x+b_6$ one gets $E^{(d)}\colon\eta^2=4x^3+db_2x^2+2d^2b_4x+d^3b_6$, and with
$X=4x$, $Y=4\eta$ the Weierstrass model
\[
E^{(d)}\colon Y^2=X^3+db_2X^2+8d^2b_4X+16d^3b_6,
\]
which has integral coefficients when the $a_i$ and $d$ are integers. For $a_1=a_3=0$ this agrees with (a) (up to the change of coordinates $u=2$).
\item For an odd prime $q$, the paper's $E^{(q)}$ is $E^{(q^*)}$, the twist by $\Q(\sqrt{q^*})$; and for a finite set $S$ of odd primes, $E_S=E^{(d_S)}$ with $d_S=\prod_{q\in S}q^*$
(\Cref{def:ellcurves:compositetwist}).
\end{enumerate}
\end{definition}

\begin{proposition}[first properties of twists]\label{prop:ellcurves:twistbasic}
Let $E\colon y^2=x^3+a_2x^2+a_4x+a_6$ and $d\in K^\times$, $\operatorname{char}K\ne2$.
\begin{enumerate}[label=(\roman*)]
\item Up to isomorphism over $K$, $E^{(d)}$ depends only on the class of $d$ in $K^\times/(K^\times)^2$; $E^{(1)}\simeq E$ and $(E^{(d)})^{(d')}\simeq E^{(dd')}$.
\item Fix $\sqrt d\in\bar K$. The map $\psi\colon E^{(d)}\to E$, $(x,y)\mapsto(x/d,\,y/(d\sqrt d))$, is an isomorphism defined over $K(\sqrt d)$, and $\psi^\sigma=[\chi_d(\sigma)]\circ\psi$ for
$\sigma\in G_K$, where $\psi^\sigma(P)=\sigma(\psi(\sigma^{-1}P))$. Equivalently $\psi(\sigma P)=\chi_d(\sigma)\,\sigma(\psi(P))$ for $P\in E^{(d)}(\bar K)$.
\item Through $\psi$, the $G_K$-set $E^{(d)}(\bar K)$ is $E(\bar K)$ with the \emph{twisted action} $\sigma*P=\chi_d(\sigma)\sigma(P)$. In particular $E^{(d)}(K)\simeq\{P\in E(K(\sqrt d)):\sigma P=-P\}$ for
the generator $\sigma$ of $\Gal(K(\sqrt d)/K)$ when $d$ is not a square.
\item The twists of $E$ in general (the curves over $K$ isomorphic to $E$ over $\bar K$) are classified by $H^1(G_K,\Aut_{\bar K}E)$; when $j\ne0,1728$, $\Aut E=\{\pm1\}$ and every twist is quadratic.
\end{enumerate}
\end{proposition}

\begin{proof}
(i) For $d'=du^2$, the change of coordinates $(u,0,0,0)$ of \Cref{prop:ellcurves:coordchange} transforms $E^{(du^2)}$ into $E^{(d)}$ (it divides the coefficients of $x^2$, $x$, $1$ by $u^2$, $u^4$,
$u^6$). The other two statements are immediate from the formula. (ii) Substituting $x=dx'$, $y=d\sqrt d\,y'$ into the equation of $E^{(d)}$ and dividing by $d^3$ gives the equation of $E$
in $(x',y')$; the inverse is $(x',y')\mapsto(dx',d\sqrt dy')$. For $\sigma\in G_K$, $\sigma(d\sqrt d)=\chi_d(\sigma)d\sqrt d$, so $\psi^\sigma(x,y)=(x/d,\chi_d(\sigma)y/(d\sqrt d))$, which is
$\psi(x,y)$ or $-\psi(x,y)=(x/d,-y/(d\sqrt d))$ according as $\chi_d(\sigma)=\pm1$. (iii) is (ii) restated; a point $P\in E(\bar K)$ corresponds to a $K$-point of $E^{(d)}$ iff $\sigma*P=P$ for all
$\sigma$, i.e.\ $\sigma P=\chi_d(\sigma)P$. (iv) We do not prove this; see \cite[Sections X.2 and X.5]{SilvermanAEC}; for the automorphism groups see \Cref{prop:ellcurves:aut}.
\end{proof}

\begin{proposition}[invariants of twists]\label{prop:ellcurves:twistinv}
\leavevmode
\begin{enumerate}[label=(\roman*)]
\item The model of \Cref{def:ellcurves:twist}(a) has $c_4(E^{(d)})=d^2c_4$, $c_6(E^{(d)})=d^3c_6$, $\Delta(E^{(d)})=d^6\Delta$ and $j(E^{(d)})=j(E)$. The integral model of
\Cref{def:ellcurves:twist}(b) has invariants $2^4d^2c_4$, $2^6d^3c_6$, $2^{12}d^6\Delta$.
\item Let $\operatorname{char}K\ne2,3$. If $j(E)=j(E')\ne0,1728$, then $E'\simeq E^{(d)}$ over $K$ with $d=\dfrac{c_6(E')\,c_4(E)}{c_6(E)\,c_4(E')}$, for any models of $E$ and $E'$.
\item Let $q$ be an odd prime. If $q\equiv1\pmod4$ then $q=q^*$. If $q\equiv3\pmod4$, the twists of $E$ by $q$ and by $q^*$ differ by the twist by $-1$, which multiplies $c_6$ by $-1$, and they
are in general different curves, the twist by $q$ being ramified at $2$ as well. The twist by $q^*$ has invariants $(q^*)^2c_4$, $(q^*)^3c_6$, $(q^*)^6\Delta$; their $q$-adic valuations are those of $q^2c_4$, $q^3c_6$, $q^6\Delta$.
\end{enumerate}
\end{proposition}

\begin{proof}
(i) For $y^2=x^3+Ax^2+Bx+C$, $b_2=4A$, $b_4=2B$, $b_6=4C$, $b_8=4AC-B^2$. Replacing $(A,B,C)$ by $(dA,d^2B,d^3C)$ multiplies $b_i$ by $d^{i/2}$ for $i=2,4,6,8$; since $c_4$, $c_6$, $\Delta$
are homogeneous of weights $4$, $6$, $12$ in this grading, they are multiplied by $d^2$, $d^3$, $d^6$, and $j$ is unchanged. For the integral model, the same computation with
$(A,B,C)=(db_2,8d^2b_4,16d^3b_6)$ gives $2^4d^2c_4$, $2^6d^3c_6$, $2^{12}d^6\Delta$ (checked by machine); this is the change of coordinates with $u=1/2$ applied to $d^2c_4$, $d^3c_6$, $d^6\Delta$.
(ii) Take short models (\Cref{lem:ellcurves:identities}(iii)), which have invariants $6^4c_4$, $6^6c_6$; the quotient defining $d$ is unchanged by such rescalings, and by changes of
coordinates, which multiply it by $u^2$. As $j\ne0,1728$, $c_4,c_6,c_4',c_6'\ne0$, and $j=j'$ is equivalent to $c_6^2/c_4^3=c_6'^2/c_4'^3$. Then $d^2c_4=c_6'^2c_4^3/(c_6^2c_4'^2)=c_4'$ and
$d^3c_6=d^2c_4\cdot dc_6/c_4=c_4'\cdot c_6'/c_4'=c_6'$. By (i), $E^{(d)}$ has the same $c_4$, $c_6$ as $E'$, hence the same short model. (iii) $q=-q^*$ when $q\equiv3\pmod4$, and $E^{(-q^*)}\simeq(E^{(q^*)})^{(-1)}$
by \Cref{prop:ellcurves:twistbasic}(i). By (i) the twist by $-1$ multiplies $c_6$ by $-1$; the field $\Q(\sqrt{-1})$ is ramified at $2$. The last statement is (i) with $v_q(q^*)=v_q(q)=1$.
\end{proof}

\begin{theorem}[the Galois representation of a twist]\label{thm:ellcurves:twistgal}
Let $\operatorname{char}K\ne2$, $d\in K^\times$ and $\ell\ne\operatorname{char}K$. Then $\psi$ of \Cref{prop:ellcurves:twistbasic}(ii) induces a $G_K$-isomorphism $T_\ell E^{(d)}\simeq T_\ell E\otimes\chi_d$; that is,
$\rho_{E^{(d)},\ell}\simeq\rho_{E,\ell}\otimes\chi_d$. Consequently $\bar\rho_{E^{(d)},\ell}\simeq\bar\rho_{E,\ell}\otimes\chi_d$, $E^{(d)}[2]\simeq E[2]$, $\det\rho_{E^{(d)},\ell}=\det\rho_{E,\ell}$,
$\ad T_\ell E^{(d)}\simeq\ad T_\ell E$, and $E^{(d)}[\ell]$ is irreducible iff $E[\ell]$ is.
\end{theorem}

\begin{proof}
$\psi$ is an isomorphism of groups (\Cref{thm:ellcurves:isog-hom}(i)), so it maps $E^{(d)}[\ell^n]$ bijectively to $E[\ell^n]$, compatibly with $[\ell]$, and by \Cref{prop:ellcurves:twistbasic}(ii)
$\psi(\sigma P)=\chi_d(\sigma)\sigma(\psi(P))$. This says that $\psi_\ell\colon T_\ell E^{(d)}\to T_\ell E\otimes\chi_d$ is $G_K$-equivariant. On $E[2]$, $\chi_d(\sigma)=\pm1$ acts trivially. Determinants are
multiplied by $\chi_d^2=1$, and $\ad(T\otimes\chi)=\ad T$ (\Cref{ch:reps}). Twisting by a character preserves the lattice of stable subspaces.
\end{proof}

\begin{corollary}[$a_p$ of a twist]\label{cor:ellcurves:twistap}
Let $E/\Q$, $d\in\Q^\times$ not a square, $D$ the discriminant of $\Q(\sqrt d)$, and $p$ a prime of good reduction for $E$ with $p\nmid D$. Then $E^{(d)}$ has good reduction at $p$ and
\[
a_p(E^{(d)})=\chi_d(p)\,a_p(E),\qquad\chi_d(p)=\Bigl(\frac Dp\Bigr).
\]
\end{corollary}

\begin{proof}
Good reduction of $E^{(d)}$ at $p$: $E$ and $E^{(d)}$ become isomorphic over the extension $\Q_p(\sqrt d)$, unramified at $p$, and good reduction is invariant under unramified extension
(\Cref{ch:reduction}). By \Cref{thm:ellcurves:twistgal} and \Cref{thm:ellcurves:frobQ}, for $\ell\ne p$, $a_p(E^{(d)})=\tr(\rho_{E,\ell}\otimes\chi_d)(\Frob_p)=\chi_d(\Frob_p)a_p(E)$, and $\chi_d(\Frob_p)=\bigl(\frac Dp\bigr)$
(\Cref{conv:ellcurves:conventions}(d)). An elementary proof when $p$ is odd and $p\nmid d\Delta$: with $E\colon y^2=f(x)$ and $E^{(d)}\colon y^2=d^3f(x/d)$, \Cref{prop:ellcurves:twistFq} gives
$a_p(E^{(d)})=-\sum_x\eta(d^3f(x/d))=-\eta(d)\sum_{x'}\eta(f(x'))=\eta(d)a_p(E)$, and $\eta(d)=\bigl(\frac dp\bigr)=\bigl(\frac Dp\bigr)$.
\end{proof}

\noindent Checked for every good $p<400$, $p=2$ included where allowed: \curve{121d1} against \curve{11a3}, \curve{121d2} against \curve{11a1} (with $\chi_{-11}$); \curve{121c2} against \curve{121a1}
($\chi_{-11}$); \curve{77c1} against \curve{539c1} and \curve{140b1} against \curve{980i1} ($\chi_{-7}$).

\begin{proposition}[the twist of good reduction at a type $\kod{I}_0^*$ prime]\label{prop:ellcurves:twistgood}
Let $q$ be an odd prime at which $E/\Q$ has type $\kod{I}_0^*$, and $E'=E^{(q^*)}$, which has good reduction at $q$ (\Cref{ch:reduction}).
\begin{enumerate}[label=(\roman*)]
\item $V_\ell E\simeq V_\ell E'\otimes\chi_{q^*}$ for every $\ell\ne q$, with $V_\ell E'$ unramified at $q$.
\item The twists of $E$ with good reduction at $q$ are the curves $E^{(q^*d)}$ with $q\nmid D$, $D$ the discriminant of $\Q(\sqrt d)$; for them $a_q(E^{(q^*d)})=\chi_d(q)a_q(E')=\pm a_q(E')$. Hence
$a_q(E')^2$, $(q+1)^2-a_q(E')^2$ and the set $\{\alpha/\beta,\beta/\alpha\}$ for the roots of $X^2-a_q(E')X+q$ depend only on $E$ and $q$.
\item Example: \curve{539d1} has type $\kod{I}_0^*$ at $7$, and $\curve{539d1}^{(-7)}\simeq\curve{11a3}$, with $a_7=-2$. So $(7+1)^2-a_7^2=60$ and
$e_3(\curve{539d1},7)=v_3(7-1)+v_3(60)=2=v_3(72)=v_3(\degphi)$ for $E=\curve{539d1}$.
\end{enumerate}
\end{proposition}

\begin{proof}
(i) is \Cref{thm:ellcurves:twistgal} (as $\chi_{q^*}^2=1$) and \Cref{thm:ellcurves:frobQ}. (ii) $E^{(q^*d)}=(E')^{(d)}$ (\Cref{prop:ellcurves:twistbasic}(i)). If $q\nmid D$, it has good reduction at $q$ and
$a_q=\chi_d(q)a_q(E')$ by \Cref{cor:ellcurves:twistap}. If $q\mid D$, it is the twist of the curve $E'$, of good reduction at $q$, by a character ramified at $q$, which has type $\kod{I}_0^*$
(\Cref{ch:reduction}), not good reduction. Replacing $a_q(E')$ by $-a_q(E')$ replaces $(\alpha,\beta)$ by $(-\alpha,-\beta)$ and does not change $\alpha/\beta$. (iii) Computed with PARI: the minimal model of the twist
of \curve{539d1} by $-7$ is $[0,-1,1,0,0]$, which is \curve{11a3}; the modular degree $\degphi=72$ of $E=\curve{539d1}$ is from Cremona's \texttt{degphi} file.
\end{proof}

\begin{definition}[composite twists]\label{def:ellcurves:compositetwist}
For a finite set $S$ of odd primes, $\chi_S=\prod_{q\in S}\chi_{q^*}$ is the quadratic character of $\Q(\sqrt{d_S})$, $d_S=\prod_{q\in S}q^*$, a fundamental discriminant $\equiv1\pmod4$; it has
conductor $\prod_{q\in S}q$ and is unramified outside $S$ (\Cref{ch:rings,ch:ant}). The twist $E_S=E^{(d_S)}$ is obtained by twisting successively by the $q^*$, $q\in S$, in any order
(\Cref{prop:ellcurves:twistbasic}(i)), $\rho_{E_S,\ell}\simeq\rho_{E,\ell}\otimes\chi_S$, and $a_p(E_S)=\chi_S(p)a_p(E)$ for every good $p\notin S$ (\Cref{cor:ellcurves:twistap}).
\end{definition}

\begin{remark}[what a twist changes]\label{rem:ellcurves:twistinvariance}
Under every quadratic twist the invariants $j$, $v_q(j)$, hence $t_\ell(E,q)$, and $\ad\rho_{E,\ell}$ are unchanged (\Cref{prop:ellcurves:twistinv}(i), \Cref{thm:ellcurves:twistgal}). The numbers $a_q$, the
discriminant, the Kodaira type at a prime where the twist ramifies, and the naive Euler factor $L_q^{\nv}$ change. So $e_\ell(E,q)$, which compares the true adjoint factor (an invariant of $\ad V$)
with the naive one, is invariant under twists unramified at $q$ (which change neither factor, by \Cref{ch:reduction}) but not under twists ramified at $q$: \curve{11a1} has type $\kod{I}_5$ at $11$
and $e_3=0$, while its twist $\curve{11a1}^{(-11)}\simeq\curve{121d2}$ has type $\kod{I}_5^*$ and $e_3=v_3(11^2-1)=1$; likewise \curve{11a3} ($\kod{I}_1$) and $\curve{11a3}^{(-11)}\simeq\curve{121d1}$
($\kod{I}_1^*$). The types are from \Cref{ch:reduction}.
\end{remark}

\begin{warning}[$q$ versus $q^*$; curve versus class]\label{warn:ellcurves:twists}
(a) For $q\equiv3\pmod4$ the twist by $\Q(\sqrt q)$ is not the paper's twist. For \curve{121a1}: $\curve{121a1}^{(-11)}\simeq\curve{121c2}$, of conductor $121$, while $\curve{121a1}^{(11)}\simeq\curve{1936k2}$,
of conductor $1936=2^4\cdot11^2$. Only $q$-adic valuations enter Lemma lem:twist (ii), and they agree (\Cref{prop:ellcurves:twistinv}(iii)).
(b) A twist of a curve is a curve, and a twist of a class is a class. With the minimal models of Cremona's tables, as checked with PARI:
\[
\curve{11a1}^{(-11)}\simeq\curve{121d2},\quad\curve{11a3}^{(-11)}\simeq\curve{121d1},\quad\curve{11a2}^{(-11)}\simeq\curve{121d3},\quad\curve{121a1}^{(-11)}\simeq\curve{121c2},
\]
\[
\curve{539c1}^{(-7)}\simeq\curve{77c1},\quad\curve{980i1}^{(-7)}\simeq\curve{140b1},\quad\curve{11a1}^{(-7)}\simeq\curve{539d2},
\]
\[
\curve{11a3}^{(-7)}\simeq\curve{539d1},\quad\curve{20a1}^{(-1)}\simeq\curve{80b1}.
\]
Statements about $a_p$, $L$-functions, $V_\ell$ and (for $\ell$ prime to the isogeny degrees) $T_\ell$ are statements about classes; statements about $j$, $\Delta$, Kodaira types, torsion and $v_q(j)$ are
about curves. ``The twist of \curve{11a1} by $\chi_{-11}$ is in the class \curve{121d}'' is right; ``it is \curve{121d1}'' is not.
\end{warning}

\begin{paperlink}\label{pl:ellcurves:twist}
The proof of Lemma lem:twist (ii) of the paper: ``the twist has a model with invariants $(q^*)^{2}c_4$, $(q^*)^{3}c_6$, $(q^*)^{6}\Delta$ and the same $j$-invariant'' is \Cref{prop:ellcurves:twistinv}(i), (iii), and its
model $y^2=x^3+q^*a_2x^2+(q^*)^2a_4x+(q^*)^3a_6$ is \Cref{def:ellcurves:twist}(a). Lemma lem:euler: in the row $\kod{I}_n^*$, ``$V=V'\otimes\chi_{q^*}$ with $V'$ Steinberg at $q$'' and ``$\ad V=\ad V'$'' are \Cref{thm:ellcurves:twistgal} (the
Steinberg property of $V'$ is \Cref{ch:reduction}); in the row $\kod{I}_0^*$, ``$V=V'\otimes\chi_{q^*}$ with $V'$ unramified'' is \Cref{prop:ellcurves:twistgood}(i). Theorems thm:twist and thm:quant: $E_S$, $\chi_S$ and $a_p(f)=\chi_S(p)a_p(g)$ are \Cref{def:ellcurves:compositetwist}. The table of Section 1: ``$E'$ the quadratic twist of
$E$ by $\Q(\sqrt{q^*})$'' is \Cref{prop:ellcurves:twistgood}, and (ii) of that proposition is why only $a_q(E')^2$ matters. Definition def:terms: ``The term $t_\ell$ is invariant under quadratic twist, and
$e_\ell(E,q)$ under quadratic twists unramified at $q$'' is \Cref{rem:ellcurves:twistinvariance}. The Example of Section 3 (ex:qual) twists \curve{121a1} and \curve{539c1}, \curve{980i1} by
$\chi_{-11}$ and $\chi_{-7}$; its identifications are at the level of classes, and \Cref{warn:ellcurves:twists}(b) gives the curves.
\end{paperlink}
\section{The \texorpdfstring{$L$}{L}-function}\label{sec:ellcurves:Lfunction}

The paper's newform $f=f_E$ is the modular form whose $L$-series is $L(E,s)$, and Lemma lem:old uses the multiplicativity of its coefficients and the recursion at good primes. This section defines
$L(E,s)$ by its Euler product, including the factors at the bad primes, and proves what follows from the results above. Analytic continuation and the functional equation are
\Cref{ch:lfunctions}, modularity is \Cref{ch:moddeg}, and the adjoint $L$-function is \Cref{ch:lfunctions}.

\begin{definition}[the $L$-function of $E$]\label{def:ellcurves:Lfunction}
Let $E/\Q$ have conductor $N$ (\Cref{ch:reduction}), given by its global minimal model. For every prime $p$ let $\tilde E$ be the reduction modulo $p$ and $a_p=p+1-\#\tilde E(\F_p)$, the singular point
included when $p\mid N$; so for $p\mid N$, $a_p=1$, $-1$ or $0$ according as the reduction is split multiplicative, nonsplit multiplicative or additive (\Cref{ch:reduction}). Define
\[
L(E,s)=\prod_{p\nmid N}\bigl(1-a_pp^{-s}+p^{1-2s}\bigr)^{-1}\prod_{p\mid N}\bigl(1-a_pp^{-s}\bigr)^{-1}=\sum_{n\ge1}a_nn^{-s}.
\]
\end{definition}

\begin{proposition}[properties of $L(E,s)$]\label{prop:ellcurves:Lfunction}
\leavevmode
\begin{enumerate}[label=(\roman*)]
\item $\abs{a_n}\le d(n)\sqrt n$, where $d(n)$ is the number of divisors of $n$; the series and the product converge absolutely and agree for $\operatorname{Re}s>3/2$.
\item $a_{mn}=a_ma_n$ for coprime $m,n$; $a_{p^{k+1}}=a_pa_{p^k}-pa_{p^{k-1}}$ for $p\nmid N$, $k\ge1$; and $a_{p^k}=a_p^k$ for $p\mid N$.
\item For $p\ne\ell$ the local factor is $L_p(V_\ell E^\vee,s)=\det\bigl(1-\Frob_p^{-1}p^{-s}\bigm|(V_\ell E^\vee)^{I_p}\bigr)^{-1}$ (\Cref{conv:ellcurves:frobenius}): at good $p$ by
\Cref{thm:ellcurves:frobQ}; at bad $p$, where $\dim(V_\ell E^\vee)^{I_p}$ is $1$ (multiplicative, with $\Frob_p^{-1}$ acting by $a_p=\pm1$) or $0$ (additive), by \Cref{ch:galrep}. At a
good prime this is also $\det(1-\Frob_pp^{-s}\mid V_\ell E)^{-1}$, but not at a multiplicative prime: there the arithmetic Frobenius acts on the line $(V_\ell E)^{I_p}$ by $pa_p$, not by $a_p$. Isogenous curves have the same $L$-function
(\Cref{thm:ellcurves:isogclass}(ii) and the invariance of the conductor), and $L(E^{(d)},s)$ has the same local factors as $\sum\chi_d(n)a_nn^{-s}$ at the primes $p\nmid ND$
(\Cref{cor:ellcurves:twistap}).
\item (Modularity.) $L(E,s)=L(f_E,s)$ for a newform $f_E$ of weight $2$ on $\Gamma_0(N)$; $L(E,s)$ extends to an entire function and satisfies a functional equation under $s\mapsto2-s$.
\end{enumerate}
\end{proposition}

\begin{proof}
(i) At a good prime, $(1-a_pX+pX^2)^{-1}=\bigl((1-\alpha_pX)(1-\beta_pX)\bigr)^{-1}$, so $a_{p^k}=\sum_{i=0}^k\alpha_p^i\beta_p^{k-i}$ and $\abs{a_{p^k}}\le(k+1)p^{k/2}$ by \Cref{thm:ellcurves:charpoly}(iii); at a
bad prime $\abs{a_{p^k}}=\abs{a_p}^k\le1$. With (ii) this gives $\abs{a_n}\le d(n)\sqrt n$. Since $d(n)=O(n^{c})$ for every $c>0$, $\sum\abs{a_n}n^{-\sigma}$ converges for $\sigma>3/2$; the product of
the absolutely convergent local series over $p\le x$ contains every $a_nn^{-s}$ with all prime factors of $n$ at most $x$, and the difference from the full series tends to $0$.
(ii) Expanding the product, $a_n$ is the product of the coefficients $a_{p^k}$ of the local series for $p^k\,\|\,n$; the recursion is the coefficient of $X^{k+1}$ in
$(1-a_pX+pX^2)\sum_ka_{p^k}X^k=1$. (iii) At a good prime $\det(1-\Frob_p^{-1}X\mid V_\ell E^\vee)=\det(1-\Frob_pX\mid V_\ell E)=1-a_pX+pX^2$ by \Cref{conv:ellcurves:frobenius} and \Cref{thm:ellcurves:frobQ}; the bad
primes are \Cref{ch:galrep}. At a multiplicative prime $V_\ell E$ restricted to the decomposition group is an extension of $\varepsilon$ by $\varepsilon\chi_\ell$, with $\varepsilon$ unramified and
$\varepsilon(\Frob_p)=a_p$ (\Cref{ch:reduction}), so the inertia invariants of $V_\ell E$ are the line $\varepsilon\chi_\ell$, with arithmetic Frobenius eigenvalue $pa_p$, and those of
$V_\ell E^\vee\simeq V_\ell E\otimes\chi_\ell^{-1}$ are the line $\varepsilon$, with geometric Frobenius eigenvalue $\varepsilon(\Frob_p)^{-1}=a_p$. (iv) We do not prove this;
it is the theorem of Wiles, Taylor-Wiles and Breuil, Conrad, Diamond and Taylor \cite{Wiles95,TaylorWiles,BCDT}, with the analytic statements from the theory of modular forms; see
\Cref{ch:moddeg,ch:lfunctions}.
\end{proof}

\noindent For \curve{11a1}, $a_1,\dots,a_{13}=1,-2,-1,2,1,2,-2,0,-2,-2,1,-2,4$; for instance $a_4=a_2^2-2=2$, $a_8=a_2a_4-2a_2=0$, $a_9=a_3^2-3=-2$, and $a_{11}=1$ ($11$ is split multiplicative).

\section{Labels and the running examples}\label{sec:ellcurves:labels}

The paper names curves by Cremona's labels and relies on Cremona's convention that the first curve of a class is the optimal one. This section explains the labels, gives their LMFDB
equivalents, and collects in one table the data of the curves that the paper and this book use as running examples. The modular degrees, Kodaira types and images are proved or computed
in \Cref{ch:moddeg,ch:reduction,ch:galrep}; here they are quoted from Cremona's files \cite{ecdata} and checked with PARI.

\begin{convention}[Cremona and LMFDB labels]\label{conv:ellcurves:labels}
\leavevmode
\begin{enumerate}[label=(\alph*)]
\item A Cremona label $N\,c\,k$ consists of the conductor $N$, an isogeny-class code $c$ (a, b, \ldots, z, ba, bb, \ldots) and the number $k$ of the curve in its class: \curve{539c1} is the first curve
of the third class of conductor $539$. In Cremona's tables curve number $1$ of each class is the $X_0(N)$-optimal curve (\Cref{ch:moddeg}), with a single exception, the class \curve{990h}, whose
optimal curve is \curve{990h3} $=[1,-1,1,-1568,-4669]$: the \texttt{degphi} file records the modular degree $1728$ of the class \curve{990h} with this equation. Optimality is not known for every class
of the tables: Section sec:data of the paper says ``Optimality is determined in the tables for every class of conductor at most $400000$ and for $359{,}009$ further classes up to
$500000$'', while its tables go up to $500000$. So below conductor $10{,}000$, the range of this chapter, optimality is determined for every class. Cremona's book \cite{Cremona97} and \cite{ARS12} print the class code in capitals:
\curve{99A1} is \curve{99a1}.
\item LMFDB labels \curve{N.ck} \cite{LMFDB} order the classes and the curves within a class differently, so the same curve can carry different letters and numbers, and the LMFDB's curve number $1$ is in
general not the optimal curve. The correspondence, from Cremona's \texttt{alllabels} file:
\begin{center}\small
\begin{tabular}{llllllll}
\toprule
Cremona & LMFDB & Cremona & LMFDB & Cremona & LMFDB & Cremona & LMFDB\\
\midrule
\curve{11a1} & \curve{11.a2} & \curve{20a1} & \curve{20.a4} & \curve{121a1} & \curve{121.a2} & \curve{539c1} & \curve{539.d2}\\
\curve{11a2} & \curve{11.a1} & \curve{27a1} & \curve{27.a3} & \curve{121c1} & \curve{121.c2} & \curve{539d1} & \curve{539.a3}\\
\curve{11a3} & \curve{11.a3} & \curve{77c1} & \curve{77.c2} & \curve{121c2} & \curve{121.c1} & \curve{726a1} & \curve{726.b2}\\
\curve{14a1} & \curve{14.a6} & \curve{99a1} & \curve{99.a2} & \curve{121d1} & \curve{121.d3} & \curve{980i1} & \curve{980.b1}\\
\curve{140b1} & \curve{140.b1} & \curve{990h3} & \curve{990.i3} & \curve{121d2} & \curve{121.d2} & & \\
\bottomrule
\end{tabular}
\end{center}
\end{enumerate}
\end{convention}

\begin{example}[the paper's curves]\label{ex:ellcurves:running}
Data from Cremona's files (\texttt{allcurves}, \texttt{allisog}, \texttt{degphi}, \texttt{galrep}), rechecked with PARI; ``images'' lists the non-surjective $\bar\rho_{E,\ell}$ in the notation of
\cite{Sutherland16}; types from PARI's \texttt{elllocalred} (\Cref{ch:reduction}).
\begin{center}\footnotesize
\begin{tabular}{lllllll}
\toprule
curve & $N$ & rank, torsion & class (degrees) & $\degphi$ & images & types\\
\midrule
\curve{11a1} & $11$ & $0$, $\Z/5$ & $5,5$ & $1$ & \texttt{5Cs.1.1} & $\kod{I}_5$ at $11$\\
\curve{20a1} & $2^2\cdot5$ & $0$, $\Z/6$ & $2,3,6$ & $1$ & \texttt{2B}, \texttt{3B.1.1} & $\kod{IV}^*$, $\kod{I}_2$ at $2,5$\\
\curve{121a1} & $11^2$ & $0$, $0$ & $11$ & $6$ & \texttt{11B.1.6} & $\kod{II}$ at $11$\\
\curve{121c1} & $11^2$ & $0$, $0$ & $11$ & $6$ & \texttt{11B.1.7} & $\kod{IV}$ at $11$\\
\curve{121d1} & $11^2$ & $0$, $0$ & $5,25$ & $24$ & \texttt{5B.4.1} & $\kod{I}_1^*$ at $11$\\
\curve{539c1} & $7^2\cdot11$ & $1$, $\Z/2$ & $2$ & $288=2^53^2$ & \texttt{2B} & $\kod{I}_3^*$, $\kod{I}_2$ at $7,11$\\
\curve{980i1} & $2^2\cdot5\cdot7^2$ & $0$, $0$ & trivial & $2880=2^63^25$ & none & $\kod{IV}^*$, $\kod{I}_1$, $\kod{I}_5^*$ at $2,5,7$\\
\curve{726a1} & $2\cdot3\cdot11^2$ & $1$, $\Z/2$ & $2$ & $144=2^43^2$ & \texttt{2B}, \texttt{3Ns} & $\kod{I}_6$, $\kod{I}_3$, $\kod{III}$ at $2,3,11$\\
\curve{99a1} & $3^2\cdot11$ & $1$, $\Z/2$ & $2$ & $4$ & \texttt{2B} & $\kod{III}$, $\kod{I}_1$ at $3,11$\\
\midrule
\curve{77c1} & $7\cdot11$ & $0$, $\Z/2$ & $2$ & $6$ & \texttt{2B} & $\kod{I}_3$, $\kod{I}_2$ at $7,11$\\
\curve{140b1} & $2^2\cdot5\cdot7$ & $0$, $0$ & trivial & $60$ & none & $\kod{IV}^*$, $\kod{I}_1$, $\kod{I}_5$ at $2,5,7$\\
\curve{539d1} & $7^2\cdot11$ & $1$, $0$ & $5,25$ & $72$ & \texttt{5B.4.1} & $\kod{I}_0^*$, $\kod{I}_1$ at $7,11$\\
\curve{14a1} & $2\cdot7$ & $0$, $\Z/6$ & $2,3,3,6,6$ & $1$ & \texttt{2B}, \texttt{3Cs.1.1} & $\kod{I}_6$, $\kod{I}_3$ at $2,7$\\
\bottomrule
\end{tabular}
\end{center}
The column ``class'' gives the first row of Cremona's isogeny matrix without the diagonal entry, that is, the degrees of the cyclic isogenies from the curve to the
other curves of its class. \curve{11a1} is the middle curve of its class, $5$-isogenous to both others (which are $25$-isogenous to each other). \curve{121d1} and \curve{539d1},
the twists of \curve{11a3}, are at an end: the class matrices of \curve{121d} and \curve{539d} are
$\left(\begin{smallmatrix}1&5&25\\5&1&5\\25&5&1\end{smallmatrix}\right)$, and the middle curves are \curve{121d2} and \curve{539d2}, the twists of \curve{11a1}. Roles in the paper: \curve{121a1}, \curve{539c1} and \curve{980i1} illustrate the three cases of Theorem thm:qual in the Example of Section 3, with
the congruent forms coming from \curve{11a1} (twisted by $\chi_{-11}$, giving the class \curve{121d}, and by $\chi_{-7}$, giving \curve{539d}) and \curve{20a1} (twisted by $\chi_{-7}$); \curve{121c2} is
$\curve{121a1}^{(-11)}$, \curve{77c1} is $\curve{539c1}^{(-7)}$, \curve{140b1} is $\curve{980i1}^{(-7)}$, and \curve{539d1} is $\curve{11a3}^{(-7)}$ (\Cref{warn:ellcurves:twists}); \curve{726a1} is the paper's
example with image $N_s(\F_3)$ (Section sec:anomaly and the appendix); \curve{99A1} is the curve of \cite{ARS12} with modular degree $4$ and congruence number $12$ (Remark after Theorem
thm:qual; \Cref{ch:moddeg} recomputes $r_E=12$); \curve{14a1} has $E[3]\simeq\Z/3\oplus\mu_3$ (Section sec:eis). Among the first nine curves, \curve{539c1}, \curve{726a1} and \curve{99a1} have rank $1$; so has \curve{539d1}.
The curves \curve{27a1}, \curve{32a1}, \curve{36a1}, \curve{121b1} with CM, and \curve{338d1}, \curve{6912d1} with image \texttt{3Ns}, appear in \Cref{sec:ellcurves:divpol,sec:ellcurves:cm}.
\end{example}

The PARI/GP commands behind these computations (\Cref{ch:computation} explains them):
\begin{lstlisting}[language=GP]
E = ellinit([1,1,1,-30,-76]);   \\ the curve 121a1
[E.c4, E.c6, E.disc, E.j]       \\ invariants
ellglobalred(E)                  \\ conductor, change to the minimal model, Tamagawa product
elllocalred(E, 11)               \\ Kodaira code and Tamagawa number at 11
elltors(E)                       \\ torsion subgroup
ellisomat(E)                     \\ the curves of the isogeny class and the matrix
factor(elldivpol(E, 3))          \\ the 3-division polynomial over Q
vector(12, i, ellap(E, prime(i)))
\end{lstlisting}

\section{Exercises}\label{sec:ellcurves:exercises}

\begin{exercise}\label{ex:ellcurves:ex-99a1}
For \curve{99a1} $=[1,-1,1,-2,0]$ compute $b_2,b_4,b_6,b_8$, $c_4$, $c_6$, $\Delta$ and $j$, check $4b_8=b_2b_6-b_4^2$ and $1728\Delta=c_4^3-c_6^2$, and find a short model $Y^2=X^3+AX+B$ with integral
$A,B$. (Answer: $c_4=3^4$, $c_6=135$, $\Delta=3^3\cdot11$, $j=3^9/11$.)
\end{exercise}

\begin{exercise}\label{ex:ellcurves:ex-valuations}
Use \Cref{prop:ellcurves:coordchange} to show that for a prime $q$ the residue of $v_q(\Delta)$ modulo $12$ and the number $v_q(j)$ are the same for all models of $E/\Q$, and deduce that a model with
integral coefficients and $v_q(\Delta)<12$ has the least $v_q(\Delta)$ among integral models. Give an example of an integral model with $v_q(\Delta)\ge12$ that is not of least $v_q(\Delta)$.
\end{exercise}

\begin{exercise}\label{ex:ellcurves:ex-14a1}
On \curve{14a1} $=[1,0,1,4,-6]$, compute the multiples of $(9,23)$. Using $\#\tilde E(\F_3)=\#\tilde E(\F_5)=6$ and \Cref{prop:ellcurves:torsred}, show $E(\Q)_{\mathrm{tors}}\simeq\Z/6$. Which point of
order $2$ is rational, and over which field are the other two points of order $2$ defined?
\end{exercise}

\begin{exercise}\label{ex:ellcurves:ex-borel}
Let $B\subset\GL_2(\F_3)$ be the upper triangular subgroup and $B_1=\{\sm1*0d\}$. Show that $B_1$ fixes one point of $\PP^1(\F_3)$ and is transitive on the other three. Explain with
\Cref{prop:ellcurves:psi3} why $\psi_3$ of \curve{20a1} (image \texttt{3B.1.1}) has the pattern $1+3$, while that of \curve{14a1} (image \texttt{3Cs.1.1}) has the pattern $1+1+2$, and identify
the root of $\psi_3$ of \curve{14a1} that belongs to the line $\mu_3$.
\end{exercise}

\begin{exercise}\label{ex:ellcurves:ex-726}
For \curve{726a1}, check that $\psi_3=(x^2-x-74)(3x^2+8x+20)$, that the splitting field of $\psi_3$ is $\Q(\sqrt{33},\sqrt{-11})$, and that it contains $\Q(\sqrt{-3})$ as \Cref{prop:ellcurves:psi3}(iii)
requires. Check that $3\mid a_p$ for the good primes $p\equiv2\pmod3$ below $100$, and explain why no such statement holds for $p\equiv1\pmod3$.
\end{exercise}

\begin{exercise}\label{ex:ellcurves:ex-121twist}
Twist \curve{121a1} $=[1,1,1,-30,-76]$ by $d=-11$ using the integral model of \Cref{def:ellcurves:twist}(b), and use \Cref{prop:ellcurves:twistinv}(ii) to show that the result is isomorphic to
\curve{121c2} $=[1,1,0,-3632,82757]$ and not to \curve{121c1}. Then twist by $d=11$ and show that the result has bad reduction at $2$, so it is not isomorphic to any curve of conductor $121$.
\end{exercise}

\begin{exercise}\label{ex:ellcurves:ex-121d}
Show that $j(\curve{121d1})=j(\curve{11a3})\ne j(\curve{11a1})$, so that \curve{121d1} is a twist of \curve{11a3} and not of \curve{11a1}. Check $a_p(\curve{121d1})=\chi_{-11}(p)a_p(\curve{11a1})$ for
$p\le50$, $p\ne11$, and explain why the numbers $a_p$ cannot distinguish \curve{11a1} from \curve{11a3}.
\end{exercise}

\begin{exercise}\label{ex:ellcurves:ex-539}
Show that the twist of \curve{539c1} by $-7$ has $j=167^3/(7^3\cdot11^2)$ and is \curve{77c1}. Compare $v_7(\Delta)$ for the two curves ($9$ against $3$) and explain the difference of $6$ with
\Cref{prop:ellcurves:twistinv}(i) and \Cref{prop:ellcurves:coordchange}(iii).
\end{exercise}

\begin{exercise}\label{ex:ellcurves:ex-Fqn}
Let $s_n=\alpha^n+\beta^n$ for the roots of $X^2-aX+\mathfrak q$. Show $s_{n+1}=as_n-\mathfrak qs_{n-1}$, and use \Cref{thm:ellcurves:charpoly} to compute $\#\curve{11a1}(\F_{27})$ from $a_3=-1$. (Answer: $20$.)
Show that $\#E(\F_{\mathfrak q})$ divides $\#E(\F_{\mathfrak q^n})$ for every $n$.
\end{exercise}

\begin{exercise}\label{ex:ellcurves:ex-I0star}
For $E'=\curve{11a3}$ and $q=7$ compute $\#E'(\F_7)$ and the number of points of the quadratic twist of the reduction, and check that their product is $(7+1)^2-a_7^2=60$
(\Cref{prop:ellcurves:twistFq}). Deduce $e_3(\curve{539d1},7)=2$, and show that replacing $E'$ by the twist $\curve{539d1}^{(-7d)}$ for any $d$ with $7\nmid\disc\Q(\sqrt d)$ gives the same value.
\end{exercise}

\begin{exercise}\label{ex:ellcurves:ex-cm}
Check that $a_p(\curve{27a1})=0$ for the primes $p\equiv2\pmod3$, $p\le50$, and that $a_p(\curve{121b1})=0$ for the primes $p\le50$ inert in $\Q(\sqrt{-11})$; explain both with
\Cref{thm:ellcurves:cm-ap}. Compute $j\bigl(\frac{1+\sqrt{-163}}2\bigr)$ numerically, show that it is $-640320^3$, and explain why $e^{\pi\sqrt{163}}$ is within $10^{-12}$ of an integer.
\end{exercise}

\begin{exercise}\label{ex:ellcurves:ex-fulltorsion}
Use \Cref{cor:ellcurves:det}(ii) to show that if $E/\Q$ has $E[m]\subseteq E(\Q)$ then $m\le2$. Find a curve over $\Q$ with $E[2]\subseteq E(\Q)$, and show that if $E[3]\subseteq E(K)$ for a number field $K$ then
$\sqrt{-3}\in K$.
\end{exercise}

\begin{exercise}\label{ex:ellcurves:ex-11a}
The class \curve{11a} has isogeny matrix $\left(\begin{smallmatrix}1&5&5\\5&1&25\\5&25&1\end{smallmatrix}\right)$ and mod $5$ images \texttt{5Cs.1.1}, \texttt{5B.1.2}, \texttt{5B.1.1} for \curve{11a1},
\curve{11a2}, \curve{11a3}. Show that $T_\ell$ of the three curves are isomorphic for $\ell\ne5$, but that $T_5(\curve{11a1})\not\simeq T_5(\curve{11a2})$. Explain, using
$\curve{11a1}/\langle(5,5)\rangle\simeq\curve{11a2}$ and $\curve{11a1}/\mu_5\simeq\curve{11a3}$, why \curve{11a3} has a rational point of order $5$ and \curve{11a2} has none.
\end{exercise}

\begin{exercise}\label{ex:ellcurves:ex-no25}
Adapt the proof of \Cref{prop:ellcurves:no27} to show, without Mazur's torsion theorem but with \Cref{thm:ellcurves:mazurkenku}(ii), that a curve $E/\Q$ with $E[5]\simeq\Z/5\oplus\mu_5$ has no rational
point of order $25$. Where does the argument use that $\ell$ is odd?
\end{exercise}

\begin{exercise}\label{ex:ellcurves:ex-cyclicdual}
Let $\varphi\colon E\to E'$ be a cyclic $n$-isogeny over $K$ ($\operatorname{char}K=0$). Show that $\hat\varphi$ is cyclic and that $\ker\hat\varphi=\varphi(E[n])$. Now let $E/\Q$ have no CM and let $E'$ be
isogenous to it. Show that $\psi\mapsto\hat\varphi\circ\psi$ embeds $\Hom(E,E')$ into $\End(E)=\Z$, so that $\Hom(E,E')\simeq\Z$, and deduce that $E$ and $E'$ are joined by a cyclic isogeny of a
unique degree: this is why the entries of Cremona's isogeny matrix are well defined.
\end{exercise}

\begin{exercise}\label{ex:ellcurves:ex-lattices}
For $\Lambda=\Z+\Z i$ and $\Lambda=\Z+\Z\zeta_3$ show $g_3=0$, respectively $g_2=0$, and compute $\End(\C/\Lambda)$. Show that the curve $y^2=x^3-x$ (\curve{32a2}, $\Delta=64>0$) has a rectangular period
lattice, and that the lattice of \curve{11a1} (\Cref{ex:ellcurves:11a1-periods}) is not rectangular. Show that for every lattice the map $z\mapsto\bar z$ carries $\Lambda$ to a lattice with the
conjugate $j$-invariant.
\end{exercise}

\begin{exercise}\label{ex:ellcurves:ex-weilconj}
Let $E/\Q$, $\ell$ odd, and $c$ a complex conjugation. Using only \Cref{thm:ellcurves:weil} and $c(\zeta)=\zeta^{-1}$ on $\mu_\ell$, show that
$\det\bar\rho_{E,\ell}(c)=-1$ (compute $e_\ell(cS,cT)$ for a basis $S,T$ of $E[\ell]$). Deduce from $c^2=1$ that $\bar\rho_{E,\ell}(c)$ is conjugate to $\operatorname{diag}(1,-1)$, and
hence that $E(\R)[\ell]$ has order $\ell$. For $\ell=2$ show that the determinant gives no information, and that $c$ acts trivially on $E[2]$ if and only if $F$ has three real
roots, if and only if $\Delta>0$; check this on \curve{32a2} and \curve{11a1}.
\end{exercise}

\begin{exercise}\label{ex:ellcurves:ex-agm}
For \curve{11a1} compute the real root $e_1$ of $F(x)=4x^3-4x^2-40x-79$, the numbers $\gamma$, $\delta$ of \Cref{prop:ellcurves:periods}(iii), and the two arithmetic-geometric means,
and check that the formula gives the periods of \Cref{ex:ellcurves:11a1-periods}. How many iterations of the AGM give $20$ correct digits? (Answer: $e_1=4.3463081582\ldots$,
$\gamma=12.0389244746\ldots$, $\delta=6.1626753526\ldots$, $\omega_1=1.2692093042\ldots$, $\operatorname{Im}\omega_2=1.4588166169\ldots$; the convergence is quadratic.)
\end{exercise}

\begin{exercise}\label{ex:ellcurves:ex-labels}
Using the table of \Cref{conv:ellcurves:labels}, find the curves whose Cremona and LMFDB numbers differ, and check with the \texttt{alllabels} and \texttt{degphi} files of \cite{ecdata} that
\curve{990h3} is the optimal curve of \curve{990h}. Why can the LMFDB's ordering not be used to identify the optimal curve?
\end{exercise}

\section{Notes and references}\label{sec:ellcurves:notes}

\paragraph{Sources followed.} Sections \ref{sec:ellcurves:weierstrass} to \ref{sec:ellcurves:weil} follow Silverman's first book \cite{SilvermanAEC}, Chapter III (Weierstrass equations in III.1,
the group law in III.2 and III.3, isogenies in III.4 to III.6, the Tate module in III.7, the Weil pairing in III.8, endomorphism rings in III.9, automorphisms in III.10), with the division
polynomials of \cite[Exercise 3.7]{SilvermanAEC} and \cite[Chapter 3]{WashingtonEC}. Section \ref{sec:ellcurves:cm} follows \cite[Chapter II]{SilvermanATAEC} and \cite{Cox}; Deuring's theorem is
\cite{Deuring41}. Section \ref{sec:ellcurves:complex} follows \cite[Chapter VI]{SilvermanAEC} and \cite[Chapter 1]{DiamondShurman}, with the period lattice of a curve over $\Q$ as in
\cite[Chapter III]{Cremona97}; the AGM algorithm is \cite[Section 3.7]{Cremona97} and \cite[Section 7.4]{Cohen93}. Section \ref{sec:ellcurves:finite} follows \cite[Chapter V]{SilvermanAEC}; the
proof of Hasse's theorem by the positivity of the degree is Silverman's. Section \ref{sec:ellcurves:overQ} follows \cite[Chapters VII and VIII]{SilvermanAEC} for torsion and Mordell-Weil;
the torsion theorem is \cite{Mazur77}, the prime-degree isogeny theorem \cite{Mazur78}, and the remaining cyclic isogenies are in Kenku's papers \cite{Kenku79,Kenku80a,Kenku80b,Kenku81}; the
complete list of cyclic degrees is \cite[Theorem 1]{Kenku} and the bound on the size of isogeny classes \cite[Theorem 2]{Kenku} (both checked against the published text); the isogeny theorem is \cite{Faltings83}. V\'elu's formulas are \cite{Velu71}. Section \ref{sec:ellcurves:twists} follows
\cite[Section X.2 and X.5]{SilvermanAEC}; Section \ref{sec:ellcurves:Lfunction} follows \cite[Appendix C, Section 16]{SilvermanAEC} and \cite[Chapter II]{Cremona97}. Other textbook accounts:
\cite{Husemoller}, \cite{Knapp}, \cite{MilneEC}; for the Galois representations, \cite[Chapter I]{SerreAbelian}.

\paragraph{Theorem and section numbers.} The chapter cites sections rather than theorem numbers. The section numbers of \cite{SilvermanAEC} (second edition), \cite{SilvermanATAEC},
\cite{WashingtonEC}, \cite{Cox}, \cite{Cohen93}, \cite{Cremona97} and \cite{DiamondShurman} given above are pointers recalled from those books; they have not been checked against the printed
editions. The table of rational cyclic isogenies in \cite{LozanoRobledo13} was not checked either.

\paragraph{The proposition on points of order $9$.} \Cref{prop:ellcurves:no27} is Proposition prop:no9 of the revised paper; the argument through the cyclic $27$-isogeny of $E/\mu_3$ is
elementary once \Cref{thm:ellcurves:mazurkenku}(iii) is known. The data of Cremona's tables below conductor $10{,}000$ agree with it, as recorded after the proof.

\paragraph{Computations.} All numbers in the chapter were computed with PARI/GP \cite{PARI} through \texttt{cypari2}, on Cremona's tables \cite{ecdata} for conductor below $10{,}000$: the invariants
of \Cref{ex:ellcurves:running-invariants}, the multiples of points, the division polynomials and their factorisations and discriminants, the numbers $a_p$ and the point counts over $\F_{p^2}$ and
$\F_{27}$, the Hasse bound for the good primes below $2000$ of seven curves, the thirteen CM $j$-values and the CM checks, the twist identifications (by computing the minimal model of the twisted
equation and comparing with the tables), the isogeny matrices, the LMFDB labels, the torsion data behind \Cref{prop:ellcurves:no27}, the consistency of the \texttt{galrep} and \texttt{allisog} files
($114{,}126$ pairs of a curve and $\ell\in\{3,5,7\}$, with $59$ discrepancies, all for CM curves, for which the \texttt{galrep} file gives no Borel label), the Kodaira types at $3$
in the classes without a $3$-isogeny (\Cref{pl:ellcurves:optimal}), and the periods of \curve{11a1}, also by the AGM. The symbolic identities of \Cref{lem:ellcurves:identities},
\eqref{eq:ellcurves:invdiff-identity} in the proof of \Cref{prop:ellcurves:invdiff}, \Cref{prop:ellcurves:coordchange}, \Cref{prop:ellcurves:addition}(iii), \Cref{prop:ellcurves:psi3}(i) and \Cref{prop:ellcurves:twistinv}(i) were checked
with the $a_i$ as indeterminates.
